\documentclass[12pt]{article}

% Idioma
\usepackage[T1]{fontenc}
\usepackage[spanish]{babel}

% Matemática
\usepackage{amsmath, amssymb, amsthm}
\usepackage{mathtools} % Para \coloneq

\usepackage[hidelinks]{hyperref}
\usepackage{varioref, cleveref}
\usepackage{algorithm}
\usepackage{algpseudocode}
\usepackage{tikz}
\usetikzlibrary{babel}
\usetikzlibrary{decorations.pathreplacing}
\usetikzlibrary{arrows.meta}

\usepackage{float}      % para poder usar [H] si querés fijar la figura
\usepackage{subcaption} % para subfiguras lado a lado
\usetikzlibrary{calc}
\usepackage{subcaption}
\usepackage{xspace}
% ---------- Entornos de teoremas ----------
\theoremstyle{plain} % estilo "teorema" (itálica)
\newtheorem{theorem}{Teorema}[section]
\newtheorem{lemma}[theorem]{Lemma}
\newtheorem{corollary}[theorem]{Corolario}

\theoremstyle{definition} % estilo "definición" (normal)
\newtheorem{definition}[theorem]{Definición}
\newtheorem{example}[theorem]{Ejemplo}

\theoremstyle{remark} % estilo "observación" (normal + título distinto)
\newtheorem{Claim}[theorem]{Claim}
\newtheorem{remark}[theorem]{Observación}

%\newtheorem{problem}{Problem}

\usepackage[most]{tcolorbox}
\usepackage{enumitem}

\newtcolorbox{problem}[1][]{
  enhanced,
  breakable,
  colback=white,
  colframe=black,
  boxrule=0.6pt,
  arc=1mm,
  left=7pt,
  right=7pt,
  top=7pt,
  bottom=7pt,
  title={Problem: \textsc{#1}},
  fonttitle=\bfseries
}






% ---------- Notación útil de grafos ----------
\newcommand{\edgecondition}{\emph{\hyperref[def:avoid:cond:edge]{edge condition}}\xspace}
\newcommand{\nonedgecondition}{\emph{\hyperref[def:avoid:cond:nonedge]{non-edge condition}}\xspace}
\newcommand{\colorcondition}{\emph{\hyperref[def:avoid:cond:color]{color condition}}\xspace}
\newcommand{\adj}{\sim}
\newcommand{\V}{V}
\newcommand{\E}{E}
\newcommand{\N}{\mathbb{N}}
\newcommand{\degG}{\operatorname{deg}}
\newcommand{\dist}{\operatorname{dist}}
\newcommand{\chiG}{\chi}
\newcommand{\segto}{\operatorname{seg}_{\to}}
\newcommand{\segfrom}{\operatorname{seg}_{\leftarrow}}
\newcommand{\segg}{\operatorname{seg}_{12}}
\newcommand{\segthree}{\operatorname{seg}_{3}}
\newcommand{\segthirteen}{\operatorname{seg}_{13}}
\newcommand{\segonetwothree}{\operatorname{seg}_{123}}
\newcommand{\minClases}{\operatorname{minClases}}
\newcommand{\chisemi}{\chi_{\mathrm{semi}}}

\newcommand{\segone}{\operatorname{bip}_{1}}
\newcommand{\Pone}{\mathcal{P}_{1}}

\newcommand{\segtwo}{\operatorname{seg}_{2}}
\newcommand{\Ptwo}{\mathcal{P}_{2}}

\newcommand{\segonetwo}{\operatorname{seg}_{12}}
\newcommand{\Ponetwo}{\mathcal{P}_{12}}

\newcommand{\Pthree}{\mathcal{P}_{3}}

\newcommand{\segonethree}{\operatorname{seg}_{13}}
\newcommand{\Ponethree}{\mathcal{P}_{13}}

\newcommand{\Ponetwothree}{\mathcal{P}_{123}}

\newcommand{\estone}{\operatorname{est}_{1}}
\newcommand{\EStone}{E_1}

\newcommand{\eric}[1]{\textcolor{green}{Eric: #1}}
\newcommand{\german}[1]{\textcolor{blue}{German: #1}}

\newcommand{\tauvc}{\tau}

\newcommand{\esttwothree}{\operatorname{est}_{23}}
\newcommand{\mathcalEtwothree}{\mathcal{E}_{23}}

\newcommand{\estonethree}{\operatorname{est}_{13}}
\newcommand{\mathcalEonethree}{\mathcal{E}_{13}}

\newcommand{\estoneonetwothree}{\operatorname{est}_{123}}
\newcommand{\mathcalEonetwothree}{\mathcal{E}_{123}}

\newcommand{\forrs}{\operatorname{for}_{rs}}
\newcommand{\forrt}{\operatorname{for}_{rt}}
\newcommand{\forst}{\operatorname{for}_{st}}


\newcommand{\mathcalFrsst}{\mathcal F_{rs\text{-}st}}
\newcommand{\mathcalFrsrt}{\mathcal F_{rs\text{-}rt}}
\newcommand{\mathcalFrtst}{\mathcal F_{rt\text{-}st}}
\newcommand{\mathcalFrsstrt}{\mathcal F_{rs\text{-}st\text{-}rt}}




\title{Tesis}
\author{German y Eric}
\date{\today}

\begin{document}
\maketitle
\tableofcontents

\section{Preliminares}




\german{reduccion de star12 con intuicion}
\german{unificacion induced parameters}
\german{escribi q no lo tenia star 12 23 13}
\german{reduccion bip 12 23 13 con intuicion}
\german{caso k=2 para dado orden}

\begin{definition}[Pattern on three vertices]
A \emph{pattern} on three ordered vertices \(r\prec s\prec t\) is a triple
\[
P=(A,N,C),
\]
where
\[
A,N\subseteq \binom{\{r,s,t\}}{2},
\qquad
A\cap N=\varnothing,
\qquad 
C=\varnothing.
\]
The set \(A\) indicates the pairs of vertices that are required to be edges, while \(N\) indicates the pairs of
vertices that are required to be non-edges. The condition \(C=\varnothing\) means that the pattern imposes no
condition on the parts of a partition. Pairs that do not belong to \(A\cup N\) are left unrestricted.
\end{definition}


\german{falta definir familia de patrones}

\begin{figure}[h] 
\centering
\begin{tikzpicture}[
  font=\small,
  every node/.style={circle, draw, inner sep=1.4pt},
  >=latex
]
\node (p1) at (0,0) {$1$};
\node (p2) at (1.7,0) {$2$};
\node (p3) at (3.4,0) {$3$};

\draw[line width=0.6pt] (p1)--(p2);
\draw[densely dashed, line width=0.6pt] (p2)--(p3);

\end{tikzpicture}
\caption{The pattern \(P=(A,N, C)\), where \(A=\{\{1,2\}\}\) and \(N=\{\{2,3\}\}\). Solid lines represent required edges, and dashed lines represent required non-edges.}
\label{fig:tiny-pattern-uncolored}
\end{figure}
\medskip

\begin{definition}[Colored pattern on three vertices]
A \emph{colored pattern} is a triple \(P=(A,N,C)\) where
\[
A,N \subseteq \binom{\{1,2,3\}}{2},
\qquad
C \in \binom{\{1,2,3\}}{2}.
\]
Here \(1,2,3\) denote positions of an ordered triple, and
\[
\binom{\{1,2,3\}}{2}=\{\{1,2\},\{1,3\},\{2,3\}\}.
\]
The set \(A\) specifies the pairs that must be edges, \(N\) specifies the pairs that must be non-edges,
and \(C\) specifies the pair of vertices that must lie in the same part of the partition.
\end{definition}


\begin{figure}[ht]
\centering
\begin{tikzpicture}[
  font=\small,
  every node/.style={circle, draw, inner sep=1.4pt},
  >=latex
]
\node[draw=red, very thick, fill=red!12] (p1) at (0,0) {$1$};
\node (p2) at (1.7,0) {$2$};
\node[draw=red, very thick, fill=red!12] (p3) at (3.4,0) {$3$};

\draw[line width=0.6pt] (p1)--(p2);
\draw[densely dashed, line width=0.6pt] (p2)--(p3);

\end{tikzpicture}
\caption{The forbidden colored pattern \(P=(A,N,C)\), where \(A=\{\{1,2\}\}\), \(N=\{\{2,3\}\}\), and \(C=\{1,3\}\). Solid lines represent required edges, dashed lines represent required non-edges, and red vertices represent the pair required to lie in the same part.}
\label{fig:tiny-pattern}
\end{figure}


\begin{definition}[Solution and value]
Let \(G\) be a graph. A \emph{solution} for $G$ is a pair \(S=(\prec,\mathcal Q)\), where \(\prec\) is a total order on \(V(G)\)
and \(\mathcal Q=\{V_1,\dots,V_k\}\) is a partition of \(V(G)\) into \(k\) parts.
The \emph{value} of \(S\) is \(|\mathcal Q|=k\).
\end{definition}


\eric{Definir avoid a pattern, y definir avoid colored pattern en base a eso.}
\german{done}
\begin{definition}[Avoiding a pattern]\label{def:avoiding-pattern}
Fix a graph \(G\), a total order \(\prec\) of \(V(G)\), and a pattern
\[
P=(A,N,\varnothing).
\]

An \emph{ordered triple} of \(G\) is a triple
\[
(v_1,v_2,v_3)
\]
of distinct vertices of \(G\) such that
\[
v_1\prec v_2\prec v_3.
\]

We say that the ordered triple \((v_1,v_2,v_3)\) \emph{avoids} \(P\) in \(G\) if at least one of the following
conditions holds:

\begin{description}
\item[\textup{(edge)}]\phantomsection\label{def:avoid:cond:edge}
There exists \(\{a,b\}\in A\) such that
\[
\{v_a,v_b\}\notin E(G).
\]

\item[\textup{(non-edge)}]\phantomsection\label{def:avoid:cond:nonedge}
There exists \(\{a,b\}\in N\) such that
\[
\{v_a,v_b\}\in E(G).
\]
\end{description}

The order \(\prec\) \emph{avoids} \(P\) in \(G\) if every ordered triple of \(G\) avoids \(P\).

For a family of patterns \(\mathcal P\), the order \(\prec\) avoids \(\mathcal P\) in \(G\) if it avoids every
pattern \(P\in\mathcal P\).


When the underlying graph is clear from the context, we omit the phrase ``in \(G\)'' and simply say that an ordered triple, an order, or a solution avoids \(P\) (or \(\mathcal P\)).


\end{definition}
\german{en algun lado escribir que los param aca estudiados la nonedge nunca se cumple, luego u
una sol no avoidea el patron si no cumple ni edge ni color (lo uso mucho eso)}

\begin{definition}[Avoiding a colored pattern]\label{def:avoiding-colored-pattern}
Fix a graph \(G\), a solution
\[
S=(\prec,\mathcal Q)
\]
for \(G\), and a colored pattern
\[
P=(A,N,C).
\]

Let
\[
v_1\prec v_2\prec v_3
\]
be an ordered triple of \(G\). We say that \((v_1,v_2,v_3)\) \emph{avoids} \(P\) with respect to \(S\) if it
avoids the pattern
\[
(A,N,\varnothing)
\]
in \(G\), according to Definition~\ref{def:avoiding-pattern}, or if the following condition holds:

\begin{description}
\item[\textup{(color)}]\phantomsection\label{def:avoid:cond:color}
If \(C=\{a,b\}\), then \(v_a\) and \(v_b\) lie in different parts of \(\mathcal Q\).
\end{description}

Equivalently, an ordered triple avoids \(P\) if at least one of
\edgecondition, \nonedgecondition, or \colorcondition{} holds.

The solution \(S\) \emph{avoids} \(P\) in \(G\) if every ordered triple of \(G\) avoids \(P\) with respect to \(S\).

For a family of colored patterns \(\mathcal P\), the solution \(S\) avoids \(\mathcal P\) in \(G\) if it avoids every
colored pattern \(P\in\mathcal P\).
\end{definition}

\german{mi duda aca es que aun no defini family of patterns}


\begin{definition}[\(P\)-partition number]
Let \(P=(A,N,\varnothing)\) be a pattern on three vertices. We define
\[
\minClases_P(G)
\]
as the minimum \(k\) such that there exists a partition
\[
V(G)=V_1\,\dot\cup\,\cdots\,\dot\cup\,V_k
\]
satisfying
\[
G[V_i]\in \mathcal C_P
\qquad\text{for every } i=1,\dots,k.
\]
\end{definition}

\eric{Este problema ya existe en la literatura, así que vamos a tener que citar resultados. O mencionar que es NP-hard, por lo menos en la intro. OK todavia no me meti en el lio de las citas}
\german{puede ser que k<=2 sea poli para nuesteos parametros?}






\begin{definition}[Parameter induced by forbidden patterns]
For a finite family of colored patterns \(\mathcal P\), define \(\minClases_{\mathcal P}(G)\) as the minimum value
of a solution for \(G\) that avoids \(\mathcal P\) in \(G\).
\end{definition}
\german{1.5 y 1.6 son parecidas peor una de color y otra de uncolored, esta bien?}
\eric{O unirlas o poner la 1.5 antes de la 1.6} \german{Done}




\german{esto es un intento de unificar la notacion de los parametros inducidos por patrones coloreados y no coloreados, pero no se si es buena idea.}
\begin{definition}[Notation for the induced parameters]
For each \(ij\in\{12,13,23\}\), we define the singleton families
\[
\mathcal E_{ij}:=\{\operatorname{star}_{ij}\},
\qquad
\mathcal B_{ij}:=\{\operatorname{bip}_{ij}\},
\qquad
\mathcal F_{ij}:=\{\operatorname{for}_{ij}\}.
\]

For every index
\[
D\in
\bigl\{
12,13,23,
12\text{-}13,
12\text{-}23,
13\text{-}23,
12\text{-}13\text{-}23
\bigr\},
\]
we denote the parameters induced by the corresponding families by
\[
\operatorname{star}_{D}(G)
:=
\minClases_{\mathcal E_D}(G),
\]
\[
\operatorname{bip}_{D}(G)
:=
\minClases_{\mathcal B_D}(G),
\]
and
\[
\operatorname{for}_{D}(G)
:=
\minClases_{\mathcal F_D}(G).
\]

For the three uncolored patterns, we similarly write
\[
\operatorname{star}(G)
:=
\minClases_{\operatorname{star}}(G),
\qquad
\operatorname{bip}(G)
:=
\minClases_{\operatorname{bip}}(G),
\qquad
\operatorname{for}(G)
:=
\minClases_{\operatorname{for}}(G).
\]
\end{definition}

\begin{remark}
We use the same symbol \(\operatorname{star}_{ij}\), \(\operatorname{bip}_{ij}\), or
\(\operatorname{for}_{ij}\) for a colored pattern, and add the argument \(G\) when referring to
the corresponding graph parameter. Thus, for example,
\[
\operatorname{for}_{12}
\]
denotes a colored pattern, whereas
\[
\operatorname{for}_{12}(G)
\]
denotes the minimum number of parts in a solution for \(G\) that avoids that pattern.
Similarly,
\[
\operatorname{for}_{12\text{-}23}(G)
=
\minClases_{\mathcal F_{12\text{-}23}}(G).
\]
\end{remark}



\begin{definition}[Solución óptima]\label{def:sol-optima}
Sea \(G\) un grafo y sea \(\mathcal P\) una familia de patrones coloreados.
Una solución \(S=(\prec,\mathcal Q)\) para $G$ se dice \emph{óptima (para \(\mathcal P\))} si \(S\) evita \(\mathcal P\) in \(G\) y
\[
|\mathcal Q|=\minClases_{\mathcal P}(G).
\]
\end{definition}


\begin{definition}[Class associated with a pattern on three vertices]
Let
\[
P=(A,N)
\]
be a pattern. We denote by \(\mathcal C_P\) the class of graphs characterized by \(P\);
that is,
\[
G\in \mathcal C_P
\]
if and only if there exists a total order \(\prec\) of \(V(G)\) such that every ordered triple of vertices of \(G\)
avoids the pattern \(P\).
\end{definition}

















\begin{theorem}\label{thm:colored-lower-bound-by-class-partition}
Let \(P=(A,N,C)\) be a colored pattern, and consider the pattern $P' = (A,N)$.

Then, for every graph \(G\),
\[
\minClases_P(G)
\;\ge\;
\minClases_{P'}(G)
.
\]
\end{theorem}

\begin{proof}
Let
\[
S=(\prec,\mathcal Q),
\qquad
\mathcal Q=\{V_1,\dots,V_k\},
\]
be a solution that avoids \(P\).

Fix a part \(V_i\), and restrict the order \(\prec\) to \(V_i\). We claim that this restricted order avoids
\(P'\) in \(G[V_i]\). Indeed, if an ordered triple inside \(V_i\) realized \(P'\), then the required edge and
non-edge conditions of \(P\) would be satisfied. Moreover, since all three vertices lie in the same part \(V_i\),
the color condition \(C\) would also be satisfied. Hence the same triple would realize \(P\), contradicting that
\(S\) avoids \(P\).

Therefore,
\[
G[V_i]\in \mathcal C_{P'}
\qquad\text{for every } i=1,\dots,k.
\]
Thus the partition \(\mathcal Q\) witnesses that
\[
\minClases_{P'}(G)\le k.
\]
Taking the minimum over all solutions avoiding \(P\), we obtain
\[
\minClases_{P'}(G)\le \minClases_P(G),
\]
as desired.
\end{proof}













\subsection{Base uncolored patterns}

We first introduce the uncolored patterns associated with the graph classes that we want to generalize.


\begin{definition}[Basic uncolored patterns]
We define the following uncolored patterns on the ordered vertices \(1\prec 2\prec 3\).
\begin{itemize}
    \item The \emph{star pattern} is
    \[
    \operatorname{star}
    :=
    \bigl(A_{\operatorname{star}},N_{\operatorname{star}}\bigr),
    \]
    where
    \[
    A_{\operatorname{star}}
    =
    \bigl\{\{1,2\}\bigr\},
    \qquad
    N_{\operatorname{star}}
    =
    \varnothing.
    \]

    \item The \emph{forest pattern} is
    \[
    \operatorname{for}
    :=
    \bigl(A_{\operatorname{for}},N_{\operatorname{for}}\bigr),
    \]
    where
    \[
    A_{\operatorname{for}}
    =
    \bigl\{\{1,3\},\{2,3\}\bigr\},
    \qquad
    N_{\operatorname{for}}
    =
    \varnothing.
    \]

    \item The \emph{bipartite pattern} is
    \[
    \operatorname{bip}
    :=
    \bigl(A_{\operatorname{bip}},N_{\operatorname{bip}}\bigr),
    \]
    where
    \[
    A_{\operatorname{bip}}
    =
    \bigl\{\{1,2\},\{2,3\}\bigr\},
    \qquad
    N_{\operatorname{bip}}
    =
    \varnothing.
    \]
\end{itemize}

\begin{figure}[h]
\centering

\begin{subfigure}[t]{0.3\textwidth}
\centering
\begin{tikzpicture}[
  scale=1,
  every node/.style={circle, draw, inner sep=1.4pt, font=\small}
]
\node (a) at (0,0) {$1$};
\node (b) at (1.5,0) {$2$};
\node (c) at (3,0) {$3$};
\draw (a)--(b);
\end{tikzpicture}

\caption{\(\operatorname{star}\) pattern.}
\label{fig:star-uncolored}
\end{subfigure}
\hfill
\begin{subfigure}[t]{0.3\textwidth}
\centering
\begin{tikzpicture}[
  scale=1,
  every node/.style={circle, draw, inner sep=1.4pt, font=\small}
]
\node (a) at (0,0) {$1$};
\node (b) at (1.5,0) {$2$};
\node (c) at (3,0) {$3$};
\draw (a)--(b);
\draw (b)--(c);
\end{tikzpicture}
\caption{\(\operatorname{bip}\) pattern.}
\label{fig:bip-uncolored}
\end{subfigure}
\hfill
\begin{subfigure}[t]{0.3\textwidth}
\centering
\begin{tikzpicture}[
  scale=1,
  every node/.style={circle, draw, inner sep=1.4pt, font=\small}
]
\node (a) at (0,0) {$1$};
\node (b) at (1.5,0) {$2$};
\node (c) at (3,0) {$3$};
\draw[bend left=35] (a) to (c);
\draw (b)--(c);
\end{tikzpicture}
\caption{\(\operatorname{for}\) pattern.}
\label{fig:forest-uncolored}
\end{subfigure}

\caption{The three basic uncolored patterns.}
\label{fig:basic-uncolored-patterns}
\end{figure}
\end{definition}





\begin{remark}
The associated classes
\[
\mathcal C_{\operatorname{star}},
\qquad
\mathcal C_{\operatorname{bip}},
\qquad
\mathcal C_{\operatorname{for}}
\]
are, respectively, stars possibly with isolated vertices, bipartite graphs, and forests.
Thus, the parameters
\[
\minClases_{\operatorname{star}},
\qquad
\minClases_{\operatorname{bip}},
\qquad
\minClases_{\operatorname{for}}
\]
are the corresponding partition parameters into these graph classes.
\end{remark}



\subsection{Colored versions}



\begin{definition}[Bipartite colored patterns]
For each
\[
\{a,b\}\in\binom{\{1,2,3\}}{2},
\]
we define the colored pattern
\[
\operatorname{bip}_{ab}:=(A_{\mathrm{bip}},N_{\mathrm{bip}},\{a,b\}).
\]
In other words, the three bipartite colored patterns are
\[
\operatorname{bip}_{12},\qquad
\operatorname{bip}_{13},\qquad
\operatorname{bip}_{23}.
\]
They all have the same required edges \(12\) and \(23\), and differ only in the pair of positions required to lie in the same part.

\begin{figure}[h]
\centering

\begin{subfigure}[t]{0.3\textwidth}
\centering
\begin{tikzpicture}[
  scale=1,
  every node/.style={circle, draw, inner sep=1.4pt, font=\small},
  same/.style={draw=red, very thick, fill=red!12}
]
\node[same] (v1) at (0,0) {$1$};
\node[same] (v2) at (1.5,0) {$2$};
\node       (v3) at (3,0) {$3$};
\draw (v1)--(v2);
\draw (v2)--(v3);
\end{tikzpicture}
\caption{\(\operatorname{bip}_{12}\)}
\label{fig:bip12}
\end{subfigure}
\hfill
\begin{subfigure}[t]{0.3\textwidth}
\centering
\begin{tikzpicture}[
  scale=1,
  every node/.style={circle, draw, inner sep=1.4pt, font=\small},
  same/.style={draw=red, very thick, fill=red!12}
]
\node[same] (v1) at (0,0) {$1$};
\node       (v2) at (1.5,0) {$2$};
\node[same] (v3) at (3,0) {$3$};
\draw (v1)--(v2);
\draw (v2)--(v3);
\end{tikzpicture}
\caption{\(\operatorname{bip}_{13}\)}
\label{fig:bip13}
\end{subfigure}
\hfill
\begin{subfigure}[t]{0.3\textwidth}
\centering
\begin{tikzpicture}[
  scale=1,
  every node/.style={circle, draw, inner sep=1.4pt, font=\small},
  same/.style={draw=red, very thick, fill=red!12}
]
\node       (v1) at (0,0) {$1$};
\node[same] (v2) at (1.5,0) {$2$};
\node[same] (v3) at (3,0) {$3$};
\draw (v1)--(v2);
\draw (v2)--(v3);
\end{tikzpicture}
\caption{\(\operatorname{bip}_{23}\)}
\label{fig:bip23}
\end{subfigure}

\caption{The three basic colored bipartite patterns.}
\label{fig:bip-basic-patterns}
\end{figure}

\end{definition}




\begin{definition}[Star colored patterns]
For each
\[
\{a,b\}\in\binom{\{1,2,3\}}{2},
\]
we define the colored pattern
\[
\operatorname{est}_{ab}:=(A_{\mathrm{est}},N_{\mathrm{est}},\{a,b\}).
\]
In other words, the three star colored patterns are
\[
\operatorname{est}_{12},\qquad
\operatorname{est}_{13},\qquad
\operatorname{est}_{23}.
\]
They all have the same required edge \(12\), and differ only in the pair of positions required to lie in the same part.

\begin{figure}[h]
\centering

\begin{subfigure}[t]{0.3\textwidth}
\centering
\begin{tikzpicture}[
  scale=1,
  every node/.style={circle, draw, inner sep=1.4pt, font=\small},
  same/.style={draw=red, very thick, fill=red!12}
]
\node[same] (v1) at (0,0) {$1$};
\node[same] (v2) at (1.5,0) {$2$};
\node       (v3) at (3,0) {$3$};
\draw (v1)--(v2);
\end{tikzpicture}
\caption{\(\operatorname{est}_{12}\)}
\label{fig:est12}
\end{subfigure}
\hfill
\begin{subfigure}[t]{0.3\textwidth}
\centering
\begin{tikzpicture}[
  scale=1,
  every node/.style={circle, draw, inner sep=1.4pt, font=\small},
  same/.style={draw=red, very thick, fill=red!12}
]
\node[same] (v1) at (0,0) {$1$};
\node       (v2) at (1.5,0) {$2$};
\node[same] (v3) at (3,0) {$3$};
\draw (v1)--(v2);
\end{tikzpicture}
\caption{\(\operatorname{est}_{13}\)}
\label{fig:est13}
\end{subfigure}
\hfill
\begin{subfigure}[t]{0.3\textwidth}
\centering
\begin{tikzpicture}[
  scale=1,
  every node/.style={circle, draw, inner sep=1.4pt, font=\small},
  same/.style={draw=red, very thick, fill=red!12}
]
\node       (v1) at (0,0) {$1$};
\node[same] (v2) at (1.5,0) {$2$};
\node[same] (v3) at (3,0) {$3$};
\draw (v1)--(v2);
\end{tikzpicture}
\caption{\(\operatorname{est}_{23}\)}
\label{fig:est23}
\end{subfigure}

\caption{The three basic colored star patterns.}
\label{fig:est-basic-patterns}
\end{figure}
\end{definition}



\begin{definition}[Forest colored patterns]
For each
\[
\{a,b\}\in\binom{\{1,2,3\}}{2},
\]
we define the colored pattern
\[
\operatorname{for}_{ab}:=(A_{\mathrm{for}},N_{\mathrm{for}},\{a,b\}).
\]
In other words, the three forest colored patterns are
\[
\operatorname{for}_{12},\qquad
\operatorname{for}_{13},\qquad
\operatorname{for}_{23}.
\]
They all have the same required edges \(13\) and \(23\), and differ only in the pair of positions required to lie in the same part.

\begin{figure}[h]
\centering

\begin{subfigure}[t]{0.3\textwidth}
\centering
\begin{tikzpicture}[
  scale=1,
  every node/.style={circle, draw, inner sep=1.4pt, font=\small},
  same/.style={draw=red, very thick, fill=red!12}
]
\node[same] (r) at (0,0) {$r$};
\node[same] (s) at (1.5,0) {$s$};
\node       (t) at (3,0) {$t$};
\draw[bend left=35] (r) to (t);
\draw (s)--(t);
\end{tikzpicture}
\caption{\(\forrs\)}
\label{fig:forrs}
\end{subfigure}
\hfill
\begin{subfigure}[t]{0.3\textwidth}
\centering
\begin{tikzpicture}[
  scale=1,
  every node/.style={circle, draw, inner sep=1.4pt, font=\small},
  same/.style={draw=red, very thick, fill=red!12}
]
\node[same] (r) at (0,0) {$r$};
\node       (s) at (1.5,0) {$s$};
\node[same] (t) at (3,0) {$t$};
\draw[bend left=35] (r) to (t);
\draw (s)--(t);
\end{tikzpicture}
\caption{\(\forrt\)}
\label{fig:forrt}
\end{subfigure}
\hfill
\begin{subfigure}[t]{0.3\textwidth}
\centering
\begin{tikzpicture}[
  scale=1,
  every node/.style={circle, draw, inner sep=1.4pt, font=\small},
  same/.style={draw=red, very thick, fill=red!12}
]
\node       (r) at (0,0) {$r$};
\node[same] (s) at (1.5,0) {$s$};
\node[same] (t) at (3,0) {$t$};
\draw[bend left=35] (r) to (t);
\draw (s)--(t);
\end{tikzpicture}
\caption{\(\forst\)}
\label{fig:forst}
\end{subfigure}

\caption{The three basic colored forest patterns.}
\label{fig:for-basic-patterns}
\end{figure}

\end{definition}




\begin{definition}[Families of star colored patterns]
We define the following families of star colored patterns:
\[
\mathcal E_{12\text{-}23}:=
\{\operatorname{est}_{12},\operatorname{est}_{23}\},
\qquad
\mathcal E_{12\text{-}13}:=
\{\operatorname{est}_{12},\operatorname{est}_{13}\},
\]
\[
\mathcal E_{13\text{-}23}:=
\{\operatorname{est}_{13},\operatorname{est}_{23}\},
\qquad
\mathcal E_{12\text{-}23\text{-}13}:=
\{\operatorname{est}_{12},\operatorname{est}_{23},\operatorname{est}_{13}\}.
\]
A solution avoids one of these families if it avoids every colored pattern in the family.
\end{definition}


\begin{definition}[Families of bipartite colored patterns]
We define the following families of bipartite colored patterns:
\[
\mathcal B_{12\text{-}23}:=
\{\operatorname{bip}_{12},\operatorname{bip}_{23}\},
\qquad
\mathcal B_{12\text{-}13}:=
\{\operatorname{bip}_{12},\operatorname{bip}_{13}\},
\]
\[
\mathcal B_{13\text{-}23}:=
\{\operatorname{bip}_{13},\operatorname{bip}_{23}\},
\qquad
\mathcal B_{12\text{-}23\text{-}13}:=
\{\operatorname{bip}_{12},\operatorname{bip}_{23},\operatorname{bip}_{13}\}.
\]
A solution avoids one of these families if it avoids every colored pattern in the family.
\end{definition}


\begin{definition}[Families of forest colored patterns]
We define the following families of forest colored patterns:
\[
\mathcal F_{12\text{-}23}:=
\{\operatorname{for}_{12},\operatorname{for}_{23}\},
\qquad
\mathcal F_{12\text{-}13}:=
\{\operatorname{for}_{12},\operatorname{for}_{13}\},
\]
\[
\mathcal F_{13\text{-}23}:=
\{\operatorname{for}_{13},\operatorname{for}_{23}\},
\qquad
\mathcal F_{12\text{-}23\text{-}13}:=
\{\operatorname{for}_{12},\operatorname{for}_{23},\operatorname{for}_{13}\}.
\]
A solution avoids one of these families if it avoids every colored pattern in the family.
\end{definition}








\begin{lemma}\label{thm:minclases-one-class}
Let \(P\) be a pattern. Then, for every graph \(G\),
\[
\minClases_P(G)=1
\quad\Longleftrightarrow\quad
G\in \mathcal C_P.
\]
\end{lemma}

\begin{proof}
First suppose that
\[
\minClases_P(G)=1.
\]
Then there exists a partition of \(V(G)\) into one part,
\[
V(G)=V_1,
\]
such that
\[
G[V_1]\in \mathcal C_P.
\]
Since \(V_1=V(G)\), we have \(G[V_1]=G\). Therefore,
\[
G\in \mathcal C_P.
\]

Conversely, suppose that
\[
G\in \mathcal C_P.
\]
Consider the trivial partition
\[
V_1:=V(G).
\]
Then
\[
G[V_1]=G\in \mathcal C_P.
\]
Hence \(G\) admits a partition into one part satisfying the definition of \(\minClases_P(G)\), so
\[
\minClases_P(G)\le 1.
\]
Since every partition has at least one part, we conclude that
\[
\minClases_P(G)=1.
\]

Thus,
\[
\minClases_P(G)=1
\quad\Longleftrightarrow\quad
G\in \mathcal C_P.
\]
\end{proof}


In words, when \(C=\varnothing\), the parameter \(\minClases_P\) is the partition version of the original
class \(\mathcal C_P\). Having value \(1\) means exactly that no partition is needed: the whole graph already
belongs to the class characterized by \(P\).
\medskip






\german{fijate si borras esto}
\begin{definition}[Patrón prohibido para \(\segone\)]
Se fija el patrón coloreado
\[
P_1=(A,N,C),
\]
donde
\[
A=\bigl\{\{1,2\},\{2,3\}\bigr\},\qquad
N=\varnothing,\qquad
C=\bigl\{\{1,2\}\bigr\}.
\]
\end{definition}

\begin{figure}[ht]
\centering
\begin{tikzpicture}[scale=1, every node/.style={circle, draw, inner sep=1.2pt, font=\small}]
\node[fill=gray!20] (a) at (0,0) {1};
\node[fill=gray!20] (b) at (1.2,0) {2};
\node             (c) at (2.4,0) {3};
\draw (a)--(b);
\draw (b)--(c);
\end{tikzpicture}
\end{figure}

\begin{definition}[\(\segone\)]
Dado un grafo \(G\), se define
\[
\segone(G)\;:=\;\minClases_{\{P_1\}}(G),
\]
es decir, el mínimo número de clases de una solución \((\prec,\mathcal{Q})\) de \(G\) que evita el patrón \(P_1\).
\end{definition}




\eric{Este es un resultado ya conocido en el paper de patrones. Si querés lo podemos dejar en la tesis, por ejemplo en un apéndice, para decir que lo probaste de nuevo. En el paper definitivamente no iría.}
\begin{theorem}\label{thm:seg1-bip}
Sea \(G=(V,E)\) un grafo. Entonces \(G\) es bipartito si y solo si \(\segone(G)=1\).
\end{theorem}

DEMO NUEVA
cambiar los nombres de la variables por la A
\begin{proof}
\noindent(\(\Rightarrow\)) Supongamos que \(G\) es bipartito. Entonces existe una partición
\(V=V_1\,\dot\cup\,V_2\) tal que \(V_1\) y \(V_2\) son independientes.
Consideremos la partición en una sola clase \(\mathcal{Q}:=\{V\}\) y se fija un orden \(\prec\) \emph{por bloques}, en el sentido de que
\[
\text{para todo } v_1\in V_1 \text{ y todo } v_2\in V_2,\qquad v_1\prec v_2.
\]

Dentro de cada bloque \(V_1\) y \(V_2\) se toma un orden arbitrario.

Sea \(P_1=(A,N,C)\) el patrón definido previamente para \(\segone\), con
\(A=\{\{1,2\},\{2,3\}\}\), \(N=\varnothing\) y \(C=\{\{1,2\}\}\).
Verifiquemos que \((\prec,\mathcal{Q})\) evita \(P_1\),
sea \(v_1\prec v_2\prec v_3\) una terna de vértices.


Supongamos que simultáneamente \(v_1\sim v_2\) y \(v_2\sim v_3\), en particular \(v_1\in V_1\) y \(v_2\in V_2\),(pues \(V_1\) es independiente y el orden ubica a todos los vértices de \(V_2\) a la derecha de los de \(V_1\)),
además como \(v_2\prec v_3\) se deduce \(v_3\in V_2\), pero esto es absurdo porque \(V_2\) es independiente, por lo tanto,
se tiene \(v_2\not\sim v_3\).

Luego \(\segone(G)\le 1\) y, como siempre vale \(\segone(G)\ge 1\), se concluye \(\segone(G)=1\).




\medskip
\noindent(\(\Leftarrow\)) Supongamos ahora que \(\segone(G)=1\),
por definición, existe una solución \((\prec,\mathcal{Q})\) con \(|\mathcal{Q}|=1\) que evita \(P_1\),
en particular, \(\mathcal{Q}=\{V\}\), de modo que para toda tripla de vértices
\(r\prec s\prec t\) no pueden ocurrir simultáneamente
\(r\sim s\) y \(s\sim t\) o sea,
\[
(r\sim s)\ \Rightarrow\ (s\not\sim t)
\qquad\text{para todo } r\prec s\prec t
\]
En particular, si un vértice tiene algún vecino a la izquierda, entonces no puede tener vecinos a la derecha 
y viceversa.


Se definen los conjuntos
\[
L:=\{\,s\in V:\ \text{no existe } t\in V \text{ con } s\prec t \text{ y } s\sim t\,\},
\]
\[
D:=\{\,s\in V:\ \text{no existe } r\in V \text{ con } r\prec s \text{ y } r\sim s\,\}.
\]
Es decir, \(L\) es el conjunto de vértices sin vecinos a la derecha, y \(D\) el de vértices sin vecinos a la izquierda.

Veamos que \(L\) es independiente, si \(r\prec s\) con \(r,s\in L\) y \(r\sim s\), entonces \(r\) tendría
un vecino a la derecha (\(s\)), contradicción con la definición de \(L\),
luego no hay aristas internas en \(L\), analogo para \(D\).



Por la observación anterior, todo vértice no aislado pertenece a \(L\) o a \(D\) y
los vértices aislados no afectan la independencia, y pueden asignarse al final a cualquiera de los dos conjuntos.
Luego obtuvimos una partición de \(V\) en dos conjuntos independientes, y por lo tanto \(G\) es bipartito.
\end{proof}

\german{supongo esto hay que borrarlo, lo de arriba}



Un posible algoritmo de backtracking para decidir si \(\segone(G)\le k\) consiste en construir
incrementalmente un orden \(v_1\prec\cdots\prec v_n\) y una partición
\[
V \;=\; V_1\,\dot\cup\,\cdots\,\dot\cup\,V_k.
\]
En cada paso se elige un vértice no ubicado, se lo agrega al final del orden parcial y se le asigna
una de las \(k\) clases.

Tras cada asignación se verifica la condición correspondiente a \(\segone\) en lo ya construido:
si el nuevo vértice queda en la posición \(t\), se chequea que no existan \(r<s<t\) tales que
\(v_r\) y \(v_s\) estén en la misma clase, \(v_r\sim v_s\) y \(v_s\sim v_t\).
Si se detecta una violación, se descarta la rama y se retrocede; si no, se continúa.


El procedimiento utilizado fue probar grafos conocidos y asi poder conjeturar y demostrar lo siguiente:

\begin{theorem}\label{thm:segone-chi}

\[
\segone(G)=\chi(G)-1.
\]
\end{theorem}

DEMO NUEVA




\begin{Claim}\label{lem:seg1-le}
\(\segone(G)\le \chi(G)-1\).
\end{Claim}

\begin{proof}
Sea \(m=\chi(G)\) y sea \(c:V\to\{1,\dots,m\}\) un \(m\)-coloreo propio de \(G\).
Para cada \(i\in\{1,\dots,m\}\) se define la \emph{clase de color}
\[
C_i:=\{\,v\in V:\ c(v)=i\,\}.
\]
A partir de estas clases se construye una partición \(\mathcal{Q}\) de \(V\) en \(m-1\) clases dada por
\[
V_1:=C_1\cup C_m,
\qquad
V_i:=C_i \ \ \text{para todo } i=2,\dots,m-1,
\]
y se escribe \(\mathcal{Q}:=\{V_1,\dots,V_{m-1}\}\). Como \(C_1,\dots,C_m\) forman una partición de \(V\),
también lo hacen \(V_1,\dots,V_{m-1}\).

Además, se fija un orden total \(\prec\) \emph{por bloques} imponiendo que
\[
C_1 \prec C_2 \prec \cdots \prec C_{m-1} \prec C_m,
\]
esto es, para todo \(i<j\), para todo \(u\in C_i\) y todo \(v\in C_j\), se cumple \(u\prec v\),
y dentro de cada bloque \(C_i\) se toma un orden arbitrario.

Probemos que \((\prec,\mathcal{Q})\) evita el patrón \(P_1=(A,N,C)\) que define \(\segone\),
donde
\[
A=\bigl\{\{1,2\},\{2,3\}\bigr\},\qquad N=\varnothing,\qquad C=\bigl\{\{1,2\}\bigr\}.
\]
Sea \(r\prec s\prec t\) una terna de vértices.
Si \(r\not\sim s\), entonces falla la arista \(\{1,2\}\in A\) y si \(r\) y \(s\) pertenecen a clases distintas de \(\mathcal{Q}\), entonces falla la condición de color
asociada a \(\{1,2\}\in C\) supongamos entonces que \(r\sim s\) y que \(r\) y \(s\) pertenecen a una misma clase de \(\mathcal{Q}\).


Como todas las clases \(V_i\) para \(i=2,\dots,m-1\) son independientes, necesariamente
\(r,s\in V_1\) y por construcción,
\[
V_1=C_1\,\dot\cup\, C_m,
\]
y el orden dentro de \(V_1\) ubica a todos los vértices de \(C_1\) antes que todos los de \(C_m\)
por lo tanto, de \(r\prec s\) y \(r,s\in V_1\) se deduce
\[
r\in C_1
\quad\text{y}\quad
s\in C_m.
\]

Pero \(C_m\) es independiente (por ser clase de un coloreo propio), lo cual contradice
\(s\sim t\), por lo tanto, \(s\not\sim t\), es decir, falla la arista \(\{2,3\}\in A\), y la terna
\(r\prec s\prec t\) evita \(P_1\).

Como la terna fue arbitraria, \((\prec,\mathcal{Q})\) evita \(P_1\) y tiene \(|\mathcal{Q}|=m-1\).
En consecuencia,
\[
\segone(G)\le m-1=\chi(G)-1,
\]
como se quería demostrar.
\end{proof}



\begin{Claim}\label{lem:le-seg1}
Se cumple \(\chi(G)-1 \le \segone(G)\).
\end{Claim}

\begin{proof}
Sea \(G\) tal que admite un orden \(v_1 \prec \cdots \prec v_n\) y una partición
\[
V \;=\; V_1 \,\dot\cup\, \cdots \,\dot\cup\, V_k,
\qquad\text{con } k=\segone(G),
\]
de modo que la solución \(S=(\prec,\{V_1,\dots,V_k\})\) evita el patrón \(P_1=(A,N,C)\) (patrón \emph{bip1}),
donde
\[
A=\bigl\{\{1,2\},\{2,3\}\bigr\},\qquad N=\varnothing,\qquad C=\{1,2\}.
\]

Se define un coloreo \(c:V\to\{0,1,\dots,k\}\) del siguiente modo. Para cada \(s\in\{1,\dots,n\}\),
se fija
\[
c(v_s)=
\begin{cases}
0, & \text{si existe } r<s \text{ tal que } v_r \in V_i,\ v_s\in V_i
      \text{ y } v_r \sim v_s,\\[2mm]
i, & \text{en caso contrario, donde \(i\in\{1,\dots,k\}\) es tal que } v_s\in V_i.
\end{cases}
\]

Se afirma que si \(c(v_s)=0\), entonces \(v_s\) no tiene vecinos a la derecha en el orden.
En efecto, por definición de \(c\) existe \(r<s\) tal que \(v_r\) y \(v_s\) pertenecen a una misma clase \(V_i\)
y además \(v_r\sim v_s\).
Si existiera \(t>s\) con \(v_s\sim v_t\), entonces la terna ordenada \((v_r,v_s,v_t)\) realizaría el patrón \(P_1\):
las posiciones \(1\) y \(2\) estarían en la misma clase (porque \(v_r,v_s\in V_i\)), y se tendrían las aristas
correspondientes a \(\{1,2\}\) y \(\{2,3\}\). Esto contradice que \(S\) evita \(P_1\).
Por lo tanto, para todo \(t>s\) se cumple \(v_s\not\sim v_t\).

Se verifica ahora que \(c\) es un coloreo propio.

\smallskip
\noindent\textbf{(a) Dos vértices de color \(0\) no pueden ser adyacentes.}
Si \(v_i\sim v_j\) con \(i<j\) y \(c(v_i)=c(v_j)=0\), entonces \(v_i\) tendría un vecino
a la derecha (a saber, \(v_j\)), contradicción.

\smallskip
\noindent\textbf{(b) Dos vértices con el mismo color \(i\in\{1,\dots,k\}\) no pueden ser adyacentes.}
Si \(c(v_a)=c(v_b)=i\) con \(a<b\), entonces \(v_a,v_b\in V_i\). Si además ocurriera \(v_a\sim v_b\),
por la definición de \(c\) el vértice \(v_b\) debería satisfacer \(c(v_b)=0\), contradicción.

Por (a) y (b), no existe arista cuyos extremos tengan el mismo color; por lo tanto,
\(c\) es un coloreo propio que utiliza a lo sumo \(k+1\) colores. En consecuencia,
\[
\chi(G)\le k+1=\segone(G)+1,
\]
como se quería demostrar.
\end{proof}


SEG2 


\begin{definition}[Patrón prohibido para \(\segtwo\)]
Se fija el patrón coloreado
\[
P_2=(A,N,C),
\]
donde
\[
A=\bigl\{\{1,2\},\{2,3\}\bigr\},\qquad
N=\varnothing,\qquad
C=\{2,3\}.
\]
\end{definition}

\begin{figure}[ht]
\centering
\begin{tikzpicture}[scale=1, every node/.style={circle, draw, inner sep=1.2pt, font=\small}]
\node                    (a) at (0,0)   {1};
\node[draw=red, very thick, fill=red!20] (b) at (1.2,0) {2};
\node[draw=red, very thick, fill=red!20] (c) at (2.4,0) {3};
\draw (a)--(b);
\draw (b)--(c);
\end{tikzpicture}
\end{figure}

\begin{definition}[\(\segtwo\)]
Dado un grafo \(G\), se define
\[
\segtwo(G)\;:=\;\minClases_{\{P_2\}}(G),
\]
es decir, el mínimo número de clases de una solución \((\prec,\mathcal{Q})\) de \(G\) que evita el patrón \(P_2\).
\end{definition}



\begin{definition}[\(\segonetwo\)]
Se fija la familia de patrones
\[
\Ponetwo := \{P_1,P_2\}.
\]
Dado un grafo \(G\), se define
\[
\segonetwo(G)\;:=\;\minClases_{\Ponetwo}(G),
\]
es decir, el mínimo número de clases de una solución \((\prec,\mathcal{Q})\) de \(G\) que evita simultáneamente
los patrones \(P_1\) y \(P_2\).
\end{definition}





\begin{lemma}[Monotonicidad por inclusión de patrones]\label{lem:monotone-patterns}
Sea \(G\) un grafo y sean \(\mathcal P,\mathcal P'\) dos familias de patrones coloreados.
Si \(\mathcal P\subseteq \mathcal P'\), entonces
\[
\minClases_{\mathcal P}(G)\;\le\;\minClases_{\mathcal P'}(G).
x\]
\end{lemma}

\begin{proof}
Toda solución que evita simultáneamente todos los patrones de \(\mathcal P'\) evita, en particular,
todos los patrones de \(\mathcal P\) (pues \(\mathcal P\subseteq \mathcal P'\)).
Por lo tanto, el conjunto de soluciones factibles para \(\minClases_{\mathcal P'}(G)\) es un subconjunto
del conjunto de soluciones factibles para \(\minClases_{\mathcal P}(G)\), y el mínimo número de clases
no puede disminuir al imponer más restricciones. En consecuencia,
\(\minClases_{\mathcal P}(G)\le \minClases_{\mathcal P'}(G)\).
\end{proof}





\begin{theorem}\label{thm:seg12-chi}
Para todo grafo \(G\),
\[
\segg(G)=\chi(G)-1.
\]
\end{theorem}

\begin{proof}
La demostración se divide en dos desigualdades.

\medskip
\begin{lemma}\label{lem:seg12-ge}
Para todo grafo \(G\), se cumple \(\segg(G)\ge \chi(G)-1\).
\end{lemma}

\begin{proof}
Por definición, \(\segone(G)=\minClases_{\{P_1\}}(G)\) y \(\segg(G)=\minClases_{\{P_1,P_2\}}(G)\).
Como \(\{P_1\}\subseteq \{P_1,P_2\}\), por el Lema \ref{lem:monotone-patterns} se obtiene
\[
\segone(G)\le \segg(G).
\]
Usando el Teorema \ref{thm:segone-chi} (i.e.\ \(\segone(G)=\chi(G)-1\)), se concluye que
\[
\chi(G)-1=\segone(G)\le \segg(G).
\]
\end{proof}





\medskip
\begin{lemma}\label{lem:seg12-le}
Para todo grafo \(G\), se cumple \(\segg(G)\le \chi(G)-1\).
\end{lemma}




\begin{proof}
Para probar \(\segg(G)\le \chi(G)-1\) se reutiliza la misma construcción del
Lema \ref{lem:segto-le}. Sea \(m=\chi(G)\) y sea \(c:V\to\{1,\dots,m\}\) una \(m\)-coloración propia.
Para cada \(i\in\{1,\dots,m\}\) se considera la clase de color
\[
C_i:=\{v\in V: c(v)=i\}.
\]
Se define la partición \(V=V_1 \,\dot\cup\, \cdots \,\dot\cup\, V_{m-1}\) por
\[
V_1:=C_1\cup C_m,
\qquad
V_i:=C_i \ \ \text{para } i=2,\dots,m-1,
\]
y se fija un orden \(\prec\) por bloques imponiendo
\[
C_1 \prec C_2 \prec \cdots \prec C_{m-1} \prec C_m,
\]
tomando además un orden arbitrario dentro de cada bloque \(C_i\).

Por el Lema \ref{lem:segto-le}, esta partición y este orden satisfacen \textup{(Seg\(_{\to}\))}.
Resta verificar que también satisfacen \textup{(Seg\(_{\leftarrow}\))}.

Sea entonces \(r<s<t\) una terna de índices del orden fijado. Supóngase que \(v_s\) y \(v_t\)
pertenecen a una misma clase \(V_i\) y que \(v_s\sim v_t\). Se probará que necesariamente
\(v_r\not\sim v_s\).

Primero se observa que no puede ocurrir \(i\in\{2,\dots,m-1\}\): en tal caso \(V_i=C_i\), y como
\(C_i\) es una clase de una coloración propia, es un conjunto independiente. Esto contradice
la hipótesis \(v_s\sim v_t\). Por lo tanto, se concluye que
\[
i=1 \quad\text{y en particular}\quad v_s,v_t\in V_1=C_1\cup C_m.
\]

Se determina ahora en qué subclase cae cada vértice. Dado que \(s<t\) y el orden por bloques ubica
todos los vértices de \(C_1\) antes que todos los vértices de \(C_m\), el único modo de tener
\(v_s\sim v_t\) con \(v_s,v_t\in C_1\cup C_m\) es que
\[
v_s\in C_1
\quad\text{y}\quad
v_t\in C_m,
\]
ya que \(C_1\) y \(C_m\) son independientes y por lo tanto no pueden contener una arista interna.

Finalmente, como \(C_1\) es el primer bloque del orden, no existe ningún vértice estrictamente a la izquierda
de un vértice de \(C_1\). En particular, si \(v_s\in C_1\) entonces no existe \(r<s\) tal que
\(v_r\sim v_s\). Por consiguiente, para la terna considerada se cumple \(v_r\not\sim v_s\), y queda verificada
\textup{(Seg\(_{\leftarrow}\))}.

Como la partición y el orden satisfacen simultáneamente \textup{(Seg\(_{\to}\))} y \textup{(Seg\(_{\leftarrow}\))},
se obtiene que \(\segg(G)\le m-1\). En conclusión,
\[
\segg(G)\le m-1=\chi(G)-1,
\]
como se quería demostrar.
\end{proof}


\medskip
Combinando los Lemas \ref{lem:seg12-ge-segto} y \ref{lem:seg12-le} se obtiene la igualdad.
\end{proof}


SEG3



\begin{definition} [\(\segthree\)]
Definimos el patrón coloreado
\[
P_3=(A,N,C),
\]
donde
\[
A=\bigl\{\{1,2\},\{2,3\}\bigr\},\qquad
N=\varnothing,\qquad
C={\{1,3}\}.
\]
\end{definition}

\begin{figure}[ht]
\centering
\begin{tikzpicture}[scale=1, every node/.style={circle, draw, inner sep=1.2pt, font=\small}]
\node[draw=red, very thick] (a) at (0,0)   {1};
\node                      (b) at (1.2,0) {2};
\node[draw=red, very thick] (c) at (2.4,0) {3};
\draw (a)--(b);
\draw (b)--(c);
\end{tikzpicture}
\end{figure}

\begin{definition}[\(\segthree\)]
Dado un grafo \(G\), definimos
\[
\segthree(G)\;:=\;\minClases_{\{\Pthree\}}(G),
\]
es decir, el mínimo número de clases de una solución \((\prec,\mathcal{Q})\) de \(G\) que evita el patrón \(\Pthree\).
\end{definition}



\begin{theorem}\label{thm:seg3-chi}
Para todo grafo \(G\),
\[
\segthree(G)=\left\lceil \frac{\chi(G)}{2}\right\rceil.
\]
\end{theorem}

\begin{proof}
Se prueban ambas desigualdades.

\medskip
\begin{Claim}\label{cl:seg3-le}
Para todo grafo \(G\),
\[
\segthree(G)\le \left\lceil \frac{\chi(G)}{2}\right\rceil.
\]
\end{Claim}

\begin{proof}
Sea \(k=\chi(G)\). Fijamos un \(k\)-coloreo propio, es decir, una partición
\[
V \;=\; C_1 \,\dot\cup\, C_2 \,\dot\cup\, \cdots \,\dot\cup\, C_k,
\]
tal que cada \(C_i\) es independiente.

Construyamos una solución \((\prec,\mathcal{Q})\) con \(|\mathcal{Q}|=\lceil k/2\rceil\) que evita el patrón
prohibido \(P_3\) (el patrón que define \(\segthree\)).

\medskip
\noindent\textbf{Caso 1: \(k\) par.} Sea \(k=2q\).
Definimos la partición \(\mathcal{Q}=\{V_1,\dots,V_q\}\) uniendo colores de a dos:
\[
V_1:=C_1\cup C_2,\quad
V_2:=C_3\cup C_4,\quad
\dots,\quad
V_q:=C_{2q-1}\cup C_{2q}.
\]
El orden \(\prec\) se fija por bloques imponiendo
\[
V_1 \prec V_2 \prec \cdots \prec V_q,
\]
y dentro de cada bloque \(V_j=C_{2j-1}\cup C_{2j}\) se ubican primero todos los vértices de \(C_{2j-1}\)
y luego todos los de \(C_{2j}\) (y dentro de cada \(C_i\), orden arbitrario).

Verifiquemos que \((\prec,\mathcal{Q})\) evita \(P_3\)
supongamos por absurdo que existe una terna \(x\prec y\prec z\) que no evita \(P_3\),
entonces, por la forma de \(P_3\), necesariamente \(x\) y \(z\) pertenecen a una misma clase \(V_j\) y además
ocurren \(x\sim y\) e \(y\sim z\),
como el orden es por bloques, se tiene \(y\in V_j\), dentro de \(V_j\) hay dos posibilidades:
\begin{itemize}
\item Si \(y\in C_{2j-1}\), entonces (por cómo se ordenó el bloque) también \(x\in C_{2j-1}\), y como
\(C_{2j-1}\) es independiente se obtiene \(x\not\sim y\), contradicción.
\item Si \(y\in C_{2j}\), entonces también \(z\in C_{2j}\), y como \(C_{2j}\) es independiente se obtiene
\(y\not\sim z\), contradicción.
\end{itemize}
Luego no existe tal terna, entonces todas evitan \(P_3\).

\medskip
\noindent\textbf{Caso 2: \(k\) impar.} Sea \(k=2q+1\), definimos
\[
V_1:=C_1\cup C_2,\quad
V_2:=C_3\cup C_4,\quad
\dots,\quad
V_q:=C_{2q-1}\cup C_{2q},\qquad
V_{q+1}:=C_{2q+1}.
\]
El orden se define por bloques \(V_1\prec\cdots\prec V_{q+1}\) como antes, dentro de cada bloque doble
\(V_j=C_{2j-1}\cup C_{2j}\) se ubica primero \(C_{2j-1}\) y luego \(C_{2j}\), y dentro de \(V_{q+1}\)
tomamos un orden arbitrario.

La verificación de que \(P_3\) se evita es igual al caso par cuando \(x\) y \(z\) caen en algún \(V_j\) con \(j\le q\).
Si \(x,z\in V_{q+1}=C_{2q+1}\), entonces también \(y\in C_{2q+1}\) y, al ser \(C_{2q+1}\) independiente,
no pueden ocurrir simultáneamente \(x\sim y\) e \(y\sim z\). Luego \(P_3\) también se evita.

\medskip
En ambos casos obtenemos una solución con \(\lceil k/2\rceil\) clases que evita \(P_3\) por lo tanto,
\[
\segthree(G)\le \left\lceil \frac{k}{2}\right\rceil
=
\left\lceil \frac{\chi(G)}{2}\right\rceil.
\]
\end{proof}

\medskip
\begin{Claim}\label{cl:le-seg3}
Para todo grafo \(G\),
\[
\left\lceil \frac{\chi(G)}{2}\right\rceil \le \segthree(G).
\]
\end{Claim}

ESTA hay que sacar afuera la obs!! preguntarle bien a ericcc
lo de que cada G[Vi] es bipartito
\begin{proof}
Sea \(k=\segthree(G)\), existe una solución \((\prec,\mathcal{Q})\) 
con \(\mathcal{Q}=\{V_1,\dots,V_k\}\) que evita el patrón \(P_3\).

Para cada \(i\in\{1,\dots,k\}\) definimos
\[
L_i:=\{\,v\in V_i:\ \text{\(v\) no tiene vecinos en \(V_i\) a la derecha de \(v\) respecto de \(\prec\)}\,\},
\]
\[
R_i:=\{\,v\in V_i:\ \text{\(v\) no tiene vecinos en \(V_i\) a la izquierda de \(v\) respecto de \(\prec\)}\,\}.
\]
Los vértices aislados en \(G[V_i]\) pueden colocarse en cualquiera de los dos conjuntos.

Notemos que todo vértice de \(V_i\) pertenece a \(L_i\cup R_i\),
si existiera \(s\in V_i\) con un vecino \(r\in V_i\) a la izquierda y otro vecino \(t\in V_i\) a la derecha,
entonces la terna \(r\prec s\prec t\) no evitaría \(P_3\).

Además, \(L_i\) es independiente, si \(x\prec y\) con \(x,y\in L_i\) y \(x\sim y\), entonces \(x\) tendría
un vecino en \(V_i\) a la derecha (a saber, \(y\)), lo cual contradice la definición de \(L_i\).
Análogamente, \(R_i\) es independiente.

Por lo tanto, para cada \(i\) se tiene una partición \(V_i=L_i\,\dot\cup\,R_i\) en dos conjuntos independientes.
En consecuencia, los \(2k\) conjuntos
\[
L_1,\,R_1,\,L_2,\,R_2,\,\dots,\,L_k,\,R_k
\]
son independientes y determinan un \(2k\)-coloreo propio de \(G\). Luego
\[
\chi(G)\le 2k.
\]
En particular \(k\ge \chi(G)/2\), y como \(k\in\mathbb{N}\),
\[
k\ge \left\lceil \frac{\chi(G)}{2}\right\rceil.
\]
\end{proof}

Combinando los Claims \ref{cl:seg3-le} y \ref{cl:le-seg3} se obtiene la igualdad.
\end{proof}



SEG13



\begin{definition}[\(\segonethree\)]
Fijamos la familia de patrones
\[
\Ponethree := \{P_1,P_3\}.
\]
Dado un grafo \(G\), se define
\[
\segonethree(G)\;:=\;\minClases_{\Ponethree}(G),
\]
es decir, el mínimo número de clases de una solución \((\prec,\mathcal{Q})\) de \(G\) que evita simultáneamente
los patrones \(P_1\) y \(P_3\).
\end{definition}





\begin{theorem}\label{thm:seg13-eq}
Para todo grafo \(G\),
\[
\segthirteen(G)=\chi(G)-1.
\]
\end{theorem}

\begin{proof}
La demostración se divide en dos desigualdades.

\medskip
\begin{lemma}\label{lem:seg13-ge-chi-1}
Para todo grafo \(G\),
\[
\segthirteen(G)\ge \chi(G)-1.
\]
\end{lemma}

\begin{proof}
Por definición, \(\segthirteen\) se obtiene imponiendo simultáneamente las restricciones que definen
\(\segone\) y \(\segthree\). En particular, toda solución que evita los patrones prohibidos de
\(\segthirteen\) evita también el patrón prohibido de \(\segone\).
Luego, por el Lema \ref{lem:monotone-patterns} (monotonía al agregar patrones prohibidos),
\[
\segthirteen(G)\ge \segone(G).
\]
Finalmente, por el Teorema \ref{thm:segone-chi} se tiene \(\segone(G)=\chi(G)-1\), y por lo tanto
\[
\segthirteen(G)\ge \chi(G)-1.
\]
\end{proof}


\medskip
\begin{lemma}\label{lem:seg13-le-chi-1}
Para todo grafo \(G\),
\[
\segthirteen(G)\le \chi(G)-1.
\]
\end{lemma}

\begin{proof}
Sea \(k=\chi(G)\). Fijamos una partición basada en el k-coloreo propio
\[
V \;=\; C_1 \,\dot\cup\, C_2 \,\dot\cup\, \cdots \,\dot\cup\, C_k,
\]
donde cada \(C_i\) es independiente.

Se construye una partición de \(V\) en \(k-1\) clases definiendo
\[
V_i := C_i \quad \text{para todo } i=1,\dots,k-2,
\qquad\text{y}\qquad
V_{k-1} := C_{k-1}\cup C_k.
\]
El orden \(\prec\) se define por bloques imponiendo
\[
V_1 \prec V_2 \prec \cdots \prec V_{k-2} \prec V_{k-1},
\]
es decir, para todo \(i<j\), para todo \(u\in V_i\) y todo \(v\in V_j\), se cumple \(u\prec v\).
Además, dentro del último bloque \(V_{k-1}=C_{k-1}\cup C_k\) se fija el orden
\[
C_{k-1} \prec C_k,
\]
esto es, primero se ubican todos los vértices de \(C_{k-1}\) y luego todos los de \(C_k\).
Por último, dentro de cada \(C_i\) se toma un orden arbitrario.


\begin{figure}[t]
\centering
\begin{tikzpicture}[
  font=\small,
  oval/.style={draw, rounded corners=12pt, minimum width=16mm, minimum height=7mm, inner sep=1.5pt},
  rowlab/.style={anchor=east}
]

\def\dy{0.90cm}   % separación vertical entre clases
\def\hg{0.60cm}   % “gap” horizontal mínimo entre óvalos

% Fila 1
\node[oval, anchor=west] (C1) at (0,0) {$C_1$};
\node[rowlab] at ($(C1.west)+(-0.35cm,0)$) {$V_1$};

% Fila 2: arranca después de que termina C1
\node[oval, anchor=west] (C2) at ($(C1.east)+(\hg,-\dy)$) {$C_2$};
\node[rowlab] at ($(C2.west)+(-0.35cm,0)$) {$V_2$};

% Puntos suspensivos
\node at ($(C2.east)+(1.0cm,-\dy)$) {$\vdots$};

% Fila k-2 (posición ilustrativa, también “arranca más a la derecha”)
\node[oval, anchor=west] (Ck2) at ($(C2.east)+(2.0cm,-2*\dy)$) {$C_{k-2}$};
\node[rowlab] at ($(Ck2.west)+(-0.35cm,0)$) {$V_{k-2}$};

% Fila k-1: arranca después de que termina C_{k-2}
\node[oval, anchor=west, minimum width=18mm] (Ckm1) at ($(Ck2.east)+(\hg,-\dy)$) {$C_{k-1}$};
\node[rowlab] at ($(Ckm1.west)+(-0.35cm,0)$) {$V_{k-1}$};

% C_k: arranca después de que termina C_{k-1}
\node[oval, anchor=west, minimum width=16mm] (Ck) at ($(Ckm1.east)+(0.80cm,0)$) {$C_k$};

\node[draw=none] at ($(Ckm1.west)+(1.9cm,-0.55cm)$) {$V_{k-1}=C_{k-1}\cup C_k$};

\end{tikzpicture}
\caption{Esquema diagonal del orden por bloques.}
\label{fig:esquema-diagonal}
\end{figure}





Comprobemos que la solución \((\prec,\{V_1,\dots,V_{k-1}\})\) evita simultáneamente los patrones
\(P_1\) y \(P_3\).

\medskip
\noindent\textbf{Verifiquemos que toda terna evita \(P_1\),}
sea \(r\prec s\prec t\). Por definición de “evitar”, alcanza con ver que se cumple al menos una
de las condiciones (1)–(3) para el patrón \(P_1=(A,N,C)\), donde
\(A=\{\{1,2\},\{2,3\}\}\), \(N=\varnothing\) y \(C=\{1,2\}\).
Es decir, alcanza con que ocurra alguna de las siguientes:
\[
\text{(1) } r\not\sim s \text{ o } s\not\sim t,
\qquad
\text{(3) } r \text{ y } s \text{ quedan en clases distintas.}
\]
(La condición (2) no aplica porque \(N=\varnothing\).)

\smallskip
Si \(r\) y \(s\) están en clases distintas, entonces vale (3) y la terna evita \(P_1\),

supongamos ahora que \(r\) y \(s\) están en una misma
clase \(V_i\), si fuera que \(i\in\{1,\dots,k-2\}\), entonces \(V_i=C_i\) es independiente, luego \(r\not\sim s\) y vale (1).

Resta el caso \(i=k-1\), o sea \(r,s\in V_{k-1}=C_{k-1}\cup C_k\).
Si \(r\) y \(s\) pertenecen al mismo color, nuevamente \(r\not\sim s\) por independencia y vale (1).
Si pertenecen a colores distintos, como dentro de \(V_{k-1}\) se fijó \(C_{k-1}\prec C_k\) y \(r\prec s\),
se tiene \(r\in C_{k-1}\) y \(s\in C_k\). En ese caso, todo vértice \(t\) con \(s\prec t\) pertenece a \(C_k\),
y como \(C_k\) es independiente, se cumple \(s\not\sim t\), es decir, vale (1).

En todos los casos la terna \((r,s,t)\) evita \(P_1\).


\medskip
\noindent\textbf{Verifiquemos que toda terna evita \(P_3\)},
sea \(r\prec s\prec t\), recordemos que \(P_3=(A,N,C)\) está dado por
\[
A=\bigl\{\{1,2\},\{2,3\}\bigr\},\qquad N=\varnothing,\qquad C=\{1,3\}.
\]
Como \(N=\varnothing\), la condición (2) nunca aplica. Por lo tanto, para que \((r,s,t)\) \emph{no} evitara \(P_3\)
tendría que fallar (1) y fallar (3) simultáneamente, es decir:
\[
r\sim s,\ \ s\sim t
\qquad\text{y}\qquad
r \text{ y } t \text{ pertenecen a una misma clase } V_i.
\]
Entonces, para probar que \((r,s,t)\) evita \(P_3\), alcanza con ver que si \(r\) y \(t\) están en una misma clase,
no pueden ocurrir a la vez \(r\sim s\) y \(s\sim t\).

Supongamos \(r,t\in V_i\), como el orden \(\prec\) es por bloques, necesariamente \(s\in V_i\).

\smallskip
\emph{Caso 1: \(i\in\{1,\dots,k-2\}\).}
Aquí \(V_i=C_i\) y \(C_i\) es independiente. Luego \(r\not\sim s\) y \(s\not\sim t\), en particular no pueden darse ambas aristas.

\smallskip
\emph{Caso 2: \(i=k-1\).}
Entonces \(r,s,t\in V_{k-1}=C_{k-1}\cup C_k\).
\begin{itemize}
  \item Si \(r\) y \(t\) caen en el mismo color (ambos en \(C_{k-1}\) o ambos en \(C_k\)), entonces por el orden interno
  \(C_{k-1}\prec C_k\) también \(s\) cae en ese mismo color, y como dicho color es independiente se tiene
  \(r\not\sim s\) y \(s\not\sim t\).
  \item Si \(r\) y \(t\) caen en colores distintos, necesariamente \(r\in C_{k-1}\) y \(t\in C_k\).
  Si \(s\in C_{k-1}\), entonces \(r\not\sim s\); si \(s\in C_k\), entonces \(s\not\sim t\).
  En ambos subcasos es imposible que ocurran simultáneamente \(r\sim s\) y \(s\sim t\).
\end{itemize}

En todos los casos, cuando \(r\) y \(t\) están en la misma clase, al menos una de las aristas \(rs\) o \(st\) falta.


\end{proof}

\medskip
Combinando los Lemas \ref{lem:seg13-ge-chi-1} y \ref{lem:seg13-le-chi-1} se obtiene la igualdad.
\end{proof}



SEG123





\begin{definition}[\(\segonetwothree\)]
Se fija la familia de patrones
\[
\Ponetwothree := \{P_1,P_2,P_3\}.
\]
Dado un grafo \(G\), se define
\[
\segonetwothree(G)\;:=\;\minClases_{\Ponetwothree}(G),
\]
es decir, el mínimo número de clases de una solución \((\prec,\mathcal{Q})\) de \(G\) que evita simultáneamente
los patrones \(P_1\), \(P_2\) y \(P_3\).
\end{definition}






\german{teo que dice que las versiones completas son "lo mismo" que las uncolored}
\german{ahora me parece super obvio, cuando lo vi por primera vez con seg me re sorprendio q no exista el k=2 en 12 23 13}
\begin{theorem}\label{thm:all-colored-two-iff-base-class}
Let
\[
P=(A,N,\varnothing)
\]
be a pattern, and let
\[
P_{12}:=(A,N,\{1,2\}),\qquad
P_{23}:=(A,N,\{2,3\}),\qquad
P_{13}:=(A,N,\{1,3\})
\]
be its three colored versions. We write
\[
P_{12\text{-}23\text{-}13}(G)
\]
for the parameter induced by the family
\[
\{P_{12},P_{23},P_{13}\}.
\]

Then, for every graph \(G\),
\[
P_{12\text{-}23\text{-}13}(G)\le 2
\quad\Longleftrightarrow\quad
G\in\mathcal C_P.
\]
\end{theorem}

\begin{proof}
Suppose first that
\[
P_{12\text{-}23\text{-}13}(G)\le 2.
\]
Then there exists a solution
\[
S=(\prec,\mathcal Q),
\qquad
|\mathcal Q|\le 2,
\]
that avoids \(P_{12}\), \(P_{23}\), and \(P_{13}\).

We claim that the order \(\prec\) avoids the pattern \(P\). Suppose otherwise that some ordered triple
\[
v_1\prec v_2\prec v_3
\]
realizes \(P\). Since the triple realizes the underlying pattern, neither \edgecondition{} nor
\nonedgecondition{} holds for any of its three colored versions.

Therefore, avoiding \(P_{12}\) forces \(v_1\) and \(v_2\) to lie in different parts of \(\mathcal Q\).
Similarly, avoiding \(P_{23}\) forces \(v_2\) and \(v_3\) to lie in different parts, and avoiding \(P_{13}\)
forces \(v_1\) and \(v_3\) to lie in different parts.

Thus, \(v_1,v_2,v_3\) must lie in three pairwise distinct parts, contradicting
\[
|\mathcal Q|\le 2.
\]
Hence every ordered triple avoids \(P\), and therefore
\[
G\in\mathcal C_P.
\]

Conversely, suppose that
\[
G\in\mathcal C_P.
\]
Then there exists a total order \(\prec\) of \(V(G)\) that avoids \(P\).
Consider the one-part partition
\[
\mathcal Q=\{V(G)\}.
\]
Every ordered triple avoids \(P\) by \edgecondition{} or \nonedgecondition{}, and therefore also avoids each of
\(P_{12}\), \(P_{23}\), and \(P_{13}\).

Thus,
\[
P_{12\text{-}23\text{-}13}(G)=1,
\]
and in particular,
\[
P_{12\text{-}23\text{-}13}(G)\le 2.
\]
\end{proof}



\[
P_{12\text{-}23\text{-}13}(G)\le 2
\quad\Longleftrightarrow\quad
P_{12\text{-}23\text{-}13}(G)=1
\quad\Longleftrightarrow\quad
G\in\mathcal C_P.
\]







\german{amo los parrafos de intuicion}
\paragraph{Intuition.}
The following reduction is designed by first considering the reverse implication.
For each edge \(xy\in E(X)\) as always, we ask: "How can we force the endpoints of every edge of a graph \(X\) to lie in different parts", we add a vertex \(w_{xy}\) so that \(x,y,w_{xy}\) form a triangle.
Since the solution must avoid the three patterns
\(\operatorname{bip}_{12}\), \(\operatorname{bip}_{23}\), and \(\operatorname{bip}_{13}\),
the three vertices of this triangle must lie in pairwise different parts, regardless of their relative order.
In particular, \(x\) and \(y\) are forced to lie in different parts.





\begin{theorem}\label{thm:bip12-23-13-np-complete}
For every fixed integer \(k\ge 3\), deciding whether
\[
\operatorname{bip}_{12\text{-}23\text{-}13}(G)\le k
\]
is NP-complete.
\end{theorem}
\german{de nuevo lo de avoidance}
\begin{proof}
The problem belongs to NP: a certificate consists of a total order \(\prec\) of \(V(G)\) and a partition
\[
\mathcal Q=\{V_1,\dots,V_q\},
\qquad q\le k,
\]
and one can verify in polynomial time whether the corresponding solution avoids
\[
\mathcal B_{12\text{-}23\text{-}13}
=
\{\operatorname{bip}_{12},
  \operatorname{bip}_{23},
  \operatorname{bip}_{13}\}.
\]

We prove NP-hardness by reducing from \(k\)-\textsc{Coloring}. Let
\[
X=(V(X),E(X))
\]
be an instance of \(k\)-\textsc{Coloring}. For every edge \(xy\in E(X)\), we add a new vertex \(w_{xy}\) and make
\(x,y,w_{xy}\) induce a triangle. Thus,
\[
V(G)
=
V(X)\cup\{w_{xy}:xy\in E(X)\},
\]
and
\[
E(G)
=
E(X)\cup
\{xw_{xy},yw_{xy}:xy\in E(X)\}.
\]

We claim that
\[
X\text{ is \(k\)-colorable}
\quad\Longleftrightarrow\quad
\operatorname{bip}_{12\text{-}23\text{-}13}(G)\le k.
\]

\medskip
\noindent\textbf{(\(\Rightarrow\))}
Suppose that \(X\) admits a proper \(k\)-coloring
\[
\varphi:V(X)\to\{1,\dots,k\}.
\]
For every edge \(xy\in E(X)\), choose a color
\[
\varphi(w_{xy})
\notin
\{\varphi(x),\varphi(y)\}.
\]
Such a color exists because \(k\ge 3\) and
\(\varphi(x)\neq\varphi(y)\). This extends \(\varphi\) to a proper \(k\)-coloring of \(G\).

Let
\[
\mathcal Q=\{V_1,\dots,V_k\}
\]
be the corresponding partition into color classes. We define an order \(\prec\) by blocks:
\[
V_1\prec V_2\prec\cdots\prec V_k,
\]
taking an arbitrary order inside each block.

We show that
\[
S=(\prec,\mathcal Q)
\]
avoids the three colored patterns. Let
\[
v_1\prec v_2\prec v_3
\]
be an ordered triple. If either \(v_1v_2\notin E(G)\) or \(v_2v_3\notin E(G)\), then the triple avoids every
bipartite colored pattern by \edgecondition.

Suppose instead that
\[
v_1v_2\in E(G)
\qquad\text{and}\qquad
v_2v_3\in E(G).
\]
Since the coloring of \(G\) is proper,
\[
\varphi(v_1)\neq\varphi(v_2)
\qquad\text{and}\qquad
\varphi(v_2)\neq\varphi(v_3).
\]
Moreover, \(\varphi(v_1)\neq\varphi(v_3)\). Indeed, if
\(\varphi(v_1)=\varphi(v_3)\), then \(v_1\) and \(v_3\) would lie in the same color block. Since
\[
v_1\prec v_2\prec v_3,
\]
the vertex \(v_2\) would lie in that same block, contradicting
\(\varphi(v_1)\neq\varphi(v_2)\).

Thus \(v_1,v_2,v_3\) lie in pairwise distinct parts of \(\mathcal Q\). Therefore the triple satisfies
\colorcondition{} for each of the colored pairs
\[
\{1,2\},\qquad
\{2,3\},\qquad
\{1,3\}.
\]
Hence \(S\) avoids \(\mathcal B_{12\text{-}23\text{-}13}\), and consequently
\[
\operatorname{bip}_{12\text{-}23\text{-}13}(G)\le k.
\]

\medskip
\noindent\textbf{(\(\Leftarrow\))}
Suppose that \(G\) admits a solution
\[
S=(\prec,\mathcal Q),
\qquad
|\mathcal Q|\le k,
\]
that avoids
\(\mathcal B_{12\text{-}23\text{-}13}\).

We show that the restriction of \(\mathcal Q\) to \(V(X)\) is a proper coloring of \(X\). Let
\[
xy\in E(X).
\]
By construction, the vertices
\[
x,\quad y,\quad w_{xy}
\]
induce a triangle in \(G\). Write them in their relative order as
\[
v_1\prec v_2\prec v_3.
\]
Since they induce a triangle,
\[
v_1v_2\in E(G)
\qquad\text{and}\qquad
v_2v_3\in E(G).
\]
Therefore \edgecondition{} does not hold for any of the three bipartite colored patterns.

Because the solution avoids \(\operatorname{bip}_{12}\), the vertices \(v_1\) and \(v_2\) must lie in different
parts. Similarly, avoiding \(\operatorname{bip}_{23}\) and \(\operatorname{bip}_{13}\) implies that
\(v_2,v_3\) and \(v_1,v_3\), respectively, lie in different parts.

Thus \(x,y,w_{xy}\) lie in pairwise distinct parts of \(\mathcal Q\). In particular, \(x\) and \(y\) lie in
different parts. Since this holds for every edge \(xy\in E(X)\), the restriction of \(\mathcal Q\) to \(V(X)\)
is a proper coloring of \(X\) with at most \(k\) colors. Hence \(X\) is \(k\)-colorable.

The construction is polynomial, so the problem is NP-hard. Since it also belongs to NP, it is NP-complete.
\end{proof}





\begin{Claim}\label{cl:seg123-le-chi}
Para todo grafo \(G\), se cumple
\[
\segonetwothree(G)\le \chi(G).
\]
\end{Claim}

\begin{proof}
Sea \(k=\chi(G)\), consideremos una partición de \(V(G)\) en \(k\) conjuntos independientes,
\[
V \;=\; C_1 \,\dot\cup\, C_2 \,\dot\cup\, \cdots \,\dot\cup\, C_k.
\]
Se considera la solución \(S=(\prec,\mathcal{Q})\) definida por
\[
\mathcal{Q}:=\{V_1,\dots,V_k\}\qquad\text{donde}\qquad V_i:=C_i \ \ \text{para todo } i\in\{1,\dots,k\},
\]
y se fija un orden \(\prec\) por bloques imponiendo
\[
C_1 \prec C_2 \prec \cdots \prec C_k,
\]
con un orden arbitrario dentro de cada bloque.

Veamos que \(S\) evita simultáneamente los patrones \(P_1,P_2,P_3\),
como cada \(V_i=C_i\) es independiente, para toda terna \(r\prec s\prec t\) se tiene:
si \(r\) y \(s\) pertenecen a una misma clase, entonces \(r\not\sim s\); y si \(s\) y \(t\) pertenecen
a una misma clase, entonces \(s\not\sim t\); y si \(r\) y \(t\) pertenecen a una misma clase, entonces
\(r\not\sim t\).
En particular, en ninguna terna pueden realizarse simultáneamente las aristas exigidas por el conjunto
\(A\) del patrón junto con la igualdad de clase exigida por \(C\). Por lo tanto, toda terna evita
\(P_1\), evita \(P_2\) y evita \(P_3\), luego \(S\) evita \(\{P_1,P_2,P_3\}\).

Con esto, \(G\) admite una solución que evita \(P_1,P_2,P_3\) con \(k\) clases, luego
\[
\segonetwothree(G)\le k=\chi(G).
\]
\end{proof}

\begin{lemma}\label{lem:seg123-ge-chi}
Para todo grafo \(G\),
\[
\segonetwothree(G)\ge \chi(G)-1.
\]
\end{lemma}



\begin{proof}
Por definición, \(\segonetwothree(G)=\minClases_{\{P_1,P_2,P_3\}}(G)\) y \(\segone(G)=\minClases_{\{P_1\}}(G)\).
Como \(\{P_1\}\subseteq \{P_1,P_2,P_3\}\), el Lema \ref{lem:monotone-patterns} de monotonía implica
\[
\segonetwothree(G)\;=\;\minClases_{\{P_1,P_2,P_3\}}(G)\ \ge\ \minClases_{\{P_1\}}(G)\;=\;\segone(G).
\]
Finalmente, por el Teorema \ref{thm:segone-chi} (esto es, \(\segone(G)=\chi(G)-1\)), se concluye
\[
\segonetwothree(G)\ge \chi(G)-1.
\]
\end{proof}


Pero, para este caso de las tres condiciones juntas, No vale el >=.

\begin{lemma}\label{lem:seg123-not-ge-chi}
No es cierto que \(\segonetwothree(G)\ge \chi(G)\) para todo grafo \(G\).
\end{lemma}

\begin{proof}
Basta exhibir un contraejemplo, vamos a considerar el grafo de Mycielski \(M_4\), para el cual
\(\chi(M_4)=4\), y para el cual se tiene además una solución con \(3\) clases que evita
simultáneamente \(P_1,P_2,P_3\), es decir,
\[
\segonetwothree(M_4)\le 3.
\]
En particular,
\[
\segonetwothree(M_4)\le 3 \;<\; 4=\chi(M_4),
\]
y por lo tanto la desigualdad \(\segonetwothree(G)\ge \chi(G)\) no vale en general.
\end{proof}



\subsection{Grafos de Mycielski}

\begin{definition}[Paso de Mycielski]
Sea \(M_j=(V_j,E_j)\) un grafo, con \(V_j=\{v_0,\dots,v_{n-1}\}\).
Se construye \(M_{j+1}\) introduciendo una copia
\[
W_j:=\{w_0,\dots,w_{n-1}\},
\]
y un vértice adicional \(w\). Se fija
\[
V(M_{j+1}) \;:=\; V_j \,\dot\cup\, W_j \,\dot\cup\, \{w\}.
\]
La adyacencia en \(M_{j+1}\) se determina imponiendo, para todo \(i,\ell\in\{0,\dots,n-1\}\),
\[
v_i \sim v_\ell \ \text{en } M_{j+1}
\quad\Longleftrightarrow\quad
v_i \sim v_\ell \ \text{en } M_j,
\]
\[
w_i \sim v_\ell \ \text{en } M_{j+1}
\quad\Longleftrightarrow\quad
v_i \sim v_\ell \ \text{en } M_j,
\]
y además
\[
w \sim w_i \ \text{en } M_{j+1}\qquad\text{para todo } i=0,\dots,n-1.
\]
No se agregan otras aristas.
\end{definition}

\begin{definition}[Familia \(M_j\)]
Se fija \(M_2:=K_2\) y, para todo \(j\ge 2\), se define \(M_{j+1}\) a partir de \(M_j\)
mediante el paso anterior.
\end{definition}

% --- Colores de las clases (sin sombreado: solo borde + texto) ---
\colorlet{Czero}{red!70!black}
\colorlet{Cone}{blue!70!black}
\colorlet{Ctwo}{green!50!black}

\begin{figure}[H]
\centering
\begin{subfigure}[t]{0.40\textwidth}
\centering
\begin{tikzpicture}[
  scale=1.65,
  every node/.style={circle, draw, inner sep=1.2pt, font=\small, fill=white}
]
% M3 = C5
\foreach \i/\ang in {0/90,1/18,2/-54,3/-126,4/162}{
  \node (v\i) at (\ang:1.05) {$v_{\i}$};
}
\foreach \i/\j in {0/1,1/2,2/3,3/4,4/0}{
  \draw (v\i) -- (v\j);
}
\end{tikzpicture}
\caption{\(M_3 \cong C_5\).}
\label{fig:M3}
\end{subfigure}
\hfill
\begin{subfigure}[t]{0.58\textwidth}
\centering
\begin{tikzpicture}[
  scale=1.55,
  every node/.style={circle, inner sep=1.2pt, font=\small, fill=white},
  V0/.style={draw=Czero, text=Czero, line width=0.7pt},
  V1/.style={draw=Cone,  text=Cone,  line width=0.7pt},
  V2/.style={draw=Ctwo,  text=Ctwo,  line width=0.7pt}
]

% v0..v4 (pentágono exterior)
\node[V2] (v0) at ( 90:1.10) {$v_0$};
\node[V2] (v1) at ( 18:1.10) {$v_1$};
\node[V0] (v2) at (-54:1.10) {$v_2$};
\node[V2] (v3) at (-126:1.10) {$v_3$};
\node[V1] (v4) at (162:1.10) {$v_4$};

% w0..w4 (copias, pentágono interior)
\node[V0] (w0) at ( 90:0.62) {$w_0$};
\node[V1] (w1) at ( 18:0.62) {$w_1$};
\node[V0] (w2) at (-54:0.62) {$w_2$};
\node[V1] (w3) at (-126:0.62) {$w_3$};
\node[V1] (w4) at (162:0.62) {$w_4$};

% vértice extra w
\node[V0] (w) at (0,0) {$w$};

% aristas del C5: v_i v_{i+1}
\draw (v0)--(v1)--(v2)--(v3)--(v4)--(v0);

% aristas w_i a vecinos de v_i (en M3=C5)
\draw (w0)--(v1) (w0)--(v4);
\draw (w1)--(v0) (w1)--(v2);
\draw (w2)--(v1) (w2)--(v3);
\draw (w3)--(v2) (w3)--(v4);
\draw (w4)--(v0) (w4)--(v3);

% aristas w - w_i
\draw (w)--(w0) (w)--(w1) (w)--(w2) (w)--(w3) (w)--(w4);

\end{tikzpicture}
\caption{\(M_4\) (con la 3-partición indicada por colores).}
\label{fig:M4}
\end{subfigure}

\caption{Los grafos \(M_3\) y \(M_4\).}
\label{fig:M3M4}
\end{figure}

Recordemos que el grafo de Mycielski \(M_4\) satisface \(\chi(M_4)=4\).
En lo que sigue se trabaja con el etiquetado
\[
V(M_4)=\{v_0,\dots,v_4,w_0,\dots,w_4,w\}.
\]


\noindent\textbf{Orden.}
\[
\begin{aligned}
w &\prec w_4 \prec v_4 \prec v_3 \prec v_1 \prec w_2 \prec w_0\\
  &\prec v_2 \prec w_3 \prec w_1 \prec v_0.
\end{aligned}
\]

\noindent\textbf{3-partición.}
\[
\begin{aligned}
V_0&=\{w,\,w_2,\,w_0,\,v_2\},\\
V_1&=\{w_4,\,v_4,\,w_3,\,w_1\},\\
V_2&=\{v_3,\,v_1,\,v_0\}.
\end{aligned}
\]

\medskip


\begin{figure}[H]
\centering
\begin{tikzpicture}[
  x=0.85cm, y=0.95cm,
  every node/.style={circle, inner sep=1.2pt, font=\small, fill=white},
  V0/.style={draw=red!70!black,  text=red!70!black,  line width=0.7pt},
  V1/.style={draw=blue!70!black, text=blue!70!black, line width=0.7pt},
  V2/.style={draw=green!50!black,text=green!50!black,line width=0.7pt}
]

% Etiquetas de filas
\node[draw=none, circle=none, fill=none, anchor=east, font=\small] at (0.2,2) {$V_0$};
\node[draw=none, circle=none, fill=none, anchor=east, font=\small] at (0.2,1) {$V_1$};
\node[draw=none, circle=none, fill=none, anchor=east, font=\small] at (0.2,0) {$V_2$};

% (Opcional) marcas de posición del orden
\foreach \x in {1,...,11}{
  \draw[gray!25] (\x, -0.35) -- (\x, 2.35);
}

% Orden: [10, 9, 4, 3, 1, 7, 5, 2, 8, 6, 0]
% En notación: 10=w, 9=w_4, 4=v_4, 3=v_3, 1=v_1, 7=w_2, 5=w_0, 2=v_2, 8=w_3, 6=w_1, 0=v_0

% Fila V0 = {w, w2, w0, v2} en posiciones 1,6,7,8
\node[V0] at (1,2) {$w$};
\node[V0] at (6,2) {$w_2$};
\node[V0] at (7,2) {$w_0$};
\node[V0] at (8,2) {$v_2$};

% Fila V1 = {w4, v4, w3, w1} en posiciones 2,3,9,10
\node[V1] at (2,1) {$w_4$};
\node[V1] at (3,1) {$v_4$};
\node[V1] at (9,1) {$w_3$};
\node[V1] at (10,1) {$w_1$};

% Fila V2 = {v3, v1, v0} en posiciones 4,5,11
\node[V2] at (4,0) {$v_3$};
\node[V2] at (5,0) {$v_1$};
\node[V2] at (11,0) {$v_0$};

% (Opcional) flecha del orden
\draw[->, line width=0.6pt] (1,2.55) -- (11,2.55);
\node[draw=none, circle=none, fill=none, font=\small] at (6,2.75) {orden \(\prec\)};

\end{tikzpicture}
\caption{Orden fijado y 3-partición de \(V(M_4)\) (cada vértice en la fila de su clase).}
\label{fig:M4-order-rows}
\end{figure}


A continuación verifiquemos que la solucion anterior evita simultáneamente \(P_1,P_2,P_3\).




\begin{lemma}\label{lem:M4-avoid-P1}
El orden \(\prec\) y la partición \(V(M_4)=V_0 \,\dot\cup\, V_1 \,\dot\cup\, V_2\) definidos arriba
evitan el patrón \(P_1\).
\end{lemma}

\begin{proof}
Sea \(r\prec s\prec t\) una terna de vértices.
Recordemos que una terna \((r,s,t)\) \emph{no} evita \(P_1\) sólo si se cumplen simultáneamente:
\[
r \text{ y } s \text{ están en una misma clase},\qquad r\sim s,\qquad s\sim t.
\]
Mostremos que esto es imposible, es decir, que toda terna satisface al menos una de las condiciones
de evitación.

De la lista de adyacencias y la partición fijada, las aristas internas de las clases son:
\[
\text{en }V_0:\ w\sim w_2,\ w\sim w_0,\qquad
\text{en }V_1:\ v_4\sim w_3,\qquad
\text{en }V_2:\ v_1\sim v_0.
\]
En particular, si \(r\) y \(s\) están en una misma clase y además \(r\sim s\), entonces necesariamente
\(s\) es uno de los vértices
\[
s\in\{w_2,\ w_0,\ w_3,\ v_0\}.
\]
Verifiquemos en cada caso que \(s\) no tiene vecinos a la derecha; esto fuerza \(s\not\sim t\) y por lo
tanto \((r,s,t)\) evita \(P_1\) (por la condición (1) de evitación, con la arista \(\{2,3\}\)).

\smallskip
\noindent\textbf{(1) \(s=w_2\).}
Se tiene \(N(w_2)=\{v_1,v_3,w\}\), y en el orden fijado
\[
w \prec v_3 \prec v_1 \prec w_2.
\]
Luego todo vecino de \(w_2\) queda a la izquierda de \(w_2\), por lo que \(w_2\not\sim t\).

\smallskip
\noindent\textbf{(2) \(s=w_0\).}
Se tiene \(N(w_0)=\{v_1,v_4,w\}\), y en el orden fijado
\[
w \prec v_4 \prec v_1 \prec w_0.
\]
Luego todo vecino de \(w_0\) queda a la izquierda de \(w_0\), por lo que \(w_0\not\sim t\).

\smallskip
\noindent\textbf{(3) \(s=w_3\).}
Se tiene \(N(w_3)=\{v_2,v_4,w\}\), y en el orden fijado
\[
w \prec v_4 \prec v_2 \prec w_3.
\]
Luego todo vecino de \(w_3\) queda a la izquierda de \(w_3\), por lo que \(w_3\not\sim t\).

\smallskip
\noindent\textbf{(4) \(s=v_0\).}
En el orden fijado, \(v_0\) es el último vértice. Por lo tanto no existe \(t\) con \(s\prec t\), y la
configuración requerida por \(P_1\) es imposible.

En todos los casos, la terna \((r,s,t)\) evita \(P_1\). Esto prueba que la solución dada evita \(P_1\).
\end{proof}


\begin{lemma}\label{lem:M4-avoid-P2}
El orden \(\prec\) y la partición \(V(M_4)=V_0 \,\dot\cup\, V_1 \,\dot\cup\, V_2\) definidos arriba
evitan el patrón \(P_2\).
\end{lemma}

\begin{proof}
Sea \(r\prec s\prec t\) una terna de vértices.
Recordemos que una terna \((r,s,t)\) \emph{no} evita \(P_2\) sólo si se cumplen simultáneamente:
\[
s \text{ y } t \text{ están en una misma clase},\qquad s\sim t,\qquad r\sim s.
\]
Mostremos que esto es imposible.

De la lista de adyacencias y la partición fijada, las aristas internas de las clases son:
\[
\text{en }V_0:\ w\sim w_2,\ w\sim w_0,\qquad
\text{en }V_1:\ v_4\sim w_3,\qquad
\text{en }V_2:\ v_1\sim v_0.
\]
En particular, si \(s\) y \(t\) están en una misma clase y además \(s\sim t\), entonces necesariamente
\(s\) es el extremo izquierdo de una arista interna de su clase, es decir,
\[
s\in\{w,\ v_4,\ v_1\}.
\]
Verifiquemos en cada caso que \(s\) no tiene vecinos a la izquierda; esto fuerza \(r\not\sim s\) y por lo
tanto \((r,s,t)\) evita \(P_2\) (por la condición (1) de evitación, con la arista \(\{1,2\}\)).

\smallskip
\noindent\textbf{(1) \(s=w\).}
En el orden fijado, \(w\) es el primer vértice. Por lo tanto no existe \(r\) con \(r\prec s\), y la
configuración requerida por \(P_2\) es imposible.

\smallskip
\noindent\textbf{(2) \(s=v_4\).}
Se tiene \(N(v_4)=\{v_0,v_3,w_0,w_3\}\), y en el orden fijado
\[
w \prec v_4 \prec v_3,\qquad
v_4 \prec w_0,\qquad
v_4 \prec v_2 \prec w_3,\qquad
v_4 \prec v_0.
\]
Luego todo vecino de \(v_4\) queda a la derecha de \(v_4\), por lo que \(r\not\sim v_4\).

\smallskip
\noindent\textbf{(3) \(s=v_1\).}
Se tiene \(N(v_1)=\{v_0,v_2,w_0,w_2\}\), y en el orden fijado
\[
w \prec v_4 \prec v_3 \prec v_1 \prec w_2,\qquad
v_1 \prec w_0,\qquad
v_1 \prec v_2,\qquad
v_1 \prec v_0.
\]
Luego todo vecino de \(v_1\) queda a la derecha de \(v_1\), por lo que \(r\not\sim v_1\).

En todos los casos, la terna \((r,s,t)\) evita \(P_2\). Esto prueba que la solución dada evita \(P_2\).
\end{proof}


\begin{lemma}\label{lem:M4-seg3}
El orden y la partición \(V(M_4)=V_0 \,\dot\cup\, V_1 \,\dot\cup\, V_2\) definidos arriba
evitan el patrón \(P_3\).
\end{lemma}

\begin{proof}
Recordemos que \(P_3=(A,N,C)\) viene dado por
\[
A=\bigl\{\{1,2\},\{2,3\}\bigr\},\qquad N=\varnothing,\qquad C=\{1,3\}.
\]
Luego, una terna \(r\prec s\prec t\) \emph{no} evita \(P_3\) si y solo si se cumple
\[
r\sim s,\quad s\sim t,\quad\text{y}\quad r,t\text{ pertenecen a una misma clase }V_i.
\]
En otras palabras, alcanza con verificar que para todo vértice \(s\) las clases de sus vecinos a la izquierda
son disjuntas de las clases de sus vecinos a la derecha.

En los extremos \(w\) y \(v_0\) no hay, respectivamente, vecinos a la izquierda ni a la derecha, y no hay nada que verificar.
Para los otros vértices tenemos:

\smallskip
\noindent\textbf{(1) \(w_4\).}
A la izquierda sólo aparece \(w\in V_0\), mientras que a la derecha aparecen \(v_3,v_0\in V_2\).

\smallskip
\noindent\textbf{(2) \(v_4\).}
No tiene vecinos a la izquierda.

\smallskip
\noindent\textbf{(3) \(v_3\).}
A la izquierda aparecen \(w_4,v_4\in V_1\), y a la derecha aparecen \(w_2,v_2\in V_0\).

\smallskip
\noindent\textbf{(4) \(v_1\).}
No tiene vecinos a la izquierda.

\smallskip
\noindent\textbf{(5) \(w_2\).}
No tiene vecinos a la derecha.

\smallskip
\noindent\textbf{(6) \(w_0\).}
No tiene vecinos a la derecha.

\smallskip
\noindent\textbf{(7) \(v_2\).}
A la izquierda aparecen \(v_3,v_1\in V_2\), y a la derecha aparecen \(w_3,w_1\in V_1\).

\smallskip
\noindent\textbf{(8) \(w_3\).}
No tiene vecinos a la derecha.

\smallskip
\noindent\textbf{(9) \(w_1\).}
A la izquierda aparecen \(w,v_2\in V_0\), y a la derecha aparece \(v_0\in V_2\).
\end{proof}

EST1



\begin{definition}[Patrón prohibido para \(\estone\)]
Se fija el patrón coloreado
\[
\EStone=(A,N,C),
\]
donde
\[
A=\bigl\{\{1,2\}\bigr\},\qquad
N=\varnothing,\qquad
C=\bigl\{\{1,2\}\bigr\}.
\]
\end{definition}

\begin{figure}[ht]
\centering
\begin{tikzpicture}[scale=1, every node/.style={circle, draw, inner sep=1.2pt, font=\small}]
\node[draw=red, very thick] (a) at (0,0)   {1};
\node[draw=red, very thick] (b) at (1.2,0) {2};
\node                      (c) at (2.4,0) {3};
\draw (a)--(b);
\end{tikzpicture}
\caption{El patrón \(\mathcal{E}_1\): se exige la arista \(12\) y se exige que \(1\) y \(2\) estén en la misma clase.}
\end{figure}

\begin{definition}[\(\estone\)]
Dado un grafo \(G\), se define
\[
\estone(G)\;:=\;\minClases_{\{\mathcal{E}_1\}}(G),
\]
es decir, el mínimo número de clases de una solución \((\prec,\mathcal{Q})\) de \(G\) que evita el patrón \(\mathcal{E}_1\).
\end{definition}

\begin{Claim}\label{cl:chi-1-le-est}
Para todo grafo \(G\) se cumple \(\chi(G)-1\le \estone(G)\).
\end{Claim}


\begin{proof}
Toda solución que evita \(\mathcal E_1\) evita, en particular, el patrón \(\mathcal P_1\); por lo tanto
\(\segone(G)\le \estone(G)\).
Como \(\segone(G)=\chi(G)-1\), se concluye \(\chi(G)-1\le \estone(G)\).
\end{proof}



\begin{lemma}\label{lem:est1-le-chi}
Para todo grafo \(G\) se cumple
\[
\estone(G)\le \chi(G).
\]
\end{lemma}

\begin{proof}
Sea \(k=\chi(G)\). Se fija una partición
\[
V(G)=C_1\,\dot\cup\,\cdots\,\dot\cup\,C_k
\]
tal que cada \(C_i\) es independiente, y se toma un orden \(\prec\) arbitrario sobre \(V(G)\).
Sea \(\mathcal{Q}:=\{C_1,\dots,C_k\}\) y considérese la solución \(S=(\prec,\mathcal{Q})\).

Como no hay aristas dentro de ninguna clase \(C_i\), toda terna de vértices distintos evita el patrón \(E_1\).
Por lo tanto, \(S\) evita \(E_1\) y tiene valor \(k\), luego \(\estone(G)\le k=\chi(G)\).
\end{proof}

% -----------

% Caracterización del caso est1(G)=chi(G)
% ------------------------------------------------------------
\begin{theorem}[Caracterización del caso \(\estone(G)=\chi(G)\)]\label{thm:est1-chi}
Sea \(G\) un grafo y sea \(k=\estone(G)\). Son equivalentes:
\begin{enumerate}
\item \(k=\chi(G)\).
\item Existe una solución \(S=(\prec,\mathcal{Q})\) de valor \(k\), con
\[
\mathcal{Q}=\{C_1,\dots,C_k\}
\qquad\text{y}\qquad
V(G)=C_1\,\dot\cup\,\cdots\,\dot\cup\,C_k,
\]
tal que \(S\) evita el patrón \(E_1\) y además, si \(v_n\) es el último vértice del orden \(\prec\),
existe un indice j pert a 1,,,k tal que (PREGUNTAR ESTOO)
\[
N(v_n)\cap C_j=\varnothing.
\]
\end{enumerate}
\end{theorem}

\begin{proof}
\noindent(1 \(\Rightarrow\) 2).
Si \(\estone(G)=\chi(G)\), entonces \(\chi(G)=k\). En la demostración del
Lema~\ref{lem:est1-le-chi} se construye una solución \(S=(\prec,\{C_1,\dots,C_k\})\) que evita \(E_1\)
y tiene valor \(k\). Además, como cada \(C_i\) es independiente, si \(v_n\in C_j\) se cumple
\(N(v_n)\cap C_j=\varnothing\).

\medskip
\noindent(2 \(\Rightarrow\) 1).
Sea \(S=(\prec,\mathcal{Q})\) como en (2), con \(\mathcal{Q}=\{C_1,\dots,C_k\}\),
probemos que cada \(C_j\) es independiente, supongamos \(x,y\in C_j\) con \(x\prec y\) y \(x\sim y\).

Si \(y\neq v_n\), existe \(t\) con \(y\prec t\). Entonces la terna \((x,y,t)\) no evita \(E_1\)
(pues \(x\sim y\) y \(x,y\) están en la misma clase), contradicción.
Luego necesariamente \(y=v_n\). Pero por hipótesis \(N(v_n)\cap C_j=\varnothing\), así que \(x\not\sim v_n\),
otra contradicción. Por lo tanto, no hay aristas dentro de \(C_j\), y \(C_j\) es independiente.

Así, \(\mathcal{Q}\) es una partición de \(V(G)\) en \(k\) conjuntos independientes, luego \(\chi(G)\le k\).
Como siempre vale \(\estone(G)=k\le \chi(G)\) por el Lema~\ref{lem:est1-le-chi}, se concluye
\(\chi(G)=k=\estone(G)\).
\end{proof}




% ------------------------------------------------------------
% Corolario (versión simple)
% ------------------------------------------------------------
\begin{corollary}\label{cor:est1-chi-1-simple}
Sea \(G\) un grafo, si en toda solución óptima \(S=(\prec,\mathcal S)\) (para \(\{E_1\}\)),
si llamamos \(v_n\) al último vértice del orden \(\prec\) y \(v_n\in C\in\mathcal S\), se cumple
\[
N(v_n)\cap C\neq\varnothing,
\]
entonces
\[
\estone(G)=\chi(G)-1.
\]
\end{corollary}

\begin{proof}
Por hipótesis no existe una solución óptima que cumpla \(N(v_n)\cap C=\varnothing\).
Entonces, por el Teorema~\ref{thm:est1-chi}, no puede ocurrir \(\estone(G)=\chi(G)\).

Como siempre se tiene \(\chi(G)-1\le \estone(G)\le \chi(G)\), se concluye necesariamente
\(\estone(G)=\chi(G)-1\).
\end{proof}





\paragraph{Intuition.}
In a solution avoiding \(\operatorname{star}_{12}\), an edge whose endpoints lie in the same part can only have its
right endpoint as the last vertex of the order. Indeed, if
\[
x\prec y,\qquad xy\in E(G),
\]
and \(x\) and \(y\) lie in the same part, then any vertex \(z\) with \(y\prec z\) would produce an ordered triple
\[
x\prec y\prec z
\]
that satisfies neither \edgecondition{} nor \colorcondition{} for \(\operatorname{star}_{12}\).
\german{por ej aca! lo de nonedge}
Therefore, to force the endpoints of every edge of a graph \(X\) to lie in different parts, it is enough to force
the last vertex of the order to lie outside \(V(X)\). We achieve this by taking the disjoint union of \(X\) with
a clique \(K_{k+1}\). \german{pense q seria copado poner estos textos de intuicion que son 
la forma en que pense la vuelta del sii, que opinas sensei??, podria ponerlos en todas las reducciones me parece super interesante
no tanto para los erics, pero si para los germancitos}

\begin{theorem}\label{thm:star12-np-complete}
For every fixed integer \(k\ge 3\), deciding whether
\[
\operatorname{star}_{12}(G)\le k
\]
is NP-complete.
\end{theorem}
\german{si las cajitas de los problemas van arriba de todo aca podria decir avoidance problem}
\begin{proof}
The problem belongs to NP: a certificate consists of an order \(\prec\) of \(V(G)\) and a partition of \(V(G)\)
into at most \(k\) parts, and one can verify in polynomial time that every ordered triple avoids
\(\operatorname{star}_{12}\).

We prove NP-hardness by reducing from \(k\)-\textsc{Coloring}. Let \(X\) be an instance of
\(k\)-\textsc{Coloring}, and construct
\[
G:=X\,\dot\cup\,K_{k+1}.
\]
Thus \(G\) is the disjoint union of \(X\) and a clique \(K_{k+1}\), there are no edges between
\(V(X)\) and \(V(K_{k+1})\).

We claim that
\[
X\text{ is \(k\)-colorable}
\quad\Longleftrightarrow\quad
\minClases_{\operatorname{star}_{12}}(G)\le k.
\]

\medskip
\noindent\textbf{(\(\Rightarrow\))}
Suppose that \(X\) has a proper \(k\)-coloring. Let
\[
V(K_{k+1})=\{z_1,\dots,z_{k+1}\}.
\]
We keep the given coloring on \(V(X)\), assign \(z_i\) to part \(i\) for \(i=1,\dots,k\), and assign
\(z_{k+1}\) to the same part as \(z_1\) (the important thing is to assign a different part for the first k+1 vertices, and the only 
one that can repeat part is the last one). Hence \(z_1z_{k+1}\) is the only edge of \(K_{k+1}\) whose endpoints
lie in the same part.

Choose any order of \(V(G)\setminus\{z_{k+1}\}\), and place \(z_{k+1}\) last. We show that the resulting solution
avoids \(\operatorname{star}_{12}\).

Let
\[
v_1\prec v_2\prec v_3
\]
be an ordered triple. If \(v_1v_2\notin E(G)\), then the triple avoids the pattern by \edgecondition.
\german{esta duda la tengo siemrpe, si es mejor empezar asi las demos de que se avoidea
o directamente tomar una terna tal que esten presentes las aristas del patron??}
Suppose that \(v_1v_2\in E(G)\). Since \(G\) is a disjoint union, the vertices \(v_1,v_2\) are either both in
\(V(X)\) or both in \(V(K_{k+1})\).

If \(v_1,v_2\in V(X)\), they lie in different parts because the coloring of \(X\) is proper. If
\(v_1,v_2\in V(K_{k+1})\), they also lie in different parts: the only monochromatic edge of the clique is
\(z_1z_{k+1}\), but \(z_{k+1}\) is the last vertex of the whole order and therefore cannot occupy the second
position of an ordered triple. Thus, in both cases, the triple avoids the pattern by \colorcondition.

Therefore,
\[
\minClases_{\operatorname{star}_{12}}(G)\le k.
\]

\medskip
\noindent\textbf{(\(\Leftarrow\))}
Suppose that \(G\) admits a solution
\[
S=(\prec,\mathcal Q)
\]
with at most \(k\) parts that avoids \(\operatorname{star}_{12}\).

Since \(K_{k+1}\) has \(k+1\) vertices and \(\mathcal Q\) has at most \(k\) parts, two vertices of the clique,
say \(a\prec b\), lie in the same part \german{podria dibujar un zapdos aca (?}. As \(K_{k+1}\) is a clique, \(ab\in E(G)\). We claim that \(b\) is the
last vertex of the order. Otherwise, there would exist a vertex \(c\) such that
\[
a\prec b\prec c.
\]
The triple \((a,b,c)\) would satisfy neither \edgecondition, because \(ab\in E(G)\), nor \colorcondition, because
\(a\) and \(b\) lie in the same part. This would contradict that \(S\) avoids
\(\operatorname{star}_{12}\).

Hence the last vertex of the order belongs to \(K_{k+1}\). Let \(z\in V(K_{k+1})\) denote this last vertex.

We show that the restriction of \(\mathcal Q\) to \(V(X)\) is a proper coloring of \(X\). Let \(xy\in E(X)\), and
assume without loss of generality that \(x\prec y\). Since \(z\notin V(X)\) and \(z\) is the last vertex,
\[
x\prec y\prec z.
\]
If \(x\) and \(y\) were assigned to the same part, then the ordered triple \((x,y,z)\) would satisfy neither
\edgecondition{} nor \colorcondition{} for \(\operatorname{star}_{12}\), a contradiction.

Therefore, the endpoints of every edge of \(X\) lie in different parts of \(\mathcal Q\). The restriction of
\(\mathcal Q\) to \(V(X)\) is consequently a proper coloring of \(X\) with at most \(k\) colors. Thus \(X\) is
\(k\)-colorable.

The construction is polynomial, so the problem is NP-hard. Since it also belongs to NP, it is NP-complete.
\end{proof}






% --- Patrón prohibido para est2 ---
\newcommand{\esttwo}{\operatorname{est}_2}
\newcommand{\EpsTwo}{\mathcal{E}_2} % patrón de est2

\begin{definition}[Patrón prohibido para \(\esttwo\)]
Se fija el patrón coloreado
\[
\EpsTwo=(A,N,C),
\]
donde
\[
A=\bigl\{\{1,2\}\bigr\},
\qquad
N=\varnothing,
\qquad
C=\{2,3\}.
\]
\end{definition}

\begin{figure}[ht]
\centering
\begin{tikzpicture}[scale=1, every node/.style={circle, draw, inner sep=1.2pt, font=\small}]
\node (a) at (0,0) {1};
\node (b) at (1.4,0) {2};
\node[draw=red, very thick, fill=red!15] (c) at (2.8,0) {3};

% arista exigida 12
\draw (a)--(b);

% "mismo color" entre 2 y 3: marco ambos
\node[draw=red, very thick, fill=red!15] (b2) at (1.4,0) {2};

\end{tikzpicture}
\end{figure}

\begin{definition}[\(\esttwo\)]
Dado un grafo \(G\), se define
\[
\esttwo(G)\;:=\;\minClases_{\{\EpsTwo\}}(G),
\]
es decir, el mínimo número de clases de una solución \((\prec,\mathcal Q)\) de \(G\) que evita el patrón \(\EpsTwo\).
\end{definition}


\begin{theorem}\label{thm:est2-vc}
Sea \(G=(V,E)\) un grafo y sea \(k\in\N\), \(\esttwo(G)\) = k si y solo si
existe un conjunto \(L\subseteq V\) con \(|L|=k\) tal que \(V\setminus L\) es un conjunto independiente.
\end{theorem}

\begin{proof}
\noindent(\(\Rightarrow\))
Sea \(S=(\prec,\mathcal Q)\) una solución  de valor \(k\) que evita el patrón \(\EpsTwo\),
donde \(\mathcal Q=\{V_1,\dots,V_k\}\) es una partición de \(V\) en \(k\) clases.
Para cada \(i\in\{1,\dots,k\}\), sea \(u_i\) el vértice más a la derecha de \(V_i\) en el orden \(\prec\).
Definimos
\[
L\coloneq\{u_1,\dots,u_k\}.
\]
Entonces \(|L|=k\), veamos que \(V\setminus L\) es independiente,
supongamos por absurdo que existen \(x,y\in V\setminus L\) con \(x\sim y\),
con \(x\prec y\).
Como \(y\notin L\), no es el último de su clase, por lo tanto existe \(t\in V\) en la misma clase que \(y\)
tal que \(y\prec t\).
Consideremos la terna \(x\prec y\prec t\), se cumple \(y\sim x\), y además \(y\) y \(t\) están en la misma clase,
luego la terna \((x,y,t)\) no evita el patrón \(\EpsTwo\), contradicción con que \(S\) lo evita.
Por lo tanto \(V\setminus L\) es independiente.

REVISAR LA de abajo


\medskip
\noindent(\(\Leftarrow\))
Supongamos que existe \(L\subseteq V\) con \(|L|=k\) tal que \(V\setminus L\) es independiente.
Escribamos \(L=\{u_1,\dots,u_k\}\),
fijemos un orden \(\prec\) que ubica primero todos los vértices de \(V\setminus L\) (en cualquier orden),
y luego todos los vértices de \(L\) (en cualquier orden).

Definamos una partición \(\mathcal Q=\{V_1,\dots,V_k\}\) de \(V\) en \(k\) clases del siguiente modo,
para cada \(i\), se toma
\[
V_i :=  X_i \cup \{u_i\},
\]
donde \(X_1,\dots,X_k\) es una partición cualquiera de \(V\setminus L\).

Verifiquemos que la solución \(S=(\prec,\mathcal Q)\) evita \(\mathcal E_2\)
sea \(r\prec s\prec t\) una terna de vértices.

Recordemos que \(\mathcal E_2=(A,N,C)\) con \(A=\{\{1,2\}\}\), \(N=\varnothing\) y \(C=\{2,3\}\).
Por la definición de \emph{evitar un patrón}, la terna \((r,s,t)\) evita \(\mathcal E_2\) si se cumple
alguna de las siguientes condiciones:
\begin{enumerate}
\item[(1)] existe alguna arista en \(A\) que no esta en E (aquí: \(r\not\sim s\));
\item[(3)] el par indicado por \(C\) no queda en la misma clase (aquí: \(s\) y \(t\) quedan en clases distintas).
\end{enumerate}

Si \(s\) y \(t\) pertenecen a clases distintas de \(\mathcal Q\), entonces se cumple (3) y \((r,s,t)\)
evita \(\mathcal E_2\), supongamos entonces que \(s\) y \(t\) pertenecen a una misma clase \(V_i\in\mathcal Q\).


Como \(s\prec t\) y ambos están en \(V_i\), se deduce que \(s\neq u_i\), es decir,
\(s\notin L\). Además, de \(r\prec s\) y del hecho de que todos los vértices de \(L\) aparecen al final de \(\prec\),
se concluye también que \(r\notin L\), luego \(r,s\in V\setminus L\), y como \(V\setminus L\) es independiente,
se tiene \(r\not\sim s\), por lo tanto se cumple (1), y \((r,s,t)\) evita \(\mathcal E_2\).
\end{proof}



\begin{definition}[Orden entre patrones]\label{def:pat-order}
Sean \(P=(A,N,C)\) y \(P'=(A',N',C')\) patrones coloreados.
Se escribe \(P\preceq P'\) si
\[
A\subseteq A'
\qquad\text{y}\qquad
C\subseteq C'.
\]
\end{definition}

\begin{lemma}[Monotonicidad respecto de \(\preceq\)]\label{lem:pat-monotone}
Si \(P\preceq P'\), entonces para todo grafo \(G\) y toda solución \(S=(\prec,\mathcal Q)\) se cumple:
\[
S \text{ evita } P \;\Rightarrow\; S \text{ evita } P'.
\]
En particular,
\[
\minClases_{P'}(G)\le \minClases_{P}(G).
\]
\end{lemma}

\begin{proof}
Supongamos \(P\preceq P'\), es decir, \(A\subseteq A'\) y \(C\subseteq C'\).
Sea \(S=(\prec,\mathcal Q)\) una solución que evita \(P\). Hay que ver que evita \(P'\).

Sea \(x\prec y\prec z\) una terna de vértices distintos y escribamos \((v_1,v_2,v_3)=(x,y,z)\).
Como \(S\) evita \(P\), la terna \((v_1,v_2,v_3)\) evita \(P\). Por definición, ocurre (1) o (3):

\smallskip
\noindent\textbf{Caso 1.} Existe \(\{i,j\}\in A\) tal que \(v_i\not\sim v_j\).
Como \(A\subseteq A'\), se tiene \(\{i,j\}\in A'\). Luego se cumple (1) para \(P'\), y la terna evita \(P'\).

\smallskip
\noindent\textbf{Caso 2.} Se tiene que \(v_i\) y \(v_j\) están en clases distintas para el par \(\{i,j\}=C\).
Como \(C\subseteq C'\), se tiene \(\{i,j\}\in C'\). Luego se cumple (3) para \(P'\), y la terna evita \(P'\).

\smallskip
En ambos casos, toda terna \(x\prec y\prec z\) evita \(P'\). Por lo tanto, \(S\) evita \(P'\).

Para la desigualdad de parámetros, sea \(k=\minClases_{\{P\}}(G)\). Entonces existe una solución con \(k\) clases
que evita \(P\). Por lo anterior, esa misma solución evita \(P'\), y en consecuencia
\(\minClases_{\{P'\}}(G)\le k=\minClases_{\{P\}}(G)\).
\end{proof}


\newcommand{\estthree}{\operatorname{est}_3}
\newcommand{\mathcalEthree}{\mathcal{E}_3}

\begin{definition}[Patrón prohibido para \(\estthree\)]
Se fija el patrón coloreado
\[
\mathcalEthree=(A,N,C),
\]
donde
\[
A=\bigl\{\{1,2\}\bigr\},\qquad
N=\varnothing,\qquad
C=\{1,3\}.
\]
\end{definition}

\begin{figure}[ht]
\centering
\begin{tikzpicture}[scale=1, every node/.style={circle, draw, inner sep=1.2pt, font=\small}]
\node[draw=red, very thick, fill=red!15] (a) at (0,0) {1};
\node (b) at (1.4,0) {2};
\node[draw=red, very thick, fill=red!15] (c) at (2.8,0) {3};

% arista exigida 12
\draw (a)--(b);
\end{tikzpicture}
\end{figure}

\begin{definition}[\(\estthree\)]
Dado un grafo \(G\), se define
\[
\estthree(G)\;:=\;\minClases_{\{\mathcalEthree\}}(G),
\]
es decir, el mínimo número de clases de una solución \((\prec,\mathcal Q)\) de \(G\) que evita el patrón \(\mathcalEthree\).
\end{definition}


\german{este paraemtro ya lo definimos en otro lado}
\begin{definition}[Parámetro \(\sigma\)]
Sea \(G=(V,E)\) un grafo. Se define
\[
\sigma(G)
:=
\min\left\{
k:\ 
V(G)=C_1\,\dot\cup\,\cdots\,\dot\cup\,C_k
\ \text{ y }\
\tau\bigl(G[C_i]\bigr)\le 1
\ \text{ para todo } i=1,\dots,k
\right\},
\]
donde \(\tau(H)\) denota el tamaño mínimo de un vertex cover de \(H\).
\end{definition}

\begin{remark}
La condición
\[
\tau\bigl(G[C_i]\bigr)\le 1
\]
significa que, dentro de la clase \(C_i\), todas las aristas internas están incidentes a un mismo vértice.
Equivalentemente, \(G[C_i]\) es una estrella, posiblemente con vértices aislados, o bien un grafo sin aristas.
\end{remark}

\begin{theorem}[Caracterización de \(\estthree\)]
Para todo grafo \(G\), se cumple
\[
\estthree(G)=\sigma(G).
\]
\end{theorem}
\german{esto se borra no? esta cubierto por 13 o rt o esta buena la demo?}

\begin{proof}
Se prueban ambas desigualdades.

\medskip
\noindent\textbf{Primero, \(\sigma(G)\le \estthree(G)\).}
Sea \(k=\estthree(G)\). Entonces existe una solución
\[
S=(\prec,\mathcal Q),
\qquad
\mathcal Q=\{V_1,\dots,V_k\},
\]
de valor \(k\) que evita el patrón \(\mathcal E_3\).

Para cada \(i\in\{1,\dots,k\}\), sea \(c_i\) el último vértice de \(V_i\) en el orden \(\prec\).
Vamos a ver que
\[
\tau(G[V_i])\le 1
\qquad\text{para todo } i.
\]
Equivale a probar que \(V_i\setminus\{c_i\}\) es independiente.

Fijemos \(i\). Supongamos que existen \(x,y\in V_i\setminus\{c_i\}\) tales que \(x\sim y\).
Sin pérdida de generalidad, tomamos \(x\prec y\).
Como \(y\neq c_i\), el vértice \(y\) no es el último de \(V_i\); por lo tanto existe \(z\in V_i\) tal que
\[
y\prec z.
\]
Entonces \(x\prec y\prec z\), con \(x\sim y\) y \(x,z\in V_i\).
Por lo tanto, la terna \((x,y,z)\) no evita \(\mathcal E_3\), contradicción.

Luego \(V_i\setminus\{c_i\}\) es independiente. Así, \(c_i\) cubre todas las aristas de \(G[V_i]\), y por lo tanto
\[
\tau(G[V_i])\le 1.
\]
Como esto vale para todo \(i\), la partición \(\mathcal Q=\{V_1,\dots,V_k\}\) testifica que
\[
\sigma(G)\le k=\estthree(G).
\]

\medskip
\noindent\textbf{Ahora, \(\estthree(G)\le \sigma(G)\).}
Sea \(k=\sigma(G)\). Entonces existe una partición
\[
V(G)=C_1\,\dot\cup\,\cdots\,\dot\cup\,C_k
\]
tal que
\[
\tau(G[C_i])\le 1
\qquad\text{para todo } i=1,\dots,k.
\]
Para cada \(i\), elegimos un vértice \(c_i\in C_i\) que cubre todas las aristas de \(G[C_i]\).
Si \(G[C_i]\) no tiene aristas, elegimos \(c_i\in C_i\) arbitrariamente.
En todos los casos se cumple que
\[
C_i\setminus\{c_i\}
\]
es independiente.

Tomamos como partición de la solución
\[
\mathcal Q:=\{C_1,\dots,C_k\}.
\]
Definimos un orden \(\prec\) por bloques:
\[
C_1\prec C_2\prec \cdots \prec C_k,
\]
y dentro de cada bloque \(C_i\) ubicamos al vértice \(c_i\) al final. El resto de los vértices de \(C_i\)
se ordena arbitrariamente.

Verifiquemos que la solución \(S=(\prec,\mathcal Q)\) evita \(\mathcal E_3\).
Sea \(r\prec s\prec t\) una terna de vértices distintos.
Para evitar \(\mathcal E_3\), alcanza con que ocurra alguna de estas dos cosas:
\[
r\not\sim s
\qquad\text{o}\qquad
r \text{ y } t \text{ están en clases distintas de } \mathcal Q.
\]

Si \(r\) y \(t\) están en clases distintas, entonces la terna evita \(\mathcal E_3\).
Supongamos entonces que \(r,t\in C_i\) para alguna clase \(C_i\).
Como el orden es por bloques, de \(r\prec s\prec t\) se deduce que también \(s\in C_i\).

Además, como \(s\prec t\), el vértice \(s\) no puede ser el último de \(C_i\). Por construcción, el último vértice
de \(C_i\) es \(c_i\), luego \(s\neq c_i\). También \(r\neq c_i\), porque \(r\prec s\).
Por lo tanto,
\[
r,s\in C_i\setminus\{c_i\}.
\]
Como \(C_i\setminus\{c_i\}\) es independiente, se tiene \(r\not\sim s\).
Entonces la terna \((r,s,t)\) evita \(\mathcal E_3\).

Así, toda terna evita \(\mathcal E_3\), y por lo tanto \(S\) es una solución de valor \(k\) que evita \(\mathcal E_3\).
En consecuencia,
\[
\estthree(G)\le k=\sigma(G).
\]

Combinando ambas desigualdades, se obtiene
\[
\estthree(G)=\sigma(G).
\]
\end{proof}


\begin{theorem}\label{thm:sigma-np-complete}
Para todo \(k\ge 3\), decidir si
\[
\sigma(G)\le k
\]
es NP-completo. En particular, calcular \(\sigma(G)\) es NP-hard.
\end{theorem}

\begin{proof}
La reducción es desde \(k\)-\textsc{Coloring}, que es NP-completo para todo \(k\ge 3\).

Sea \(X\) una instancia de \(k\)-\textsc{Coloring}. Construimos el grafo
\[
G \;:=\; X \,\dot\cup\, K_{2k},
\]
es decir, la unión disjunta de \(X\) con una clique de tamaño \(2k\).

Vamos a probar que
\[
X \text{ es \(k\)-coloreable}
\quad\Longleftrightarrow\quad
\sigma(G)\le k.
\]

\medskip
\noindent\textbf{(\(\Rightarrow\))} Supongamos que \(X\) es \(k\)-coloreable. Entonces existe una partición
\[
V(X)=I_1\,\dot\cup\,\cdots\,\dot\cup\,I_k
\]
tal que cada \(I_i\) es independiente en \(X\).

Por otro lado, particionamos los \(2k\) vértices de \(K_{2k}\) en \(k\) pares
\[
P_1,\dots,P_k.
\]
Cada \(P_i\) induce una arista.

Definimos
\[
C_i:=I_i\cup P_i
\qquad\text{para } i=1,\dots,k.
\]
Como \(G\) es la unión disjunta de \(X\) y \(K_{2k}\), no hay aristas entre \(I_i\) y \(P_i\).
Además, \(I_i\) es independiente y \(P_i\) induce una única arista. Por lo tanto, \(G[C_i]\) es una arista
junto con vértices aislados, y entonces
\[
\tau(G[C_i])\le 1.
\]
Así, la partición
\[
V(G)=C_1\,\dot\cup\,\cdots\,\dot\cup\,C_k
\]
muestra que \(\sigma(G)\le k\).

\medskip
\noindent\textbf{(\(\Leftarrow\))} Supongamos ahora que \(\sigma(G)\le k\). Entonces existe una partición
\[
V(G)=C_1\,\dot\cup\,\cdots\,\dot\cup\,C_q
\]
con \(q\le k\), tal que
\[
\tau(G[C_i])\le 1
\qquad\text{para todo } i=1,\dots,q.
\]

Miremos la clique \(K_{2k}\). Una clase \(C_i\) no puede contener tres vértices de \(K_{2k}\), porque tres vértices
de una clique inducen un \(K_3\), y
\[
\tau(K_3)=2.
\]
Por lo tanto, cada clase contiene a lo sumo dos vértices de \(K_{2k}\).

Como \(K_{2k}\) tiene \(2k\) vértices y hay \(q\le k\) clases, se deduce que necesariamente \(q=k\) y que cada clase
contiene exactamente dos vértices de \(K_{2k}\). En particular, cada clase \(C_i\) contiene una arista proveniente
de la clique.

Ahora veamos qué ocurre con los vértices de \(X\). Si para alguna clase \(C_i\) existiera una arista \(xy\) de \(X\)
con \(x,y\in C_i\), entonces \(G[C_i]\) contendría dos aristas disjuntas:
\[
\text{una arista dentro de }K_{2k}
\qquad\text{y}\qquad
xy \text{ dentro de }X.
\]
Dos aristas disjuntas no pueden cubrirse con un solo vértice, luego
\[
\tau(G[C_i])\ge 2,
\]
contradiciendo la elección de la partición.

Por lo tanto, para cada \(i\), el conjunto \(C_i\cap V(X)\) es independiente en \(X\). En consecuencia,
\[
V(X)=\bigl(C_1\cap V(X)\bigr)\,\dot\cup\,\cdots\,\dot\cup\,\bigl(C_k\cap V(X)\bigr)
\]
define un \(k\)-coloreo propio de \(X\). Luego \(X\) es \(k\)-coloreable.

\medskip
Hemos probado que
\[
X \text{ es \(k\)-coloreable}
\quad\Longleftrightarrow\quad
\sigma\bigl(X\,\dot\cup\,K_{2k}\bigr)\le k.
\]
Por lo tanto, decidir si \(\sigma(G)\le k\) es NP-hard para todo \(k\ge 3\).

Finalmente, el problema está en NP: dada una partición de \(V(G)\) en a lo sumo \(k\) clases, se puede verificar en
tiempo polinomial si para cada clase \(C_i\) se cumple \(\tau(G[C_i])\le 1\), chequeando si todas las aristas de
\(G[C_i]\) están incidentes a un mismo vértice.

En conclusión, para todo \(k\ge 3\), decidir si \(\sigma(G)\le k\) es NP-completo.
\end{proof}



\german{creo q en algun lugar dijimos o tenemos q decir que para todos estos parametros decidir si su P-parametro <=2 es  poli}
\begin{theorem}\label{thm:sigma-dos-poly}
El problema de decidir si \(\sigma(G)\le 2\) se puede resolver en tiempo polinomial.
\end{theorem}

\begin{proof}
Consideramos el siguiente algoritmo.

\begin{center}
\begin{tabbing}
\hspace{1cm}\=\hspace{1cm}\=\hspace{1cm}\=\kill
\textsc{Decidir-\(\sigma\)-dos}(\(G\)) \\
\> \textbf{para cada} \(S\subseteq V(G)\) con \(|S|\le 2\) \textbf{hacer} \\
\>\> \textbf{si} \(G-S\) es bipartito \textbf{entonces} \\
\>\>\> \textbf{devolver} \textbf{true} \\
\> \textbf{devolver} \textbf{false}
\end{tabbing}
\end{center}

El algoritmo es polinomial, ya que hay
\[
1+|V(G)|+\binom{|V(G)|}{2}
\]
subconjuntos \(S\subseteq V(G)\) de tamaño a lo sumo \(2\), y para cada uno se puede verificar en tiempo lineal
si \(G-S\) es bipartito mediante BFS o DFS.

Veamos ahora que el algoritmo es correcto.

\medskip
\noindent\textbf{(\(\Rightarrow\))} Supongamos que \(\sigma(G)\le 2\). Entonces existe una partición
\[
V(G)=V_1\,\dot\cup\,V_2
\]
tal que
\[
\tau(G[V_i])\le 1
\qquad\text{para } i=1,2.
\]
Para cada \(i\in\{1,2\}\), tomamos un conjunto \(S_i\subseteq V_i\) que cubre todas las aristas de \(G[V_i]\) y
tal que \(|S_i|\le 1\). Definimos
\[
S:=S_1\cup S_2.
\]
Entonces \(|S|\le 2\). Además, como \(S_i\) cubre todas las aristas internas de \(G[V_i]\), el conjunto
\(V_i\setminus S\) es independiente para \(i=1,2\). Por lo tanto,
\[
G-S
\]
es bipartito, con bipartición
\[
(V_1\setminus S)\,\dot\cup\,(V_2\setminus S).
\]
Luego el algoritmo encuentra un conjunto \(S\) de tamaño a lo sumo \(2\) tal que \(G-S\) es bipartito, y por lo tanto
devuelve \textbf{true}.

\medskip
\noindent\textbf{(\(\Leftarrow\))} Supongamos ahora que existe \(S\subseteq V(G)\) con \(|S|\le 2\) tal que \(G-S\)
es bipartito. Sea
\[
V(G)\setminus S = A\,\dot\cup\,B
\]
una bipartición de \(G-S\), donde \(A\) y \(B\) son independientes.

Como \(|S|\le 2\), podemos escribir
\[
S=S_1\,\dot\cup\,S_2
\]
con \(|S_1|\le 1\) y \(|S_2|\le 1\). Definimos
\[
V_1:=A\cup S_1,
\qquad
V_2:=B\cup S_2.
\]
Entonces
\[
V(G)=V_1\,\dot\cup\,V_2.
\]

Veamos que cada clase induce un grafo con vertex cover de tamaño a lo sumo \(1\).
Como \(A\) es independiente y \(|S_1|\le 1\), todas las aristas internas de \(G[V_1]\), si existen, están incidentes
al único vértice de \(S_1\). Por lo tanto,
\[
\tau(G[V_1])\le 1.
\]
Del mismo modo, como \(B\) es independiente y \(|S_2|\le 1\), se tiene
\[
\tau(G[V_2])\le 1.
\]
Así, la partición \(V(G)=V_1\,\dot\cup\,V_2\) testifica que
\[
\sigma(G)\le 2.
\]

Concluimos que el algoritmo devuelve \textbf{true} si y solo si \(\sigma(G)\le 2\). Por lo tanto, decidir si
\(\sigma(G)\le 2\) se puede hacer en tiempo polinomial.
\end{proof}


%\begin{definition}[Semi-coloreo]
%Sea \(G=(V,E)\) un grafo. Un \emph{semi-coloreo} de \(G\) con \(k\) colores es una partición
%\[
%V \;=\; C_1 \,\dot\cup\, \cdots \,\dot\cup\, C_k
%\]
%tal que, para todo \(i\in\{1,\dots,k\}\), a lo sumo un vértice de \(C_i\) tiene algún vecino dentro de \(C_i\).
%Equivalentemente,
%\[
%\bigl|\{\,v\in C_i:\ N(v)\cap C_i\neq\varnothing\,\}\bigr|\le 1
%\qquad\text{para todo } i=1,\dots,k.
%\]
%Se define \(\chi_{\mathrm{semi}}(G)\) como el mínimo \(k\in\N\) para el cual \(G\) admite un semi-coloreo con \(k\) colores.
%\end{definition}




%\begin{theorem}\label{thm:est3-semicoloreo}
%Para todo grafo \(G\),
%\[
%\estthree(G)=\chisemi(G).
%\]
%\end{theorem}

%\begin{proof}
%Se demuestran ambas desigualdades.

%\medskip
%\noindent\textbf{(i) \(\chisemi(G)\le \estthree(G)\).}
%Sea \(k=\estthree(G)\). sea \(S=(\prec,\mathcal Q)\) una solución  de valor \(k\) que evita \(\mathcal E_3\),
%donde \(\mathcal Q=\{V_1,\dots,V_k\}\).
%Para cada \(i\in\{1,\dots,k\}\), sea \(c_i\) el último vértice de \(V_i\) en el orden \(\prec\).


%Veamos que \(V_i\setminus\{c_i\}\) es independiente. 
%Supongamos que existen \(x,y\in V_i\setminus\{c_i\}\)
%con \(x\prec y\) y \(x\sim y\). Como \(y\neq c_i\), existe \(t\in V_i\) con \(y\prec t\). Entonces la terna \(x\prec y\prec t\) satisface que \(x\sim y\) y además \(x\) y \(t\) están en la misma clase
%\(V_i\), lo cual contradice que \(S\) evita \(\mathcal E_3\). El absurdo viene de asumir que existen \(x,y\in V_i\setminus\{c_i\}\) adyacentes.

%En consecuencia, la partición \(V=V_1\,\dot\cup\,\cdots\,\dot\cup\,V_k\) es un semi-coloreo con \(k\) colores,
%en cada clase \(V_i\), a lo sumo el vértice \(c_i\) puede tener vecinos dentro de \(V_i\).






%\medskip
%\noindent\textbf{(ii) \(\estthree(G)\le \chisemi(G)\).}
%Sea \(k=\chisemi(G)\). tomemos un semi-coloreo con \(k\) colores, es decir, una partición
%\[
%V \;=\; C_1\,\dot\cup\,\cdots\,\dot\cup\,C_k
%\]
%tal que, para cada \(i\), el conjunto
%\(
%\{v\in C_i:\ N(v)\cap C_i\neq\varnothing\}
%\)
%tiene a lo sumo un elemento. Para cada \(i\), si existe, denotemos por \(c_i\in C_i\) al vértice (único) que tiene algún
%vecino dentro de \(C_i\), si no existe tal vértice, se elige \(c_i\in C_i\) arbitrario.

%Se construye una solución \(S=(\prec,\mathcal Q)\) de valor \(k\) que evita \(\mathcal E_3\) tomando
%\[
%\mathcal Q:=\{C_1,\dots,C_k\},
%\]
%y fijando un orden \(\prec\) por bloques imponiendo
%\[
%C_1\prec C_2\prec\cdots\prec C_k,
%\]
%y dentro de cada bloque \(C_i\) ubicando al final al vértice \(c_i\) (y ordenando el resto arbitrariamente).

%Verifiquemos que \(S\) evita \(\mathcal E_3\).
%Sea  \(r\prec s\prec t\) una terna de vértices.
%Recordemos que \(\mathcal E_3=(A,N,C)\) con \(A=\{\{1,2\}\}\), \(N=\varnothing\) y \(C=\{1,3\}\).
%Como \(N=\varnothing\), la condición (2) nunca aplica. Por lo tanto, para que \((r,s,t)\) evite \(\mathcal E_3\)
%alcanza con que ocurra (1) o (3):
%\begin{enumerate}
%\item[(1)] falta alguna arista exigida por \(A\) (aquí: \(r\not\sim s\));
%\item[(3)] el par indicado por \(C\) no queda en la misma clase (aquí: \(r\) y \(t\) quedan en clases distintas).
%\end{enumerate}

%Si \(r\) y \(t\) pertenecen a clases distintas de \(\mathcal Q\), entonces se cumple (3) y \((r,s,t)\) evita \(\mathcal E_3\).
%Supongamos entonces que \(r\) y \(t\) pertenecen a una misma clase \(C_i\).

%Por construcción, dentro de \(C_i\) el vértice \(c_i\) se ubicó al final, y todo vértice de \(C_i\setminus\{c_i\}\)
%no tiene vecinos dentro de \(C_i\). Como \(r\prec s\prec t\) y \(r,t\in C_i\), se tiene \(r\in C_i\setminus\{c_i\}\),
%y por lo tanto \(r\) no es adyacente a ningún vértice de \(C_i\), en particular \(r\not\sim s\).
%Luego se cumple (1), y \((r,s,t)\) evita \(\mathcal E_3\).

%En consecuencia, toda terna \(r\prec s\prec t\) evita \(\mathcal E_3\), y \(S\) evita \(\mathcal E_3\).

%\end{proof}



\newcommand{\estonetwo}{\operatorname{est}_{12}}
\newcommand{\mathcalEonetwo}{\mathcal{E}_{12}}

\begin{definition}[\(\estonetwo\)]
Fijamos la familia de patrones
\[
\mathcal{E}_{12}:=\{\mathcal{E}_1,\mathcal{E}_2\}.
\]
Dado un grafo \(G\), se define
\[
\estonetwo(G)\;:=\;\minClases_{\mathcal{E}_{12}}(G),
\]
es decir, el mínimo número de clases de una solución \((\prec,\mathcal Q)\) de \(G\) que evita simultáneamente
\(\mathcal{E}_1\) y \(\mathcal{E}_2\).
\end{definition}






\begin{theorem}\label{thm:est12-vc}
Para todo grafo \(G\),
\[
\estonetwo(G)=\tauvc(G).
\]
donde \(\tau(G)\) denota el tamaño mínimo de un vertex cover de \(G\).
\end{theorem}



\begin{proof}
\medskip
\noindent\textbf{(i) \(\tauvc(G)\le \estonetwo(G)\).}
Como \(\estonetwo\) exige evitar simultáneamente \(\mathcal E_1\) y \(\mathcal E_2\), toda solución que evita
\(\{\mathcal E_1,\mathcal E_2\}\) evita en particular \(\mathcal E_2\).
Por el Lema \ref{lem:monotone-patterns} (monotonía por inclusión de patrones),
\[
\estonetwo(G)\ge \esttwo(G).
\]
Por el teorema \ref{thm:est2-vc}, \(\esttwo(G)=\tauvc(G)\), concluimos \(\tauvc(G)\le \estonetwo(G)\).


\medskip
\noindent\textbf{(ii) \(\estonetwo(G)\le \tauvc(G)\).}
Sea \(L\subseteq V(G)\) un vertex cover mínimo, y sea \(k=|L|=\tauvc(G)\).
Entonces \(I:=V(G)\setminus L\) es independiente.

Elegimos un vértice \(w\in L\) y escribimos \(L\setminus\{w\}=\{\ell_2,\dots,\ell_k\}\).
Definimos una partición \(\mathcal Q=\{V_1,\dots,V_k\}\) por
\[
V_1:= I\cup\{w\},
\qquad
V_i:=\{\ell_i\}\ \ \text{para } i=2,\dots,k,
\]
y fijamos un orden \(\prec\) poniendo primero todos los vértices de \(I\) (en cualquier orden),
luego \(\ell_2\prec\cdots\prec \ell_k\), y por último \(w\).




\begin{figure}[t]
\centering
\begin{tikzpicture}[
  font=\small,
  oval/.style={draw, rounded corners=14pt, minimum height=7mm, inner sep=2.5pt},
  dot/.style={circle, draw, minimum size=6mm, inner sep=0pt},
  >=latex
]
\def\dy{0.85cm}
\def\hg{0.75cm}
\def\xlab{-0.3} % columna fija para los labels V_i

% V1: I
\node at (\xlab,0) {$V_1$};
\node[oval, anchor=west, minimum width=30mm] (I) at (0,0) {$I$};

% V2: ell_2
\node at (\xlab,-\dy) {$V_2$};
\node[dot, anchor=west] (l2) at ($(I.east)+(\hg,-\dy)$) {$\ell_2$};

% V3: ell_3
\node at (\xlab,-2*\dy) {$V_3$};
\node[dot, anchor=west] (l3) at ($(l2.east)+(\hg,-\dy)$) {$\ell_3$};

% puntos
\node at ($(l3.east)+(0.8cm,-\dy)$) {$\vdots$};

% Vk: ell_k (alineado con la columna fija)
\node at (\xlab,-5*\dy) {$V_k$};
\node[dot, anchor=west] (lk) at ($(l3.east)+(1.1cm,-3*\dy)$) {$\ell_k$};

% w al final, misma fila que V1
\node[dot] (w) at ($(lk.east)+(1.6cm,5*\dy)$) {$w$};

% texto del orden abajo
\node[draw=none, anchor=north] at ($(I.west)+(2.2cm,-6.0*\dy)$)
{\(\prec:\) primero \(I\), luego \(\ell_2,\ell_3,\dots,\ell_k\), y al final \(w\).};

\end{tikzpicture}
\caption{Ilustración de la solución: \(V_1=I\cup\{w\}\) y \(V_i=\{\ell_i\}\) para \(i\ge 2\), con orden por bloques.}
\label{fig:est12-solucion}
\end{figure}



Verifiquemos que la solución \(S=(\prec,\mathcal Q)\) evita \(\mathcal E_1\) y \(\mathcal E_2\).

\smallskip
\noindent\textbf{Verifiquemos que se evita \(\mathcal E_1\).}
Sea \(r\prec s\prec t\). Como \(N=\varnothing\), alcanza con que ocurra (1) o (3).
Si \(r\) y \(s\) están en clases distintas, ocurre (3) y la terna evita \(\mathcal E_1\).
Supongamos que \(r\) y \(s\) están en la misma clase. Entonces necesariamente \(r,s\in V_1\).
Como \(s\prec t\), no puede ocurrir \(s=w\) (pues \(w\) es el último vértice del orden), luego \(s\in I\).
En particular \(r\in I\). Como \(I\) es independiente, se tiene \(r\not\sim s\), o sea ocurre (1).
Por lo tanto, toda terna evita \(\mathcal E_1\).

\smallskip
\noindent\textbf{Verifiquemos que se evita \(\mathcal E_2\).}
Sea \(r\prec s\prec t\). Nuevamente alcanza con (1) o (3).
Si \(s\) y \(t\) están en clases distintas, ocurre (3) y la terna evita \(\mathcal E_2\).
Supongamos que \(s\) y \(t\) están en la misma clase. Entonces necesariamente \(s,t\in V_1\).
Como \(t\) está a la derecha de \(s\), seguro que \(s\) no es \(w\),
luego \(s\in I\) más aun todos los vértices que están antes que \(s\) también pertenecen a \(I\),
en particular, \(r\in I\). Como \(I\) es independiente, se tiene \(r\not\sim s\), es decir ocurre (1).
Por lo tanto, toda terna evita \(\mathcal E_2\).

Con esto, \(S\) evita simultáneamente \(\mathcal E_1\) y \(\mathcal E_2\), y tiene valor \(|\mathcal Q|=k\).
Luego \(\estonetwo(G)\le k=\tauvc(G)\).
\end{proof}




\begin{definition}[\(\esttwothree\)]
Fijamos la familia de patrones
\[
\mathcal{E}_{23}:=\{\mathcal{E}_2,\mathcal{E}_3\}.
\]
Dado un grafo \(G\), se define
\[
\esttwothree(G)\;:=\;\minClases_{\mathcal{E}_{23}}(G),
\]
es decir, el mínimo número de clases de una solución \((\prec,\mathcal Q)\) de \(G\) que evita simultáneamente
\(\mathcal{E}_2\) y \(\mathcal{E}_3\).
\end{definition}


\begin{theorem}\label{thm:est23-vc}
Para todo grafo \(G\),
\[
\esttwothree(G)=\tau(G),
\]
donde \(\tau(G)\) denota el tamaño mínimo de un vertex cover de \(G\).
\end{theorem}

\begin{proof}


\medskip
\noindent\textbf{(i) \(\tau(G)\le \esttwothree(G)\).}
Por definición, \(\esttwothree(G)=\minClases_{\{\mathcal E_2,\mathcal E_3\}}(G)\) y
\(\esttwo(G)=\minClases_{\{\mathcal E_2\}}(G)\).
Como \(\{\mathcal E_2\}\subseteq \{\mathcal E_2,\mathcal E_3\}\), por el Lema~\ref{lem:monotone-patterns} se tiene
\[
\esttwothree(G)\ge \esttwo(G).
\]
Usando el Teorema~\ref{thm:est2-vc} (i.e.\ \(\esttwo(G)=\tau(G)\)), se concluye \(\tau(G)\le \esttwothree(G)\).

\medskip
\noindent\textbf{(ii) \(\esttwothree(G)\le \tau(G)\).}
Sea \(L\subseteq V(G)\) un vertex cover mínimo y sea \(k=|L|=\tau(G)\).
Entonces \(I:=V(G)\setminus L\) es independiente. Enumeremos
\[
L=\{\ell_1,\dots,\ell_k\}.
\]
Definamos una partición \(\mathcal Q=\{V_1,\dots,V_k\}\) por
\[
V_1:= I\cup\{\ell_1\},
\qquad
V_i:=\{\ell_i\}\ \ \text{para } i=2,\dots,k.
\]
Se fija un orden \(\prec\) por bloques imponiendo
\[
V_1 \prec V_2 \prec \cdots \prec V_k,
\]
y dentro de \(V_1\) se ubican primero todos los vértices de \(I\) (en cualquier orden) y al final \(\ell_1\).


Verifiquemos que la solución \(S=(\prec,\mathcal Q)\) evita simultáneamente \(\mathcal E_2\) y \(\mathcal E_3\).

\smallskip
\noindent\textbf{Verifiquemos que se evita \(\mathcal E_2\).}
Sea \(r\prec s\prec t\). Como \(N=\varnothing\), basta ver que ocurre (1) o (3) para \(\mathcal E_2\):
\[
\text{(1) } r\not\sim s
\qquad\text{o}\qquad
\text{(3) } s \text{ y } t \text{ están en clases distintas.}
\]
Si \(s\) y \(t\) están en clases distintas, ocurre (3). Supongamos entonces que \(s\) y \(t\) están en una misma clase.
Como \(V_2,\dots,V_k\) tienen solo un elemento, necesariamente \(s,t\in V_1\). Además, \(s\prec t\) implica \(s\neq \ell_1\),
luego \(s\in I\). Como \(V_1\) es el primer bloque y \(r\prec s\), también \(r\in V_1\) y como \(\ell_1\) está al final
de \(V_1\), se deduce \(r\in I\). Como \(I\) es independiente, \(r\not\sim s\), o sea ocurre (1).
Por lo tanto, toda terna evita \(\mathcal E_2\).

\smallskip
\noindent\textbf{Verifiquemos que se evita \(\mathcal E_3\).}
Sea \(r\prec s\prec t\). Nuevamente, como \(N=\varnothing\), basta ver que ocurre (1) o (3) para \(\mathcal E_3\):
\[
\text{(1) } r\not\sim s
\qquad\text{o}\qquad
\text{(3) } r \text{ y } t \text{ están en clases distintas.}
\]
Si \(r\) y \(t\) están en clases distintas, ocurre (3). Supongamos entonces que \(r\) y \(t\) están en una misma clase.
Otra vez, como \(V_2,\dots,V_k\) son singletons, necesariamente \(r,t\in V_1\), y por lo tanto también \(s\in V_1\).
Como \(r\prec s\), se tiene \(r\neq \ell_1\), luego \(r\in I\). Además, \(s\neq \ell_1\) (pues \(s\prec t\)),
luego \(s\in I\). Como \(I\) es independiente, \(r\not\sim s\), o sea ocurre (1).
Por lo tanto, toda terna evita \(\mathcal E_3\).

Con esto, \(S\) evita \(\{\mathcal E_2,\mathcal E_3\}\) y tiene valor \(|\mathcal Q|=k\). En consecuencia,
\(\esttwothree(G)\le k=\tau(G)\).
\end{proof}




\begin{definition}[\(\estonethree\)]
Se fija la familia de patrones
\[
\mathcal{E}_{13}:=\{\mathcal{E}_1,\mathcal{E}_3\}.
\]
Dado un grafo \(G\), se define
\[
\estonethree(G)\;:=\;\minClases_{\mathcal{E}_{13}}(G),
\]
es decir, el mínimo número de clases de una solución \((\prec,\mathcal Q)\) de \(G\) que evita simultáneamente
\(\mathcal{E}_1\) y \(\mathcal{E}_3\).
\end{definition}



\begin{lemma}\label{lem:est13-eq-est1}
Para todo grafo \(G\),
\[
\estonethree(G)=\estone(G).
\]
\end{lemma}

\begin{proof}
Como \(\{\mathcal E_1\}\subseteq\{\mathcal E_1,\mathcal E_3\}\), por el Lema~\ref{lem:monotonia-patrones} se tiene
\[
\estonethree(G)\ge \estone(G).
\]
Resta probar \(\estonethree(G)\le \estone(G)\).

Sea \(k=\estone(G)\). Fijamos una solución \(S=(\prec,\mathcal Q)\) de valor \(k\) que evita \(\mathcal E_1\), con
\(\mathcal Q=\{V_1,\dots,V_k\}\). Sea \(u\) el último vértice del orden \(\prec\).

\begin{Claim} Si \(x\prec y\), \(x,y\in V_i\) y \(x\sim y\), entonces \(y=u\).
\end{Claim}

\begin{proof}
Si \(y\neq u\), existe \(t\) con \(y\prec t\). Entonces la terna \((x,y,t)\) no evita \(\mathcal E_1\)
(porque \(x\sim y\) y \(x,y\) están en la misma clase), lo cual es una contradicción. Luego \(y=u\).
\end{proof}

Sea \(V_j\) la clase que contiene a \(u\) ESTO esta OK???. Se construye un nuevo orden \(\prec'\) por bloques respecto de \(\mathcal Q\)
poniendo primero todos los vértices de \(V_1,\dots,V_{j-1},V_{j+1},\dots,V_k\) (en cualquier orden por bloques)\eric{definir orden por bloques de una familia de conjuntos disjuntos dos a dos}, y al final
los vértices de \(V_j\). Dentro de \(V_j\), se ubican primero los vértices de \(V_j\setminus\{u\}\) (en cualquier orden) y
al final \(u\). Definimos \(S':=(\prec',\mathcal Q)\).

\medskip
\noindent\textbf{Verifiquemos que \(S'\) evita \(\mathcal E_1\).}
Sea \(r\prec' s\prec' t\).
Si \(r\) y \(s\) están en clases distintas, se cumple (3)\eric{citar la def de evitar} y la terna evita \(\mathcal E_1\).\eric{Sacar esta primera frase y asumir directamente lo siguiente}
Si \(r,s\in V_i\) y \(r\sim s\), por la observación \eric{citar} necesariamente \(s=u\). Pero \(u\) es el último vértice del orden
\(\prec'\), así que no existe \(t\) con \(s\prec' t\). Luego no existe ninguna terna que viole \(\mathcal E_1\), y \(S'\)
evita \(\mathcal E_1\).

\medskip
\noindent\textbf{Verifiquemos que \(S'\) evita \(\mathcal E_3\).}
Sea \(r\prec' s\prec' t\). Como \(N=\varnothing\), alcanza con que ocurra (1) o (3) para \(\mathcal E_3\),
o bien \(r\not\sim s\), o bien \(r\) y \(t\) están en clases distintas.

Si \(r\) y \(t\) están en clases distintas, se cumple (3) y no hay nada que ver.
Supongamos entonces que \(r,t\in V_i\). Como el orden \(\prec'\) es por bloques, también \(s\in V_i\).

Como \(s\prec' t\), el vértice \(s\) no puede ser el último del orden \(\prec'\). Pero el último vértice es \(u\),
luego \(s\neq u\). Por la observación anterior (toda arista intraclase tiene como extremo derecho a \(u\)),
se deduce que \(r\not\sim s\) (ESTO ES POR EL CONTRARECIPROCO DE LA AFIRMACION
LO EXPLICO MEJOR? O ASI SE ENTIENDE?). En consecuencia se cumple (1), y la terna \((r,s,t)\) evita \(\mathcal E_3\).


En consecuencia, \(S'\) evita simultáneamente \(\mathcal E_1\) y \(\mathcal E_3\), tiene valor \(k\), y por lo tanto
\[
\estonethree(G)\le k=\estone(G).
\]
\end{proof}





\begin{definition}[\(\estoneonetwothree\)]
Fijamos familia de patrones
\[
\mathcal{E}_{123}:=\{\mathcal{E}_1,\mathcal{E}_2,\mathcal{E}_3\}.
\]
Dado un grafo \(G\), se define
\[
\estoneonetwothree(G)\;:=\;\minClases_{\mathcal{E}_{123}}(G),
\]
es decir, el mínimo número de clases de una solución \((\prec,\mathcal Q)\) de \(G\) que evita simultáneamente
\(\mathcal{E}_1\), \(\mathcal{E}_2\) y \(\mathcal{E}_3\).
\end{definition}

\bigskip

\begin{theorem}\label{thm:star-all-vertex-cover}
For every graph \(G\),
\[
\operatorname{star}_{12\text{-}23\text{-}13}(G)
=
\begin{cases}
1, & \text{if } \tau(G)\le 1,\\[1mm]
\tau(G)+1, & \text{if } \tau(G)\ge 2,
\end{cases}
\]
where \(\tau(G)\) denotes the minimum size of a vertex cover of \(G\).
\end{theorem}

\begin{proof}
We distinguish two cases.

\medskip
\noindent\textbf{Case 1: \(\tau(G)\le 1\).}

If \(G\) has no edges, then any order together with the one-part partition
\[
\mathcal Q=\{V(G)\}
\]
avoids all three colored patterns, since every ordered triple satisfies the
\edgecondition.

Suppose now that \(\tau(G)=1\), and let \(z\in V(G)\) cover every edge of \(G\).
Choose an order in which \(z\) is the last vertex, and again take the one-part partition
\[
\mathcal Q=\{V(G)\}.
\]
Let
\[
v_1\prec v_2\prec v_3
\]
be an ordered triple. Since every edge of \(G\) is incident with \(z\), and \(z\) is the last vertex of the whole
order, the pair \(\{v_1,v_2\}\) is not an edge. Hence the triple satisfies the \edgecondition  for each of
the three star colored patterns.
\german{esto vale para todos los star}
Therefore,
\[
\operatorname{star}_{12\text{-}23\text{-}13}(G)\le 1.
\]
Since every solution has at least one part,
\[
\operatorname{star}_{12\text{-}23\text{-}13}(G)=1.
\]

\medskip
\noindent\textbf{Case 2: \(\tau(G)\ge 2\).}

We first prove the upper bound. Let \(L\) be a minimum vertex cover of \(G\), and let
\[
I:=V(G)\setminus L.
\]
Since \(L\) is a vertex cover, \(I\) is independent. Write
\[
L=\{z_1,\dots,z_{\tau(G)}\}.
\]
Consider the partition
\[
\mathcal Q
=
\bigl\{I,\{z_1\},\dots,\{z_{\tau(G)}\}\bigr\},
\]
which has \(\tau(G)+1\) parts. Define an order by placing first all vertices of \(I\), in arbitrary order, and then
all vertices of \(L\), also in arbitrary order.

Let
\[
v_1\prec v_2\prec v_3
\]
be an ordered triple. If \(v_1v_2\notin E(G)\), then the triple avoids all three patterns by the
\edgecondition.

Suppose that \(v_1v_2\in E(G)\). Since \(I\) is independent and all its vertices precede the vertices of \(L\),
the vertices \(v_1,v_2,v_3\) lie in pairwise distinct parts of \(\mathcal Q\). Indeed, either \(v_1\in I\) and
\(v_2,v_3\in L\), or all three vertices belong to \(L\).
Thus the \colorcondition holds for each of the colored pairs
\[
\{1,2\},\qquad \{1,3\},\qquad \{2,3\}.
\]
Hence the solution avoids all three patterns, and therefore
\[
\operatorname{star}_{12\text{-}23\text{-}13}(G)
\le \tau(G)+1.
\]

We now prove the lower bound. Let
\[
k:=\operatorname{star}_{12\text{-}23\text{-}13}(G),
\]
and let
\[
S=(\prec,\mathcal Q),
\qquad
\mathcal Q=\{V_1,\dots,V_k\},
\]
be a solution of value \(k\) that avoids the three star colored patterns.

For every \(i\in\{1,\dots,k\}\), let \(c_i\) be the last vertex of \(V_i\) in the order \(\prec\), and define
\[
L:=\{c_1,\dots,c_k\}.
\]

We first claim that \(L\) is a vertex cover of \(G\). Suppose otherwise that there are adjacent vertices
\[
x,y\in V(G)\setminus L,
\qquad
x\prec y.
\]
Since \(y\notin L\), it is not the last vertex of its part. Hence there exists a vertex \(z\) in the same part as
\(y\) such that
\[
x\prec y\prec z.
\]
The required edge \(xy\) is present, also \(y\) and \(z\) lie in the same part, so  neither the \colorcondition nor the \edgecondition
holds. Thus the triple does not avoid \(\operatorname{star}_{23}\), a contradiction.

Therefore \(L\) is a vertex cover, and
\[
\tau(G)\le k.
\]
Since \(\tau(G)\ge 2\), we have \(k\ge 2\).

Let \(c_h\) be the first vertex of \(L\) in the order \(\prec\). The idea of this demostration is
to show that we can remove \(c_h\) from \(L\) and still have a vertex cover. We claim that
\[
N_G(c_h)\cap \bigl(V(G)\setminus L\bigr)=\varnothing.
\]
Suppose, towards a contradiction, that there exists
\[
x\in V(G)\setminus L
\]
such that \(x\sim c_h\). Let \(V_i\) be the part containing \(x\), and let \(c_i\) be its last vertex.

If \(i=h\), then \(x\prec c_h\). Since \(k\ge 2\), there is another vertex \(c_j\in L\), and the choice of \(c_h\)
gives
\[
x\prec c_h\prec c_j.
\]
The edge \(xc_h\) is present and \(x,c_h\) lie in the same part. Hence neither the \textup{(edge)} condition nor
the \textup{(color)} condition holds for \(\operatorname{star}_{12}\), a contradiction.

Assume now that \(i\neq h\). Since \(c_h\) is the first vertex of \(L\),
\[
c_h\prec c_i.
\]
Also, \(x\prec c_i\), because \(c_i\) is the last vertex of \(V_i\).

If \(x\prec c_h\), then
\[
x\prec c_h\prec c_i.
\]
The edge \(xc_h\) is present and \(x,c_i\) lie in the same part, so the triple does not avoid
\(\operatorname{star}_{13}\).

If \(c_h\prec x\), then
\[
c_h\prec x\prec c_i.
\]
The edge \(c_hx\) is present and \(x,c_i\) lie in the same part, so the triple does not avoid
\(\operatorname{star}_{23}\).

Both alternatives give a contradiction. Therefore \(c_h\) has no neighbor outside \(L\).

Since \(L\) is a vertex cover and \(c_h\) has no neighbor in \(V(G)\setminus L\), the set
\[
L\setminus\{c_h\}
\]
is also a vertex cover of \(G\). Consequently,
\[
\tau(G)\le k-1,
\]
and hence
\[
k\ge \tau(G)+1.
\]

Combining both inequalities gives
\[
\operatorname{star}_{12\text{-}23\text{-}13}(G)=\tau(G)+1.
\]
\end{proof}




\subsection{Decision problems}

Let \(\mathcal P\) be a family of colored patterns. We consider the following three decision problems.

\eric{Poner los problemas en cajitas que queda mas lindo DONE}
\eric{Usar textsc en cada nombre de problema DONE}  
\begin{problem}[\(\mathcal P\)-Avoidance]
\textbf{Input:} A graph \(G\) and an integer \(k\).

\textbf{Question:} Does there exist a solution
\[
S=(\prec,\mathcal Q),
\qquad
\mathcal Q=\{V_1,\dots,V_k\},
\]
for \(G\) such that \(S\) avoids every colored pattern in \(\mathcal P\)?
\end{problem}

\eric{Cada vez que digamos "solution" hay que decir para que grafo es. OK}

\eric{Si las partes de la partición pueden ser vacías, podemos poner $\{V_1,\dots,V_k\}$ en cambio. OK}

\begin{problem}[\(\mathcal P\)-Avoidance with a given order]
\textbf{Input:} A graph \(G\), a total order \(\prec\) of \(V(G)\), and an integer \(k\).

\textbf{Question:} Does there exist a partition
\[
\mathcal Q=\{V_1,\dots,V_k\},
\]
of \(V(G)\)
such that the solution
\[
S=(\prec,\mathcal Q)
\]
for \(G\) avoids every colored pattern in \(\mathcal P\)?
\end{problem}



\begin{problem}[\(\mathcal P\)-Avoidance with a given partition]
\textbf{Input:} A graph \(G\) and a partition
\[
\mathcal Q=\{V_1,\dots,V_k\}
\]
of \(V(G)\).

\textbf{Question:} Does there exist a total order \(\prec\) of \(V(G)\) such that the solution
\[
S=(\prec,\mathcal Q)
\]
for \(G\) avoids every colored pattern in \(\mathcal P\)?
\end{problem}

When \(\mathcal P=\{P\}\) consists of a single colored pattern, we simply write \textsc{\(P\)-Avoidance} instead of
\textsc{\(\{P\}\)-Avoidance}.



In particular, we will study these three variants for the single patterns
\[
\forrs,\qquad \forrt,\qquad \forst,
\]
and for the pattern families
\[
\mathcal F_{rs\text{-}st},\qquad
\mathcal F_{rs\text{-}rt},\qquad
\mathcal F_{rt\text{-}st},\qquad
\mathcal F_{rs\text{-}st\text{-}rt}.
\]



\german{puse las dos formas con el algoritmo y sin, eligamos cual tmb podria poner el algoritmo dentro de la segunda demo}
\german{por alguna extrana razon los algoritmos q escribi estan ambos al final}


\begin{algorithm}[h]
\caption{\textsc{Two-Parts-Given-Order}}
\label{alg:two-parts-given-order}
\begin{algorithmic}[1]
\Require A graph \(G\), a total order \(\prec\) of \(V(G)\), a pattern
\(P=(A,N,\varnothing)\), and a set \(D\subseteq\{12,13,23\}\).
\Ensure \textbf{YES} if \((G,\prec)\) admits a valid solution with at most two parts; otherwise \textbf{NO}.

\State Construct a graph \(R\) with
\[
V(R)\gets V(G),
\qquad
E(R)\gets\varnothing.
\]

\ForAll{ordered triples \(v_1\prec v_2\prec v_3\) of \(G\)}
    \If{\((v_1,v_2,v_3)\) realizes \(P\)}
        \If{\(12\in D\)}
            \State \(E(R)\gets E(R)\cup\{\{v_1,v_2\}\}\)
        \EndIf
        \If{\(13\in D\)}
            \State \(E(R)\gets E(R)\cup\{\{v_1,v_3\}\}\)
        \EndIf
        \If{\(23\in D\)}
            \State \(E(R)\gets E(R)\cup\{\{v_2,v_3\}\}\)
        \EndIf
    \EndIf
\EndFor

\If{\(R\) is bipartite}
    \State \Return \textbf{YES}
\Else
    \State \Return \textbf{NO}
\EndIf
\end{algorithmic}
\end{algorithm}

The graph \(R\) is called the \emph{constraint graph}. Each edge \(v_iv_j\in E(R)\) records a constraint requiring
\(v_i\) and \(v_j\) to lie in different parts of the solution.






\begin{proof}
Let
\[
R=R_{P,D}(G,\prec)
\]
be the constraint graph constructed by Algorithm~\ref{alg:two-parts-given-order}.

Suppose first that \(R\) is bipartite, and let
\[
\mathcal Q=\{V_1,V_2\}
\]
be the partition induced by a proper \(2\)-coloring of \(R\). We claim that
\[
S=(\prec,\mathcal Q)
\]
avoids every colored pattern in \(\mathcal P_D\).

Let \(v_1\prec v_2\prec v_3\) be an ordered triple. If the triple avoids the base pattern \(P\), then it also
avoids every colored version of \(P\) by \edgecondition{} or \nonedgecondition.

Suppose instead that the triple realizes \(P\). For every \(ij\in D\), the algorithm adds the edge
\[
v_iv_j
\]
to \(R\). Since the coloring of \(R\) is proper, \(v_i\) and \(v_j\) lie in different parts of \(\mathcal Q\).
Hence the triple avoids \(P_{ij}\) by \colorcondition. Therefore \(S\) avoids \(\mathcal P_D\).

Conversely, suppose that there exists a partition
\[
\mathcal Q=\{V_1,V_2\}
\]
such that \(S=(\prec,\mathcal Q)\) avoids \(\mathcal P_D\). We show that \(\mathcal Q\) defines a proper
\(2\)-coloring of \(R\).

Let \(v_iv_j\in E(R)\). By construction, there exists an ordered triple
\[
v_1\prec v_2\prec v_3
\]
that realizes \(P\), with \(ij\in D\). Since the triple realizes the base pattern, neither
\edgecondition{} nor \nonedgecondition{} holds. Therefore, in order to avoid \(P_{ij}\), the triple must satisfy
\colorcondition. Hence \(v_i\) and \(v_j\) lie in different parts of \(\mathcal Q\).

Thus every edge of \(R\) has endpoints in different parts, so \(R\) is bipartite.

Finally, the algorithm examines \(O(|V(G)|^3)\) ordered triples, adds at most three constraints per triple, and checks
bipartiteness in linear time in the size of \(R\). Therefore its total running time is
\[
O(|V(G)|^3).
\]
\end{proof}



\paragraph{Intuition.}
When \(k=2\), every constraint requiring two vertices to lie in different parts has the form
\[
c(v_i)\neq c(v_j),
\qquad
c:V(G)\to\{1,2\}.
\]
Such a constraint can be represented by an edge \(v_iv_j\) in an auxiliary graph. Therefore, all the color
constraints imposed by the ordered triples that realize \(P\) can be collected in a single constraint graph
\(R_{P,D}(G,\prec)\). A valid partition into at most two parts then exists exactly when this constraint graph is
bipartite.

\begin{theorem}\label{thm:given-order-two-parts}
Let
\[
P=(A,N,\varnothing)
\]
be a pattern on three vertices, and let
\[
D\subseteq\{12,13,23\}.
\]
For every \(ij\in D\), define the colored pattern
\[
P_{ij}:=(A,N,\{i,j\}),
\]
and let
\[
\mathcal P_D:=\{P_{ij}:ij\in D\}.
\]
Thus, \(\mathcal P_D\) is the family of colored patterns obtained from \(P\) by using the pairs indexed by \(D\)
as color constraints.

The \(\mathcal P_D\)-Avoidance problem with a given order, restricted to instances with
\[
k=2,
\]
can be solved in polynomial time. More precisely, given a graph \(G\) and a total order \(\prec\) of \(V(G)\),
one can construct a graph
\[
R_{P,D}(G,\prec)
\]
such that
\[
(G,\prec)\text{ admits a valid solution with at most two parts}
\quad\Longleftrightarrow\quad
R_{P,D}(G,\prec)\text{ is bipartite}.
\]
\end{theorem}

\begin{proof}
The constraint graph
\[
R=R_{P,D}(G,\prec)
\]
is defined by
\[
V(R):=V(G),
\]
and
\[
E(R):=
\bigcup_{\substack{v_1\prec v_2\prec v_3\\
(v_1,v_2,v_3)\text{ realizes P in \(G\)}}} 
\bigl\{
\{v_i,v_j\}:ij\in D
\bigr\}.
\]

We prove the equivalence.

Suppose first that there exists a partition
\[
\mathcal Q=\{V_1,V_2\}
\]
such that the solution
\[
S=(\prec,\mathcal Q)
\]
avoids every colored pattern in \(\mathcal P_D\) in \(G\).
Color each vertex of \(R\) according to its part in \(\mathcal Q\).

Let \(v_iv_j\in E(R)\). By the definition of \(R\), there exists an ordered triple
\[
v_1\prec v_2\prec v_3
\]
that realizes \(P\) in \(G\), with \(ij\in D\). Since \(P\) is not avoided, but \(\mathcal P_D\) is,
the triple avoid \(P_{ij}\) only by \colorcondition. Therefore \(v_i\) and \(v_j\) lie in different parts of
\(\mathcal Q\). Hence the induced coloring of \(R\) is proper, and \(R\) is bipartite.


Conversely, suppose that \(R\) is bipartite, and let
\[
c:V(R)\to\{1,2\}
\]
be a proper \(2\)-coloring of \(R\). Since \(V(R)=V(G)\), we take the color classes of \(c\) as a partition
\(\mathcal Q\) of \(V(G)\) into at most two parts.

We claim that the solution
\[
S=(\prec,\mathcal Q)
\]
avoids every colored pattern in \(\mathcal P_D\) in \(G\).
Consider an ordered triple
\[
v_1\prec v_2\prec v_3.
\]
If the triple avoids the base pattern \(P\), then it also avoids every colored pattern \(P_{ij}\) by
\edgecondition{} or \nonedgecondition.

Suppose instead that the triple realizes \(P\). For every \(ij\in D\), the edge \(v_iv_j\) belongs to \(E(R)\).
Since \(c\) is a proper coloring of \(R\),
\[
c(v_i)\neq c(v_j).
\]
Thus \(v_i\) and \(v_j\) lie in different parts of \(\mathcal Q\), and the triple avoids \(P_{ij}\) by
\colorcondition. Therefore \(S\) avoids every pattern in \(\mathcal P_D\).

To construct \(R\), we enumerate the \(O(|V(G)|^3)\) ordered triples of \(G\), test whether each triple realizes
\(P\), and add at most \(|D|\le 3\) edges. We then test whether \(R\) is bipartite using BFS or DFS.
Using an adjacency matrix for \(G\), the total running time is
\[
O(|V(G)|^3).
\]
\end{proof}









\begin{theorem}\label{thm:forrs-given-order-np-complete}
The \(\forrs\)-Avoidance problem with a given order is NP-complete.
More precisely, for every fixed \(k\ge 3\), deciding whether there exists a partition into at most \(k\) parts
that avoids \(\forrs\) with respect to a given order is NP-complete.
\end{theorem}




\begin{proof}[Proof sketch]
Reduce from \(k\)-\textsc{Coloring}. Let \(X\) be an instance with
\(V(X)=\{x_1,\dots,x_n\}\). For each edge \(x_i x_j\in E(X)\), with \(i<j\),
add a witness vertex \(w_{ij}\), adjacent only to \(x_i\) and \(x_j\), and do
not keep the original edge \(x_i x_j\). Thus,
\[
V(G)=V(X)\cup \{w_{ij}: x_i x_j\in E(X),\ i<j\},
\]
\eric{También fijarse si usamos siempre : o mid para conjuntos por comprensión}
and
\[
E(G)=\{x_iw_{ij},x_jw_{ij}: x_i x_j\in E(X),\ i<j\}.
\]
The given order is
\[
x_1\prec \cdots \prec x_n \prec \text{all vertices } w_{ij}.
\]

If \(X\) has a proper \(k\)-coloring, keep the colors on the original vertices
and assign the witnesses arbitrarily. Then each witness \(w_{ij}\) has earlier
neighbors \(x_i,x_j\), which have different colors.

Conversely, if a \(\forrs\)-avoiding partition exists, then for every edge
\(x_i x_j\in E(X)\), with \(i<j\), the triple
\[
x_i\prec x_j\prec w_{ij}
\]
forces \(x_i\) and \(x_j\) to lie in different parts. Hence the restriction to
\(V(X)\) is a proper \(k\)-coloring of \(X\).
\end{proof}






\begin{theorem}\label{thm:forrs-general-np-complete}
The \(\forrs\)-Avoidance problem is NP-complete.
More precisely, for every fixed \(k\ge 3\), deciding whether a graph admits a
\(\forrs\)-avoiding solution with at most \(k\) parts is NP-complete.
\end{theorem}





\begin{proof}[Proof sketch]
Reduce from \(k\)-\textsc{Coloring}. Given an instance \(X\), for each edge
\(xy\in E(X)\), add \(k+1\) common witnesses
\[
W_{xy}:=\{w^1_{xy},\dots,w^{k+1}_{xy}\},
\]
each adjacent only to \(x\) and \(y\), and do not keep the original edge \(xy\).
Thus,
\[
V(G)=V(X)\cup \bigcup_{xy\in E(X)} W_{xy},
\]
and
\[
E(G)=\{xw^i_{xy},yw^i_{xy}: xy\in E(X),\ 1\le i\le k+1\}.
\]

If \(X\) has a proper \(k\)-coloring, color the original vertices accordingly,
assign the witnesses arbitrarily, and order first the original vertices and then
the witnesses.

Conversely, suppose that a \(\forrs\)-avoiding solution with at most \(k\) parts
exists. If two adjacent original vertices \(x,y\) were in the same part, say
\(x\prec y\), then all \(k+1\) witnesses for \(xy\) must appear before \(y\).
Two of them lie in the same part, and together with \(y\) they form an occurrence
of \(\forrs\), a contradiction. Hence the restriction to \(V(X)\) is a proper
\(k\)-coloring of \(X\).
\end{proof}



\begin{definition}[\(\mathcal P\)-removable vertex]
Let \(G\) be a graph, let
\[
\mathcal Q=\{V_1,\dots,V_k\}
\]
be a fixed partition of \(V(G)\), and let \(H\) be an induced subgraph of \(G\). Let \(\mathcal P\) be a family of colored patterns.

A vertex \(v\in V(H)\) is called \emph{\(\mathcal P\)-removable} in \(H\) with respect to \(\mathcal Q\) if no pattern
in \(\mathcal P\) can be realized with \(v\) playing the role of the last vertex of the ordered triple.

More explicitly, \(v\) is \(\mathcal P\)-removable if, for every pattern \(P\in\mathcal P\), and for every two distinct
vertices \(x,y\in V(H)\setminus\{v\}\), the ordered triple
\[
(x,y,v)
\]
avoids \(P\) in \(H\). \german{importante agregar in H no? tmb lo cambie en la def de arriba de todo}\eric{Sip}.
\german{hace falta poner las dos triplas? con poner (x,y,z) alcanza no? pq intercambian el papel}\eric{No hace falta, tenes razon}
\end{definition}



Equivalently, placing \(v\) at the end of the current induced subgraph \(H\) does not create any forbidden occurrence
whose last vertex is \(v\). \german{esta ok esto?}




\begin{lemma}\label{lem:removable-for12}
Let \(G\) be a graph, let
\[
\mathcal Q=\{V_1,\dots,V_k\}
\]
be a fixed partition of \(V(G)\), and let \(H\) be an induced subgraph of \(G\).
A vertex \(v\in V(H)\) is \(\mathcal F_{12}\)-removable in \(H\) with respect to \(\mathcal Q\) if and only if
\[
|N_H(v)\cap V_i|\le 1
\qquad\text{for every } i=1,\dots,k.
\]

In other words, \(v\) is removable if it has no two neighbors in the same part of the partition.

\end{lemma}

\begin{proof}

Suppose that
\[
|N_H(v)\cap V_i|\le 1
\qquad\text{for every } i=1,\dots,k.
\]
Let \(x,y\in V(H)\setminus\{v\}\) be two distinct vertices adjacents to \(v\), then by the assumption they lie in different
parts of \(\mathcal Q\). Hence the color condition for \(\operatorname{for}_{12}\) fails, and the triples avoid the
pattern. Therefore \(v\) is \(\mathcal F_{12}\)-removable.

Conversely, suppose that there exists a part \(V_i\) such that
\[
|N_H(v)\cap V_i|\ge 2.
\]
Choose two distinct vertices \(x,y\in N_H(v)\cap V_i\). Then \(x\sim v\), \(y\sim v\), and \(x,y\) lie in the same
part. Thus the ordered triple \((x,y,v)\) realizes \(\operatorname{for}_{12}\). Therefore \(v\) is not
\(\mathcal F_{12}\)-removable.

\end{proof}


\begin{lemma}\label{lem:removable-for13-23}
Let \(G\) be a graph, let
\[
\mathcal Q=\{V_1,\dots,V_k\}
\]
be a fixed partition of \(V(G)\), and let \(H\) be an induced subgraph of \(G\).
For a vertex \(v\in V(H)\), let \(\mathcal Q(v)\) denote the part of \(\mathcal Q\) that contains \(v\).
Then \(v\) is \(\mathcal F_{13\text{-}23}\)-removable in \(H\) with respect to \(\mathcal Q\) if and only if
\[
\deg_H(v)\le 1
\qquad\text{or}\qquad
N_H(v)\cap \mathcal Q(v)=\varnothing.
\]
\end{lemma}

That is, \(v\) is removable if it has at most one neighbor in the current graph, or if none of its current neighbors
lies in its own part. \german{esto va afuera o adentro del lemma?}


\begin{proof}
Suppose that \(\deg_H(v)\le 1\). Then no pair of distinct vertices \(x,y\in V(H)\setminus\{v\}\) can both be
adjacent to \(v\). Hence every ordered triple ending at \(v\) misses one of the required edges, and therefore avoids
both \(\operatorname{for}_{13}\) and \(\operatorname{for}_{23}\).

Now suppose instead that
\[
N_H(v)\cap \mathcal Q(v)=\varnothing.
\]
Let \(x,y\in V(H)\setminus\{v\}\) be two distinct vertices adjacents to \(v\). Then neither of them lies in the same part as \(v\).
Thus the color condition fails for both \(\operatorname{for}_{13}\) and \(\operatorname{for}_{23}\). Hence \(v\) is
\(\mathcal F_{13\text{-}23}\)-removable.

Conversely, suppose that
\[
\deg_H(v)\ge 2
\]
and that \(v\) has a neighbor \(x\in N_H(v)\cap \mathcal Q(v)\). Since \(\deg_H(v)\ge 2\), choose another neighbor
\(y\in N_H(v)\), with \(y\neq x\). Then the ordered triple \((x,y,v)\) realizes \(\operatorname{for}_{13}\): both
\(x\) and \(y\) are adjacent to \(v\), and \(x\) lies in the same part as \(v\). Therefore \(v\) is not
\(\mathcal F_{13\text{-}23}\)-removable.


\end{proof}







\begin{lemma}\label{lem:removable-for12-23-13}
Let \(G\) be a graph, let
\[
\mathcal Q=\{V_1,\dots,V_k\}
\]
be a fixed partition of \(V(G)\), and let \(H\) be an induced subgraph of \(G\).
For a vertex \(v\in V(H)\), let \(\mathcal Q(v)\) denote the part of \(\mathcal Q\) that contains \(v\).
Then \(v\) is \(\mathcal F_{12\text{-}23\text{-}13}\)-removable in \(H\) with respect to \(\mathcal Q\) if and only if
either
\[
\deg_H(v)\le 1,
\]
or both of the following conditions hold:
\[
N_H(v)\cap \mathcal Q(v)=\varnothing
\]
and
\[
|N_H(v)\cap V_i|\le 1
\qquad\text{for every } i=1,\dots,k.
\]

Equivalently, if \(v\) has at least two neighbors, then all its neighbors must lie in parts different from
\(\mathcal Q(v)\), and no two of them may lie in the same part.
\end{lemma}



\begin{proof}
Suppose that \(\deg_H(v)\le 1\). Then no ordered triple ending at \(v\) can contain two required edges incident
with \(v\). Hence \(v\) is removable.

Now suppose that
\[
N_H(v)\cap \mathcal Q(v)=\varnothing
\]
and
\[
|N_H(v)\cap V_i|\le 1
\qquad\text{for every } i=1,\dots,k.
\]

Let \(x,y\in V(H)\setminus\{v\}\) be two distinct vertices adjacents to \(v\), then \(x\) and \(y\)
lie in different parts, and neither lies in the same part as \(v\). Therefore the color condition fails for each of
\(\operatorname{for}_{12}\), \(\operatorname{for}_{13}\), and \(\operatorname{for}_{23}\). Hence \(v\) is
\(\mathcal F_{12\text{-}23\text{-}13}\)-removable.

Conversely, suppose that \(v\) is \(\mathcal F_{12\text{-}23\text{-}13}\)-removable and that \(\deg_H(v)\ge 2\).
We prove the two conditions.

First, \(v\) cannot have a neighbor in its own part. Indeed, if
\[
x\in N_H(v)\cap \mathcal Q(v),
\]
then, since \(\deg_H(v)\ge 2\), we can choose another neighbor \(y\in N_H(v)\), with \(y\neq x\). The ordered triple
\((x,y,v)\) realizes \(\operatorname{for}_{13}\), contradicting removability. Thus
\[
N_H(v)\cap \mathcal Q(v)=\varnothing.
\]

Second, no part can contain two neighbors of \(v\). If for some \(i\) there exist distinct vertices
\[
x,y\in N_H(v)\cap V_i,
\]
then the ordered triple \((x,y,v)\) realizes \(\operatorname{for}_{12}\), since \(x\sim v\), \(y\sim v\), and \(x,y\)
lie in the same part. This contradicts removability. Therefore
\[
|N_H(v)\cap V_i|\le 1
\qquad\text{for every } i=1,\dots,k.
\]

\end{proof}





\eric{Pensar si poner un guion despues de two o no}
\begin{definition}[First-two-symmetric family]
Let \(\mathcal P\) be a family of colored patterns. We say that \(\mathcal P\) is
\emph{first-two-symmetric} if, for every graph \(H\), every partition \(\mathcal Q\) of \(V(H)\),
and every three distinct vertices \(x,y,z\in V(H)\), the ordered triple
\[
(x,y,z)
\]
avoids \(\mathcal P\) in \(H\) if and only if the ordered triple
\[
(y,x,z)
\]
avoids \(\mathcal P\) in \(H\).

In other words, once the last vertex of the ordered triple is fixed, the order of the first two vertices does not matter.
\end{definition}




\begin{algorithm}[ht]
\caption{\textsc{Given-Partition-\(\mathcal P\)}}
\label{alg:generic-given-partition}
\begin{algorithmic}[1]
\Require A graph \(G=(V(G),E(G))\), a partition \(\mathcal Q=\{V_1,\dots,V_k\}\) of \(V(G)\), and a family of colored patterns \(\mathcal P\).
\Ensure \textbf{YES} and an order \(\prec\) if such an order exists; otherwise \textbf{NO}.

\State \(A \gets V(G)\)
\State \(L \gets [\,]\) \Comment{\(L\) stores the vertices in deletion order}

\While{\(A\neq\varnothing\)}
    \If{there exists a \(\mathcal P\)-removable vertex \(v\in A\) in \(G[A]\) with respect to \(\mathcal Q\)}
        \State append \(v\) to \(L\)
        \State \(A\gets A\setminus\{v\}\)
    \Else
        \State \Return \textbf{NO}
    \EndIf
\EndWhile

\State Let \(\prec\) be the order obtained by reading \(L\) from right to left
\State \Return \((\textbf{YES},\prec)\)
\end{algorithmic}
\end{algorithm}


\begin{lemma}[Correctness of Algorithm~\ref{alg:generic-given-partition}]
\label{lem:generic-given-partition-correctness}
Let \(\mathcal P\) be a first-two symmetric family of colored patterns and let \(G\) be a graph.
Then Algorithm~\ref{alg:generic-given-partition} returns \textbf{YES} if and only if there exists an order
\(\prec\) of \(V(G)\) such that
\[
(\prec,\mathcal Q)
\]
avoids every pattern in \(\mathcal P\) in \(G\).
\end{lemma}


\begin{proof}
First suppose that the algorithm returns \textbf{YES}. Let \(\prec\) be the order obtained by reading the deletion
list \(L\) from right to left. We show that \((\prec,\mathcal Q)\) avoids \(\mathcal P\) in \(G\).

Consider any ordered triple
\[
x\prec y\prec z
\]
in the final order. When \(z\) was deleted by the algorithm, both \(x\) and \(y\) were still active, because they
appear before \(z\) in the final order. Since \(z\) was \(\mathcal P\)-removable in the current induced subgraph,
the triple \((x,y,z)\) avoids every pattern in \(\mathcal P\). Therefore every ordered triple avoids \(\mathcal P\),
and so \((\prec,\mathcal Q)\) is a valid solution.

Conversely, suppose that there exists an order \(\prec\) of \(V(G)\) such that \((\prec,\mathcal Q)\) avoids
\(\mathcal P\). We prove that the algorithm returns \textbf{YES}.

At any step, let \(H=G[A]\) be the current induced subgraph. The restriction of \(\prec\) to \(A\) still avoids
\(\mathcal P\), because deleting vertices cannot create a forbidden ordered triple. Let \(v\) be the last vertex of
this restricted order.

We claim that \(v\) is \(\mathcal P\)-removable in \(H\). Take any two distinct vertices
\(x,y\in V(H)\setminus\{v\}\). In the restricted order, either
\[
x\prec y\prec v
\]
or
\[
y\prec x\prec v.
\]
The corresponding ordered triple avoids \(\mathcal P\), because the restricted order is valid. Since \(\mathcal P\)
is first-two symmetric, the other ordering of the first two vertices also avoids \(\mathcal P\). Hence \((x,y,v)\)
avoids \(\mathcal P\) for every pair \(x,y\), so \(v\) is \(\mathcal P\)-removable.

Thus, whenever a valid order exists for the current induced subgraph, a \(\mathcal P\)-removable vertex exists.
Therefore the algorithm cannot return \textbf{NO}.
\end{proof}

\begin{remark}\label{lem:swap-first-two}
Let \(\mathcal P\) be a family of colored patterns. Suppose that whenever \(P\in\mathcal P\), the pattern obtained from \(P\) by swapping positions \(1\) and \(2\) also belongs to \(\mathcal P\). Then \(\mathcal P\) is first-two symmetric.
\end{remark}


\begin{lemma}\label{lem:families-first-two-symmetric}
The families
\[
\mathcal F_{12},
\qquad
\mathcal F_{13\text{-}23},
\qquad
\mathcal F_{12\text{-}23\text{-}13}
\]
are first-two symmetric.\eric{$\mathcal{F}_{12}$ no existe, o definirla o usar $\{for_{12}\}$.}
\end{lemma}

\begin{proof}
For \(\mathcal F_{12}\), swapping the first two vertices preserves the pattern \(\operatorname{for}_{12}\):
the required edges are still the two edges incident with the third vertex, and the colored pair \(\{1,2\}\)
is unchanged.

For \(\mathcal F_{13\text{-}23}\), swapping the first two vertices exchanges the two patterns
\(\operatorname{for}_{13}\) and \(\operatorname{for}_{23}\). Since both patterns belong to the family, avoidance of the
family is unchanged.

Finally, \(\mathcal F_{12\text{-}23\text{-}13}\) contains all three patterns, and is therefore also invariant under
swapping the first two vertices. \german{dudo de si esta demo tiene q ser mas formal}\eric{como quieras, esta bien ahora}
\end{proof}





\begin{corollary}\label{cor:given-partition-poly-cases}
Let \(G=(V(G),E(G))\) be a graph, let
\[
\mathcal Q=\{V_1,\dots,V_k\}
\]
be a fixed partition of \(V(G)\), and set
\[
n:=|V(G)|,
\qquad
m:=|E(G)|.
\]
The given-partition avoidance problem is polynomial-time solvable for the following families:
\[
\mathcal F_{12},
\qquad
\mathcal F_{13\text{-}23},
\qquad
\mathcal F_{12\text{-}23\text{-}13}.
\]
\german{esto verlo bien con eric pq ni idea si esta bien}
More precisely, Algorithm~\ref{alg:generic-given-partition} can be implemented in time
\[
O(nk+m)
\]
for \(\mathcal F_{12}\), in time
\[
O(n+m)
\]
for \(\mathcal F_{13\text{-}23}\), and in time
\[
O(nk+m) \german{este segun eric es O(n + mlog(n)) asi que dudo de todos jaja}
\]
for \(\mathcal F_{12\text{-}23\text{-}13}\).
\end{corollary}

\begin{proof}
By Lemma~\ref{lem:families-first-two-symmetric}, the families
\[
\mathcal F_{12},
\qquad
\mathcal F_{13\text{-}23},
\qquad
\mathcal F_{12\text{-}23\text{-}13}
\]
are first-two symmetric. Therefore, by Lemma~\ref{lem:generic-given-partition-correctness},
Algorithm~\ref{alg:generic-given-partition} is correct for each of these families.

It remains to justify the running times.

\medskip
\noindent\textbf{The family \(\mathcal F_{12}\).}
By Lemma~\ref{lem:removable-for12}, a vertex \(v\) is \(\mathcal F_{12}\)-removable in the current induced subgraph
\(H\) if and only if
\[
|N_H(v)\cap V_i|\le 1
\qquad\text{for every } i=1,\dots,k.
\]
We maintain counters
\[
c(v,i):=|N_H(v)\cap V_i|
\]
for every active vertex \(v\) and every part \(V_i\), together with the number of parts \(V_i\) for which
\(c(v,i)\ge 2\). Then \(v\) is removable exactly when this number is zero.

The counters can be initialized in time \(O(nk+m)\). When a vertex \(u\in V_j\) is deleted, only the counters
\(c(v,j)\) for active neighbors \(v\) of \(u\) change. Thus each edge is processed at most once during the deletion
process. Using a queue of currently removable vertices, the total running time is
\[
O(nk+m).
\]

\medskip
\noindent\textbf{The family \(\mathcal F_{13\text{-}23}\).}
By Lemma~\ref{lem:removable-for13-23}, a vertex \(v\) is \(\mathcal F_{13\text{-}23}\)-removable in the current
induced subgraph \(H\) if and only if
\[
\deg_H(v)\le 1
\qquad\text{or}\qquad
N_H(v)\cap \mathcal Q(v)=\varnothing,
\]
where \(\mathcal Q(v)\) denotes the part of \(\mathcal Q\) containing \(v\).

We maintain the current degree \(\deg_H(v)\) of each active vertex \(v\), and also the number
\[
s(v):=|N_H(v)\cap \mathcal Q(v)|
\]
of active neighbors of \(v\) in its own part. When a vertex is deleted, these quantities are updated only for its
active neighbors. Hence each edge is processed at most once, and the total running time is
\[
O(n+m).
\]

\medskip
\noindent\textbf{The family \(\mathcal F_{12\text{-}23\text{-}13}\).}
By Lemma~\ref{lem:removable-for12-23-13}, a vertex \(v\) is \(\mathcal F_{12\text{-}23\text{-}13}\)-removable in
the current induced subgraph \(H\) if and only if either
\[
\deg_H(v)\le 1,
\]
or both
\[
N_H(v)\cap \mathcal Q(v)=\varnothing
\]
and
\[
|N_H(v)\cap V_i|\le 1
\qquad\text{for every } i=1,\dots,k.
\]
Thus we combine the counters used in the previous two cases: the degree \(\deg_H(v)\), the same-part counter
\(s(v)=|N_H(v)\cap \mathcal Q(v)|\), and the counters
\[
c(v,i):=|N_H(v)\cap V_i|.
\]
Again, these can be initialized in time \(O(nk+m)\), and each edge is processed at most once during the deletion
process. Hence the total running time is
\[
O(nk+m).
\]

Therefore the given-partition avoidance problem is polynomial-time solvable for all three families, with the stated
running times.
\end{proof}



\eric{Me gustó la idea de sacar factor común, podemos quizá sacar la property para afuera y simplemente que la property sea "si no importa el orden de los dos primeros vertices de la tripla, entonces el algoritmo 1 es correcto." asdf}

\eric{1. Decir la property: "no me importa el orden de los dos primeros vertices de la tripla."
2. presentas el algoritmo 1 diciendo que lo vas a usar para ....
3. Decir el lema que si tenes la property, entonces el algoritmo 1 es correcto.
4. Decir que las tres familias de patrones tienen la property.
5. Corolario: las tres familias de patrones tienen algoritmo poli. Pones la complejidad de instanciar el algoritmo 1 para cada uno en particular.}
\german{OK}






\begin{theorem}\label{thm:forst-given-order-np-complete}
The \(\forst\)-Avoidance problem with a given order is NP-complete.
More precisely, for every fixed \(k\ge 3\), deciding whether there exists a partition into at most \(k\) parts
that avoids \(\forst\) with respect to a given order is NP-complete.
\end{theorem}

\begin{proof}[Proof sketch]
Reduce from \(k\)-\textsc{Coloring}. Let \(X\) be an instance with an arbitrary order on its vertices.
For each edge \(xy\in E(X)\), with \(x\prec_X y\), add a witness vertex \(w_{xy}\), adjacent only to \(y\), and keep
the original edge \(xy\). Thus,
\[
V(G)=V(X)\cup \{w_{xy}: xy\in E(X),\ x\prec_X y\},
\]
and
\[
E(G)=E(X)\cup \{w_{xy}y: xy\in E(X),\ x\prec_X y\}.
\]


\begin{figure}[h]
\centering
\begin{tikzpicture}[
  scale=1,
  every node/.style={circle, draw, inner sep=1.4pt, font=\small}
]

\node (w) at (0,0) {$w_{xy}$};
\node (x) at (2.2,0) {$x$};
\node (y) at (4.4,0) {$y$};

\draw (x)--(y);
\draw (w) to[bend left=35] (y);

\node[draw=none, font=\small] at (2.2,-1.1) {$w_{xy}\prec x\prec y$};

\end{tikzpicture}
\caption{Gadget used in the reduction for \(\forst\) with a given order.}
\label{fig:forst-given-order-gadget}
\end{figure}


The given order \(\prec\) puts first all witnesses \(w_{xy}\), in arbitrary order, and then all original vertices,
following the order \(\prec_X\).

If \(X\) has a proper \(k\)-coloring, keep the colors on the original vertices and assign each witness \(w_{xy}\)
a color different from the color of \(y\). Then no occurrence of \(\forst\) appears.

Conversely, suppose that a \(\forst\)-avoiding partition exists. For each edge \(xy\in E(X)\), with \(x\prec_X y\),
we have
\[
w_{xy}\prec x\prec y,
\qquad
w_{xy}\sim y,
\qquad
x\sim y.
\]
Avoiding \(\forst\) therefore forces \(x\) and \(y\) to lie in different parts. Hence the restriction of the
partition to \(V(X)\) is a proper \(k\)-coloring of \(X\).
\end{proof}





\begin{theorem}\label{thm:forst-general-np-complete}
The \(\forst\)-Avoidance problem is NP-complete.
More precisely, for every fixed \(k\ge 3\), deciding whether a graph admits a
\(\forst\)-avoiding solution with at most \(k\) parts is NP-complete.
\end{theorem}

\begin{proof}[Proof sketch]
We reduce from \(k\)-\textsc{Coloring}. Given an instance \(X\), for every vertex
\(x\in V(X)\) we create an independent block
\[
B_x:=\{x_1,x_2\}.
\]
For every edge \(xy\in E(X)\), we join \(B_x\) and \(B_y\) completely, and we add no other edges. Thus
\[
V(G)=\bigcup_{x\in V(X)} B_x,
\]
and
\[
E(G)=\{x_i y_j:\ xy\in E(X),\ i,j\in\{1,2\}\}.
\]

If \(X\) has a proper \(k\)-coloring, put both vertices of \(B_x\) in the color class of \(x\). Since adjacent
vertices of \(X\) receive different colors, no edge of \(G\) lies inside a part, and hence the solution avoids
\(\forst\) for any order.

Conversely, suppose that \(G\) has a \(\forst\)-avoiding solution
\[
S=(\prec,\mathcal Q),
\qquad |\mathcal Q|\le k.
\]
For each \(x\in V(X)\), let \(\ell_x\) be the last vertex of the block \(B_x=\{x_1,x_2\}\) in the order \(\prec\).
We color \(x\) with the part of \(\mathcal Q\) that contains \(\ell_x\).

We claim that this gives a proper \(k\)-coloring of \(X\). Let \(xy\in E(X)\). We need to show that
\(\ell_x\) and \(\ell_y\) lie in different parts of \(\mathcal Q\). Suppose, to the contrary, that they lie in the
same part. Without loss of generality, assume
\[
\ell_x \prec \ell_y.
\]
Let \(x'\) be the other vertex of \(B_x\). Since \(\ell_x\) is the last vertex of \(B_x\), we have
\[
x'\prec \ell_x\prec \ell_y.
\]
Moreover, because \(xy\in E(X)\), the construction makes \(B_x\) complete to \(B_y\). Hence
\[
x'\sim \ell_y
\qquad\text{and}\qquad
\ell_x\sim \ell_y.
\]
Finally, \(\ell_x\) and \(\ell_y\) lie in the same part by assumption. Therefore the ordered triple
\[
x'\prec \ell_x\prec \ell_y
\]
realizes the forbidden pattern \(\forst\), contradicting that \(S\) avoids \(\forst\).

Thus, for every edge \(xy\in E(X)\), the vertices \(x\) and \(y\) receive different colors. Hence \(X\) is
\(k\)-colorable.
\end{proof}




\begin{theorem}\label{thm:forst-given-partition-np-complete}
The \(\forst\)-Avoidance problem with a given partition is NP-complete.
\end{theorem}

\begin{proof}[Proof idea]
We reduce from \textsc{NAE-3SAT}\eric{Definirlo citando. cuando hagamos lo de references bib}. Let \(\Phi\) be an instance of \textsc{NAE-3SAT}.
A \emph{place} of \(\Phi\) is a pair
\[
p=(i,j),
\]
corresponding to the \(i\)-th literal of the clause \(C_j\).

For each place \(p\), we create two vertices
\[
T_p,\ F_p.
\]
They form one part of the fixed partition:
\[
V_p:=\{T_p,F_p\}.
\]
We add the internal edge
\[
T_pF_p.
\]
Thus the given partition is
\[
\mathcal Q:=\{V_p:\ p\text{ is a place of }\Phi\}.
\]

\medskip
\noindent\textbf{Consistency edges.}
Let \(p,q\) be two distinct places corresponding to occurrences of the same variable.

If the variable appears with the same polarity in \(p\) and \(q\) (both positive or both negative), we add the edges
\[
T_pT_q,
\qquad
F_pF_q.
\]
If the variable appears with opposite polarity in \(p\) and \(q\), we add the edges
\[
T_pF_q,
\qquad
F_pT_q.
\]
These edges are intended to force all occurrences of the same variable to be oriented consistently.

\medskip
\noindent\textbf{Clause edges.}
Let
\[
C_j=(L_{j1},L_{j2},L_{j3})
\]
be a clause, and let
\[
p_1=(1,j),\qquad p_2=(2,j),\qquad p_3=(3,j)
\]
be its three places. We add the cyclic edges
\[
T_{p_1}F_{p_2},
\qquad
T_{p_2}F_{p_3},
\qquad
T_{p_3}F_{p_1}.
\]

No other edges are added.

The intended correspondence is that an order of the vertices chooses, for each place \(p\), which of \(T_p\) and \(F_p\)
comes first. The consistency edges force equal variables to receive compatible choices, while the clause edges rule out
the two forbidden assignments in which all literals of a clause receive the same truth value.
Thus, \(\Phi\) is NAE-satisfiable if and only if the constructed instance admits an order avoiding \(\forst\) with
respect to the fixed partition \(\mathcal Q\).

DESPUES la escribo el ida y vuelta, que es medio largo
\end{proof}





\begin{theorem}\label{thm:forrt-given-order-np-complete}
The \(\forrt\)-Avoidance problem with a given order is NP-complete.
More precisely, for every fixed \(k\ge 3\), deciding whether there exists a partition into at most \(k\) parts
that avoids \(\forrt\) with respect to a given order is NP-complete.
\end{theorem}

\begin{proof}[Proof sketch]
We reduce from \(k\)-\textsc{Coloring}. Let \(X\) be an instance with
\[
V(X)=\{x_1,\dots,x_n\},
\]
ordered arbitrarily as
\[
x_1\prec x_2\prec\cdots\prec x_n.
\]
For each edge \(x_i x_j\in E(X)\), with \(i<j\), we create a witness vertex \(w_{ij}\). We keep the original
edge \(x_i x_j\), and we make \(w_{ij}\) adjacent only to \(x_j\). Thus,
\[
V(G)=V(X)\cup \{w_{ij}: x_i x_j\in E(X),\ i<j\},
\]
and
\[
E(G)=E(X)\cup \{w_{ij}x_j : x_i x_j\in E(X),\ i<j\}.
\]

For each \(j\), let
\[
W_j:=\{w_{ij}: i<j \text{ and } x_i x_j\in E(X)\}.
\]
The given order is
\[
x_1\prec W_2\prec x_2\prec W_3\prec x_3\prec\cdots\prec W_n\prec x_n,
\]
where the vertices inside each \(W_j\) are ordered arbitrarily. Hence, for every edge \(x_i x_j\in E(X)\), with \(i<j\),
we have
\[
x_i\prec w_{ij}\prec x_j,
\qquad
x_i\sim x_j,
\qquad
w_{ij}\sim x_j.
\]

If \(X\) has a proper \(k\)-coloring, keep the colors on the original vertices and assign each witness \(w_{ij}\) a color
different from the color of \(x_j\) (for instance, the color of \(x_i\)). This gives a partition that avoids \(\forrt\).

Conversely, if a \(\forrt\)-avoiding partition exists, then for every edge \(x_i x_j\in E(X)\), with \(i<j\), the triple
\[
x_i\prec w_{ij}\prec x_j
\]
forces \(x_i\) and \(x_j\) to lie in different parts. Therefore, the restriction of the partition to \(V(X)\) is a proper
\(k\)-coloring of \(X\).
\end{proof}

\begin{figure}[h]
\centering
\begin{tikzpicture}[
  scale=1,
  every node/.style={circle, draw, inner sep=1.4pt, font=\small}
]

\node (x) at (0,0) {$x_i$};
\node (w) at (2.2,0) {$w_{ij}$};
\node (y) at (4.4,0) {$x_j$};

\draw[bend left=35] (x) to (y);
\draw[bend right=35] (w) to (y);

\node[draw=none, font=\small] at (2.2,-1.1) {$x_i\prec w_{ij}\prec x_j$};

\end{tikzpicture}
\caption{Gadget used in the reduction for \(\forrt\) with a given order.}
\label{fig:forrt-given-order-gadget}
\end{figure}









\begin{theorem}\label{thm:rt-colored-equals-uncolored}
Let
\[
P=(A,N,\{r,t\})
\]
be a colored pattern, and consider the pattern
\[
P'=(A,N).
\]
Then, for every graph \(G\),
\[
\minClases_P(G)=\minClases_{P'}(G).
\]
\end{theorem}

\begin{proof}
The inequality
\[
\minClases_P(G)\ge \minClases_{P'}(G)
\]
follows from Theorem~\ref{thm:colored-lower-bound-by-class-partition}.

It remains to prove the reverse inequality. Let
\[
k=\minClases_{P'}(G).
\]
By definition of \(\minClases_{P'}\), there exists a partition
\[
V(G)=V_1\,\dot\cup\,\cdots\,\dot\cup\,V_k
\]
such that, for every \(i=1,\dots,k\), the graph \(G[V_i]\) belongs to the class associated with \(P'\).
Equivalently, for each \(i\), there exists an order \(\prec_i\) of \(V_i\) such that \(G[V_i]\) avoids the pattern
\(P'=(A,N)\).

We define a global order \(\prec\) by blocks:
\[
V_1\prec V_2\prec\cdots\prec V_k,
\]
using the order \(\prec_i\) inside each block \(V_i\). Consider the solution
\[
S=(\prec,\{V_1,\dots,V_k\}).
\]
We show that \(S\) avoids \(P\).

Let \(\varphi:\{r,s,t\}\to V(G)\) be an ordered occurrence, so
\[
\varphi(r)\prec \varphi(s)\prec \varphi(t).
\]

Suppose that \(\varphi(r)\) and \(\varphi(t)\) lie in the same part \(V_i\). Since the order \(\prec\) is by
blocks and
\[
\varphi(r)\prec \varphi(s)\prec \varphi(t),
\]
we also have \(\varphi(s)\in V_i\). Thus the whole occurrence lies inside \(G[V_i]\). But the order
\(\prec_i\) was chosen so that \(G[V_i]\) avoids \(P'\). Hence the occurrence fails one of the edge/non-edge
conditions of \(P'\), and therefore it also avoids \(P\).

Thus \(S\) avoids \(P\), and it has \(k\) parts. Hence
\[
\minClases_P(G)\le k=\minClases_{P'}(G).
\]
Combining both inequalities, we obtain
\[
\minClases_P(G)=\minClases_{P'}(G).
\]
\end{proof}











\begin{corollary}\label{lem:forrt-given-partition-forest}
Let \(G=(V,E)\) be a graph and let
\[
\mathcal Q=\{V_1,\dots,V_k\}
\]
be a fixed partition of \(V\). There exists an order \(\prec\) such that
\[
S=(\prec,\mathcal Q)
\]
avoids \(\forrt\) if and only if \(G[V_i]\) is a forest for every \(i=1,\dots,k\).
\end{corollary}

\eric{Corolario de lo de antes, no necesita demo. ahi esta ok? }




\german{CREO que hay que seguir chequeando desde aca}

\eric{El siguiente teorema es general para cualquier patron P que sea polinomial de calcular la clase (que creo que son todos).}

\begin{theorem}\label{thm:forrt-given-partition-poly}
The \(\forrt\)-Avoidance problem with a given partition can be solved in polynomial time.
\end{theorem}

\begin{proof}
ACA PODEMOS PONER EL ALGORITMO QUE YA HICIMOS 
By Lemma~\ref{lem:forrt-given-partition-forest}, it is enough to check whether \(G[V_i]\) is a forest for every
part \(V_i\in\mathcal Q\).

This can be done in polynomial time by running a DFS or BFS on each induced subgraph \(G[V_i]\) and checking whether
it is acyclic. Equivalently, one checks whether every connected component of \(G[V_i]\) has one fewer edge than
vertices.

Therefore, the algorithm returns \textbf{YES} if and only if each \(G[V_i]\) is a forest, which is equivalent to the
existence of an order \(\prec\) such that \((\prec,\mathcal Q)\) avoids \(\forrt\).
\end{proof}





\begin{theorem}\label{thm:forrs-st-given-order-np-complete}
The \(\mathcal F_{rs\text{-}st}\)-Avoidance problem with a given order is NP-complete.
More precisely, for every fixed \(k\ge 3\), deciding whether there exists a partition into at most \(k\) parts
that avoids both \(\forrs\) and \(\forst\) with respect to a given order is NP-complete.
\end{theorem}

\begin{proof}[Proof sketch]
We reduce from \(k\)-\textsc{Coloring}. Let \(X\) be an instance with
\[
V(X)=\{x_1,\dots,x_n\}.
\]
For each edge \(x_i x_j\in E(X)\), with \(i<j\), add one witness vertex \(w_{ij}\), adjacent exactly to
\(x_i\) and \(x_j\), and do not keep the original edge \(x_i x_j\). Thus,
\[
V(G)=V(X)\cup \{w_{ij}: x_i x_j\in E(X),\ i<j\},
\]
and
\[
E(G)=\{x_iw_{ij},x_jw_{ij}: x_i x_j\in E(X),\ i<j\}.
\]
The given order is
\[
x_1\prec x_2\prec \cdots \prec x_n \prec \text{all vertices } w_{ij},
\]
where the witnesses are ordered arbitrarily.

If \(X\) has a proper \(k\)-coloring, keep the colors on the original vertices.
For each edge \(x_i x_j\in E(X)\), with \(i<j\), assign the witness \(w_{ij}\) the color of \(x_i\).
We use the fixed order
\[
x_1\prec x_2\prec\cdots\prec x_n\prec \text{all vertices }w_{ij}.
\]

We claim that the resulting solution avoids both \(\forrs\) and \(\forst\).
Indeed, the only possible relevant triples are of the form
\[
x_i\prec x_j\prec w_{ij},
\]
with
\[
x_i\sim w_{ij}
\qquad\text{and}\qquad
x_j\sim w_{ij}.
\]
Since \(x_i x_j\in E(X)\), the proper coloring of \(X\) gives that \(x_i\) and \(x_j\) lie in different parts.
Thus the triple avoids \(\forrs\).
Moreover, \(w_{ij}\) was assigned the color of \(x_i\), so \(w_{ij}\) and \(x_j\) also lie in different parts.
Thus the triple avoids \(\forst\).
Therefore the constructed solution avoids \(\mathcal F_{rs\text{-}st}\).
Conversely, suppose that there is a solution avoiding both \(\forrs\) and \(\forst\) with at most \(k\) parts.
For each edge \(x_i x_j\in E(X)\), with \(i<j\), we have
\[
x_i\prec x_j\prec w_{ij},
\qquad
x_i\sim w_{ij},
\qquad
x_j\sim w_{ij}.
\]
Avoiding \(\forrs\) forces \(x_i\) and \(x_j\) to lie in different parts. Therefore, the restriction of the
partition to \(V(X)\) is a proper \(k\)-coloring of \(X\).
\end{proof}



\begin{theorem}\label{thm:forrs-st-general-np-complete}
The \(\mathcal F_{rs\text{-}st}\)-Avoidance problem is NP-complete.
More precisely, for every fixed \(k\ge 3\), deciding whether a graph admits a solution with at most \(k\) parts
that avoids both \(\forrs\) and \(\forst\) is NP-complete.
\end{theorem}

\begin{proof}[Proof sketch]
We reduce from \(k\)-\textsc{Coloring}. Fix an arbitrary order of the vertices of \(X\). For each edge
\(xy\in E(X)\), with \(x\) before \(y\) in this order, we add \(k+1\) common witnesses
\[
W_{xy}:=\{w^1_{xy},\dots,w^{k+1}_{xy}\},
\]
each adjacent exactly to \(x\) and \(y\). The original edge \(xy\) is not kept. Thus
\[
V(G)=V(X)\cup \bigcup_{xy\in E(X)} W_{xy},
\]
and
\[
E(G)=\{xw^i_{xy},yw^i_{xy}: xy\in E(X),\ 1\le i\le k+1\}.
\]

If \(X\) has a proper \(k\)-coloring, keep the colors on \(V(X)\). Choose an order with all original vertices first
and all witnesses afterwards. For each edge \(xy\), with \(x\) before \(y\), assign every witness \(w^i_{xy}\) the
color of \(x\). Then the triple \(x\prec y\prec w^i_{xy}\) avoids \(\forrs\) because \(x\) and \(y\) have different
colors, and avoids \(\forst\) because \(y\) and \(w^i_{xy}\) have different colors.

Conversely, suppose that there is a solution avoiding both \(\forrs\) and \(\forst\) with at most \(k\) parts.
If two adjacent original vertices \(x,y\) were in the same part, assume \(x\prec y\). Then no witness \(w^i_{xy}\)
can appear after \(y\), since \(x\prec y\prec w^i_{xy}\) would realize \(\forrs\). Hence all \(k+1\) witnesses for
\(xy\) appear before \(y\). By the pigeonhole principle, two of them lie in the same part; together with \(y\), they
realize \(\forrs\), a contradiction. Therefore adjacent vertices of \(X\) lie in different parts, and the restriction
to \(V(X)\) is a proper \(k\)-coloring of \(X\).
\end{proof}



\medskip

\begin{theorem}\label{thm:forrs-st-given-partition-np-complete}
The \(\mathcal F_{rs\text{-}st}\)-Avoidance problem with a given partition is NP-complete.
\end{theorem}

\begin{proof}[Proof sketch]
We reduce from the following NP-complete variant of \textsc{Betweenness with precedence constraints}:
an instance is \((A,B,T)\), where \(A\) is a set of elements, \(B\) is a set of precedence constraints
\((u,v)\) requiring \(u\prec v\), and \(T\) is a set of disjoint triples \((a,b,c)\) requiring \(b\) to lie
between \(a\) and \(c\), that is, either \(a\prec b\prec c\) or \(c\prec b\prec a\).

First, for each \((u,v)\in B\), add a new element \(z_{uv}\) and replace \(u\prec v\) by
\[
u\prec z_{uv},
\qquad
z_{uv}\prec v.
\]
Let \((A',B',T)\) be the resulting instance.

We construct a graph \(G\) and a fixed partition \(\mathcal Q\).

The base vertices of \(G\) are the elements of \(A'\). For every pair \((u,v)\in B'\), add two auxiliary vertices
\(p_{uv}\) and \(q_{uv}\). Thus
\[
V(G)=A'\cup \{p_{uv},q_{uv}:(u,v)\in B'\}.
\]

The partition \(\mathcal Q\) is defined as follows. For every triple \((a,b,c)\in T\), put \(a\) and \(c\) in the
same part, and put \(b\) in a different part. Since the triples are disjoint, this creates no conflicts.
Every base vertex not appearing in a triple is put in its own part. Finally, for each \((u,v)\in B'\), put
\(p_{uv}\) and \(q_{uv}\) in the same part as \(u\).

The edges of \(G\) are the following. For every triple \((a,b,c)\in T\), add the triangle
\[
ab,\quad bc,\quad ac.
\]
For every pair \((u,v)\in B'\), add
\[
up_{uv},\quad vp_{uv},\quad uq_{uv},\quad vq_{uv}.
\]
No other edges are added.

The triple gadgets force betweenness: since \(a\) and \(c\) are in the same part and \(a,b,c\) form a triangle,
placing \(b\) after both \(a,c\) creates an occurrence of \(\forrs\), while placing \(b\) before both \(a,c\)
creates an occurrence of \(\forst\). Hence \(b\) must lie between \(a\) and \(c\).

The pair gadgets force precedence. For \((u,v)\in B'\), the vertices \(p_{uv},q_{uv}\) are in the part of \(u\) and
are common neighbors of \(u\) and \(v\). If an order had \(v\prec u\), then either one auxiliary appears after \(u\),
creating an occurrence of \(\forst\), or both auxiliaries appear before \(u\), in which case the two auxiliaries together
with \(u\) create an occurrence of \(\forrs\). Therefore every valid order must satisfy \(u\prec v\).

Conversely, any order satisfying \((A',B',T)\) can be extended to the auxiliary vertices by placing, for each
\((u,v)\in B'\),
\[
u\prec \cdots \prec v \prec p_{uv}\prec q_{uv}.
\]
The subdivision step ensures that \(u\) and \(v\) are in different parts of \(\mathcal Q\), so the pair gadgets do not
create occurrences of either \(\forrs\) or \(\forst\). The triple gadgets are valid because each \(b\) lies between
the corresponding \(a\) and \(c\).

Thus the original instance is satisfiable if and only if the constructed instance \((G,\mathcal Q)\) admits an order
avoiding both \(\forrs\) and \(\forst\). The construction is polynomial, and the problem is in NP because an order can
be checked in polynomial time. Hence the problem is NP-complete.
\end{proof}














\subsubsection{Families containing \(\operatorname{for}_{12}\): given order}

For \(D\subseteq\{12,13,23\}\), we write
\[
\mathcal F_D:=\{\operatorname{for}_{ab}:ab\in D\}.
\]

\paragraph{Common construction.}
Fix \(k\ge 3\), and let \(X\) be an instance of \(k\)-\textsc{Coloring}, with
\[
V(X)=\{x_1,\dots,x_n\}.
\]
For every edge \(x_i x_j\in E(X)\), with \(i<j\), we add one witness vertex \(w_{ij}\), adjacent exactly to
\(x_i\) and \(x_j\). We do not keep the original edge \(x_i x_j\). Thus,
\[
V(G)=V(X)\cup \{w_{ij}:x_i x_j\in E(X),\ i<j\},
\]
and
\[
E(G)=\{x_iw_{ij},x_jw_{ij}:x_i x_j\in E(X),\ i<j\}.
\]
The given order \(\prec\) is
\[
x_1\prec x_2\prec\cdots\prec x_n\prec \text{all witnesses }w_{ij},
\]
where the witnesses are ordered arbitrarily.

In this ordered graph, the only triples with the two required edges of a forest pattern are
\[
x_i\prec x_j\prec w_{ij}
\]
for some edge \(x_i x_j\in E(X)\), with \(i<j\).






\begin{figure}[h]
\centering
\begin{tikzpicture}[
  scale=1,
  every node/.style={circle, draw, inner sep=1.4pt, font=\small}
]

\node (x) at (0,0) {$x_i$};
\node (y) at (2.2,0) {$x_j$};
\node (w) at (4.4,0) {$w_{ij}$};

\draw[bend left=35] (x) to (w);
\draw (y)--(w);

\node[draw=none, font=\small] at (2.2,-1.1) {$x_i\prec x_j\prec w_{ij}$};

\end{tikzpicture}
\caption{Gadget used in the construction of \(G\).}
\label{fig:forrs-rt-given-order-gadget}
\end{figure}







\eric{Te permito que sea un lema}
\begin{lemma}\label{obs:common-12-given-order-triples} \german{puede q esto se use en mas lados}
In the common construction above, let
\[
u\prec v\prec z
\]
be an ordered triple such that
\[
u\sim z
\qquad\text{and}\qquad
v\sim z.
\]
Then there exists an edge \(x_i x_j\in E(X)\), with \(i<j\), and a witness \(w_{ij}\) such that
\[
(u,v,z)=(x_i,x_j,w_{ij}).
\]


\end{lemma}

\begin{proof}
There are no edges among vertices of \(V(X)\), and there are no edges among witnesses. Moreover, every witness
comes after all original vertices in the order. Therefore, if \(z\) has two earlier neighbors, then \(z\) must be
a witness vertex \(w_{ij}\). By construction, the only neighbors of \(w_{ij}\) are \(x_i\) and \(x_j\), and since
\(i<j\), they appear in the order as
\[
x_i\prec x_j\prec w_{ij}.
\]
Thus \((u,v,z)=(x_i,x_j,w_{ij})\).
\end{proof}


\begin{lemma}\label{lem:given-order-12-common-converse} \german{la vuelta comun}
Let \(D\subseteq\{12,13,23\}\) with \(12\in D\). If the ordered graph \((G,\prec)\) constructed above admits a
partition into at most \(k\) parts avoiding \(\mathcal F_D\), then \(X\) is \(k\)-colorable.
\end{lemma}

\begin{proof}
Since \(12\in D\), every solution avoiding \(\mathcal F_D\) also avoids \(\operatorname{for}_{12}\).
Let
\[
S=(\prec,\mathcal Q)
\]
be such a solution, with \(|\mathcal Q|\le k\).

We show that the restriction of \(\mathcal Q\) to \(V(X)\) is a proper coloring of \(X\).
Let \(x_i x_j\in E(X)\), with \(i<j\). By construction, there exists a witness \(w_{ij}\) adjacent to both \(x_i\) and \(x_j\), and the order is
\[
x_i\prec x_j\prec w_{ij}.
\]
If \(x_i\) and \(x_j\) were assigned to the same part of \(\mathcal Q\), then this ordered triple

would not avoid \(\operatorname{for}_{12}\), contradicting the fact that \(S\) avoids
\(\mathcal F_D\).

Therefore, for every edge \(x_i x_j\in E(X)\), the vertices \(x_i\) and \(x_j\) are assigned to different parts of
\(\mathcal Q\). Thus the restriction of \(\mathcal Q\) to \(V(X)\) gives a proper coloring of \(X\) with at most
\(k\) colors.
\end{proof}

\medskip




\begin{theorem}\label{thm:F12-given-order-np-complete}
For every fixed \(k\ge 3\), the \(\mathcal F_{12}\)-Avoidance problem with a given order is NP-complete.
\end{theorem}


\begin{proof}
The problem is in NP. For NP-hardness, use the common construction above.

Suppose that \(X\) has a proper \(k\)-coloring \(\varphi\). We define a partition of \(V(G)\) into at most \(k\)
parts by assigning each original vertex \(x_k\) to the part \(\varphi(x_k)\), and assigning each witness \(w_{ij}\)
arbitrarily.

We claim that this partition, together with the fixed order \(\prec\), avoids \(\mathcal F_{12}\).
Let
\[
u\prec v\prec z
\]
be an ordered triple such that \(u\) and \(v\) are both adjacent to \(z\), by Observation~\ref{obs:common-12-given-order-triples}, there is an edge
\(x_i x_j\in E(X)\), with \(i<j\), and a witness \(w_{ij}\) such that
\[
(u,v,z)=(x_i,x_j,w_{ij}).
\]
Since \(\varphi\) is a proper coloring of \(X\), the vertices \(x_i\) and \(x_j\) are assigned to different parts.
Therefore by the \colorcondition, the triple avoids
\(\operatorname{for}_{12}\). 

Thus the constructed solution avoids \(\mathcal F_{12}\).

The converse follows from Lemma~\ref{lem:given-order-12-common-converse} with \(D=\{12\}\).
\end{proof}



\begin{theorem}\label{thm:F12-13-given-order-np-complete}
For every fixed \(k\ge 3\), the \(\mathcal F_{12\text{-}13}\)-Avoidance problem with a given order is NP-complete.
\end{theorem}

\begin{proof}
The problem is in NP. For NP-hardness, use the common construction above.

Suppose that \(X\) has a proper \(k\)-coloring \(\varphi\). We assign each original vertex \(x_k\) to the part
\(\varphi(x_k)\). For each edge \(x_i x_j\in E(X)\), with \(i<j\), we assign the witness \(w_{ij}\) to the part
\[
\varphi(x_j).
\]

We claim that the resulting solution avoids both \(\operatorname{for}_{12}\) and \(\operatorname{for}_{13}\).

Let
\[
u\prec v\prec z
\]
be an ordered triple such that \(u\) and \(v\) are both adjacent to \(z\), by Observation~\ref{obs:common-12-given-order-triples}, there is an edge
\(x_i x_j\in E(X)\), with \(i<j\), and a witness \(w_{ij}\) such that
\[
(u,v,z)=(x_i,x_j,w_{ij}).
\]
Since \(\varphi\) is proper, \(x_i\) and \(x_j\) are assigned to different parts. Hence the \colorcondition
\(C=\{1,2\}\) of \(\operatorname{for}_{12}\) activates.

Moreover, \(w_{ij}\) is assigned to the part \(\varphi(x_j)\), while \(x_i\) is assigned to the part
\(\varphi(x_i)\). Since \(\varphi(x_i)\neq\varphi(x_j)\), the vertices \(x_i\) and \(w_{ij}\) are assigned to
different parts. Hence by the \colorcondition  condition the solution avoids \(\operatorname{for}_{13}\).

Thus the solution avoids \(\mathcal F_{12\text{-}13}\). \german{asi? }

The converse follows from Lemma~\ref{lem:given-order-12-common-converse} with \(D=\{12,13\}\).
\end{proof}





\begin{theorem}\label{thm:F12-23-given-order-np-complete}
For every fixed \(k\ge 3\), the \(\mathcal F_{12\text{-}23}\)-Avoidance problem with a given order is NP-complete.
\end{theorem}


\begin{proof}
The problem is in NP. For NP-hardness, use the common construction above.

Suppose that \(X\) has a proper \(k\)-coloring \(\varphi\). We assign each original vertex \(x_k\) to the part
\(\varphi(x_k)\). For each edge \(x_i x_j\in E(X)\), with \(i<j\), we assign the witness \(w_{ij}\) to the part
\[
\varphi(x_i).
\]

We claim that the resulting solution avoids both \(\operatorname{for}_{12}\) and \(\operatorname{for}_{23}\).
Let
\[
u\prec v\prec z
\]
be an ordered triple such that \(u\) and \(v\) are both adjacent to \(z\), by Observation~\ref{obs:common-12-given-order-triples}, there is an edge
\(x_i x_j\in E(X)\), with \(i<j\), and a witness \(w_{ij}\) such that
\[
(u,v,z)=(x_i,x_j,w_{ij}).
\]
Since \(\varphi\) is proper, \(x_i\) and \(x_j\) are assigned to different parts. Hence the color condition
\(C=\{1,2\}\) of \(\operatorname{for}_{12}\) activates.

Moreover, \(w_{ij}\) is assigned to the part \(\varphi(x_i)\), while \(x_j\) is assigned to the part
\(\varphi(x_j)\). Since \(\varphi(x_i)\neq\varphi(x_j)\), the vertices \(x_j\) and \(w_{ij}\) are assigned to
different parts. Hence the color condition \(C=\{2,3\}\) of \(\operatorname{for}_{23}\) fails.

Thus the solution avoids \(\mathcal F_{12\text{-}23}\).

The converse follows from Lemma~\ref{lem:given-order-12-common-converse} with \(D=\{12,23\}\).
\end{proof}





\begin{theorem}\label{thm:F12-23-13-given-order-np-complete}
For every fixed \(k\ge 3\), the \(\mathcal F_{12\text{-}23\text{-}13}\)-Avoidance problem with a given order is NP-complete.
\end{theorem}


\begin{proof}
The problem is in NP. For NP-hardness, use the common construction above.

Suppose that \(X\) has a proper \(k\)-coloring \(\varphi\). We assign each original vertex \(x_k\) to the part
\(\varphi(x_k)\). For each edge \(x_i x_j\in E(X)\), with \(i<j\), choose a part different from both
\[
\varphi(x_i)
\qquad\text{and}\qquad
\varphi(x_j),
\]
and assign \(w_{ij}\) to that part. This is possible because \(k\ge 3\) and
\[
\varphi(x_i)\neq\varphi(x_j).
\]

We claim that the resulting solution avoids
\[
\mathcal F_{12\text{-}23\text{-}13}.
\]
Let
\[
u\prec v\prec z
\]
be an ordered triple such that \(u\) and \(v\) are both adjacent to \(z\), by Observation~\ref{obs:common-12-given-order-triples}, there is an edge
\(x_i x_j\in E(X)\), with \(i<j\), and a witness \(w_{ij}\) such that
\[
(u,v,z)=(x_i,x_j,w_{ij}).
\]
The choice of the part of \(w_{ij}\), together with the fact that \(\varphi\) is proper, implies that the three
vertices \(x_i\), \(x_j\), and \(w_{ij}\) are assigned to pairwise different parts. Therefore each of the color
conditions \(C=\{1,2\}\), \(C=\{1,3\}\), and \(C=\{2,3\}\) activates. Hence the triple avoids
\(\operatorname{for}_{12}\), \(\operatorname{for}_{13}\), and \(\operatorname{for}_{23}\).

Thus the solution avoids \(\mathcal F_{12\text{-}23\text{-}13}\).

The converse follows from Lemma~\ref{lem:given-order-12-common-converse} with \(D=\{12,23,13\}\).
\end{proof}







\subsubsection{Families containing \(\operatorname{for}_{12}\): general case}

\paragraph{Common construction.}
Fix \(k\ge 3\), and let \(X\) be an instance of \(k\)-\textsc{Coloring}.
For every edge \(xy\in E(X)\), we add \(k+1\) witness vertices
\[
W_{xy}:=\{w^1_{xy},\dots,w^{k+1}_{xy}\},
\]
each adjacent exactly to \(x\) and \(y\). We do not keep the original edge \(xy\). Thus,
\[
V(G)=V(X)\cup \bigcup_{xy\in E(X)} W_{xy},
\]
and
\[
E(G)=\{xw^a_{xy},yw^a_{xy}:xy\in E(X),\ a=1,\dots,k+1\}.
\]

When \(X\) is properly \(k\)-colored, we will use the following order in the forward direction. We choose an arbitrary
order of the original vertices, put all original vertices first, and then put all witness vertices.


\begin{figure}[ht]
\centering
\begin{tikzpicture}[
  scale=1,
  every node/.style={circle, draw, inner sep=1.4pt, font=\small}
]

\node (x) at (0,0) {$x_i$};
\node (y) at (2.2,0) {$x_j$};

\node (w1) at (4.8,1.2) {$w^{1}_{ij}$};
\node (w2) at (4.8,0.2) {$w^{2}_{ij}$};
\node[draw=none] (vd) at (4.8,-0.8) {$\vdots$};
\node (wk) at (4.8,-1.8) {$w^{k+1}_{ij}$};

\draw[bend left=20] (x) to (w1);
\draw (y) -- (w1);

\draw[bend left=10] (x) to (w2);
\draw (y) -- (w2);

\draw[bend right=10] (x) to (wk);
\draw (y) -- (wk);

\node[draw=none, font=\small] at (2.6,-3.0)
{$x_i\prec x_j\prec w^{1}_{ij},w^{2}_{ij},\dots,w^{k+1}_{ij}$};

\end{tikzpicture}
\caption{Gadget used in the general construction: for each edge \(x_ix_j\), we add \(k+1\) common witnesses.}
\label{fig:forrs-general-gadget}
\end{figure}





\begin{remark}\label{obs:common-12-general-forward-triples}
Consider the forward order described above. Let
\[
u\prec v\prec z
\]
be an ordered triple such that
\[
u\sim z
\qquad\text{and}\qquad
v\sim z.
\]
Then there exists an edge \(xy\in E(X)\), with \(x\prec y\), and an index \(a\in\{1,\dots,k+1\}\), such that
\[
(u,v,z)=(x,y,w^a_{xy}).
\]
\end{remark}

\begin{proof}
There are no edges among original vertices and no edges among witness vertices. Moreover, in the forward order all
witnesses appear after all original vertices. Hence, if \(z\) has two earlier neighbors, then \(z\) must be a witness
\(w^a_{xy}\). Its only neighbors are \(x\) and \(y\), and these appear before \(w^a_{xy}\) in the chosen order. If
\(x\prec y\), then the triple is exactly
\[
x\prec y\prec w^a_{xy}.
\]
\end{proof}

\begin{lemma}\label{lem:general-12-common-converse}
Let \(D\subseteq\{12,13,23\}\) with \(12\in D\). If the graph \(G\) constructed above admits a solution with at most
\(k\) parts avoiding \(\mathcal F_D\), then \(X\) is \(k\)-colorable.
\end{lemma}

\begin{proof}
Since \(12\in D\), every solution avoiding \(\mathcal F_D\) also avoids \(\operatorname{for}_{12}\).
Let
\[
S=(\prec,\mathcal Q)
\]
be such a solution, with \(|\mathcal Q|\le k\).

We show that the restriction of \(\mathcal Q\) to \(V(X)\) is a proper coloring of \(X\). Let \(xy\in E(X)\).
Suppose, that \(x\) and \(y\) are assigned to the same part of \(\mathcal Q\).
Without loss of generality, assume
\[
x\prec y.
\]

No witness \(w^a_{xy}\in W_{xy}\) can appear after \(y\). Indeed, if
\[
x\prec y\prec w^a_{xy},
\]
then the triple \(x\prec y\prec w^a_{xy}\) does not avoid \(\operatorname{for}_{12}\) because \(w^a_{xy}\) is adjacent to both \(x\) and \(y\).

Therefore, all \(k+1\) witnesses in \(W_{xy}\) appear before \(y\). Since there are at most \(k\) parts,  by the pigeonhole principle, two of them,
say \(w^a_{xy}\) and \(w^b_{xy}\), are assigned to the same part of \(\mathcal Q\). Suppose they appear in the order
\[
w^a_{xy}\prec w^b_{xy}\prec y.
\]
Both witnesses are adjacent to \(y\), and they are assigned to the same part. Hence the triple
\[
w^a_{xy}\prec w^b_{xy}\prec y
\]
realizes \(\operatorname{for}_{12}\), again a contradiction.

Thus, for every edge \(xy\in E(X)\), the vertices \(x\) and \(y\) are assigned to different parts. Hence the
restriction of \(\mathcal Q\) to \(V(X)\) gives a proper coloring of \(X\) with at most \(k\) colors.
\end{proof}





\begin{theorem}\label{thm:F12-general-np-complete}
For every fixed \(k\ge 3\), the \(\mathcal F_{12}\)-Avoidance problem is NP-complete.
\end{theorem}

\begin{proof}
The problem is in NP. For NP-hardness, use the common construction above.

Suppose that \(X\) has a proper \(k\)-coloring \(\varphi\). We choose the forward order described above. We assign
each original vertex \(x\in V(X)\) to the part \(\varphi(x)\), and we assign each witness vertex arbitrarily.

We claim that the resulting solution avoids \(\mathcal F_{12}\). Let
\[
u\prec v\prec z
\]
be an ordered triple such that \(u\) and \(v\) are both adjacent to \(z\). By Observation~\ref{obs:common-12-general-forward-triples}, there exist
\(xy\in E(X)\), with \(x\prec y\), and \(a\in\{1,\dots,k+1\}\), such that
\[
(u,v,z)=(x,y,w^a_{xy}).
\]
Since \(\varphi\) is a proper coloring of \(X\), the vertices \(x\) and \(y\) are assigned to different parts.
Therefore the color condition \(C=\{1,2\}\) of \(\operatorname{for}_{12}\) fails, and the triple avoids
\(\operatorname{for}_{12}\).

Thus the constructed solution avoids \(\mathcal F_{12}\).

The converse follows from Lemma~\ref{lem:general-12-common-converse} with \(D=\{12\}\).
\end{proof}



\begin{theorem}\label{thm:F12-13-general-np-complete}
For every fixed \(k\ge 3\), the \(\mathcal F_{12\text{-}13}\)-Avoidance problem is NP-complete.
\end{theorem}

\begin{proof}
The problem is in NP. For NP-hardness, use the common construction above.

Suppose that \(X\) has a proper \(k\)-coloring \(\varphi\). We choose the forward order described above. We assign
each original vertex \(x\in V(X)\) to the part \(\varphi(x)\). For each edge \(xy\in E(X)\), with \(x\prec y\), we
assign every witness \(w^a_{xy}\in W_{xy}\) to the part
\[
\varphi(y).
\]

We claim that the resulting solution avoids both \(\operatorname{for}_{12}\) and \(\operatorname{for}_{13}\).
Let
\[
u\prec v\prec z
\]
be an ordered triple such that \(u\) and \(v\) are both adjacent to \(z\).
By Observation~\ref{obs:common-12-general-forward-triples}, there exist \(xy\in E(X)\), with \(x\prec y\),
and \(a\in\{1,\dots,k+1\}\), such that
\[
(u,v,z)=(x,y,w^a_{xy}).
\]
Since \(\varphi\) is proper, \(x\) and \(y\) are assigned to different parts, so the color condition \(C=\{1,2\}\)
of \(\operatorname{for}_{12}\) fails. Moreover, \(w^a_{xy}\) is assigned to the part \(\varphi(y)\), while \(x\) is
assigned to the part \(\varphi(x)\); hence \(x\) and \(w^a_{xy}\) are assigned to different parts, so the color
condition \(C=\{1,3\}\) of \(\operatorname{for}_{13}\) fails.

Thus the constructed solution avoids \(\mathcal F_{12\text{-}13}\).

The converse follows from Lemma~\ref{lem:general-12-common-converse} with \(D=\{12,13\}\).
\end{proof}





\begin{theorem}\label{thm:F12-23-general-np-complete}
For every fixed \(k\ge 3\), the \(\mathcal F_{12\text{-}23}\)-Avoidance problem is NP-complete.
\end{theorem}

\begin{proof}
The problem is in NP. For NP-hardness, use the common construction above.

Suppose that \(X\) has a proper \(k\)-coloring \(\varphi\). We choose the forward order described above. We assign
each original vertex \(x\in V(X)\) to the part \(\varphi(x)\). For each edge \(xy\in E(X)\), with \(x\prec y\), we
assign every witness \(w^a_{xy}\in W_{xy}\) to the part
\[
\varphi(x).
\]

We claim that the resulting solution avoids both \(\operatorname{for}_{12}\) and \(\operatorname{for}_{23}\).
Let
\[
u\prec v\prec z
\]
be an ordered triple such that \(u\) and \(v\) are both adjacent to \(z\).
By Observation~\ref{obs:common-12-general-forward-triples}, there exist \(xy\in E(X)\), with \(x\prec y\),
and \(a\in\{1,\dots,k+1\}\), such that
\[
(u,v,z)=(x,y,w^a_{xy}).
\]
Since \(\varphi\) is proper, \(x\) and \(y\) are assigned to different parts, so the color condition \(C=\{1,2\}\)
of \(\operatorname{for}_{12}\) fails. Moreover, \(w^a_{xy}\) is assigned to the part \(\varphi(x)\), while \(y\) is
assigned to the part \(\varphi(y)\); hence \(y\) and \(w^a_{xy}\) are assigned to different parts, so the color
condition \(C=\{2,3\}\) of \(\operatorname{for}_{23}\) fails.

Thus the constructed solution avoids \(\mathcal F_{12\text{-}23}\).

The converse follows from Lemma~\ref{lem:general-12-common-converse} with \(D=\{12,23\}\).
\end{proof}







\begin{theorem}\label{thm:F12-23-13-general-np-complete}
For every fixed \(k\ge 3\), the \(\mathcal F_{12\text{-}23\text{-}13}\)-Avoidance problem is NP-complete.
\end{theorem}

\begin{proof}
The problem is in NP. For NP-hardness, use the common construction above.

Suppose that \(X\) has a proper \(k\)-coloring \(\varphi\). We choose the forward order described above. We assign
each original vertex \(x\in V(X)\) to the part \(\varphi(x)\). For each edge \(xy\in E(X)\), with \(x\prec y\), choose
a part different from both
\[
\varphi(x)
\qquad\text{and}\qquad
\varphi(y),
\]
and assign every witness \(w^a_{xy}\in W_{xy}\) to that part. This is possible because \(k\ge 3\) and
\(\varphi(x)\neq\varphi(y)\).

We claim that the resulting solution avoids
\[
\mathcal F_{12\text{-}23\text{-}13}.
\]
Let
\[
u\prec v\prec z
\]
be an ordered triple such that \(u\) and \(v\) are both adjacent to \(z\). If \(u\) and \(v\) are not both adjacent to \(z\). By Observation~\ref{obs:common-12-general-forward-triples}, there exist \(xy\in E(X)\), with \(x\prec y\),
and \(a\in\{1,\dots,k+1\}\), such that
\[
(u,v,z)=(x,y,w^a_{xy}).
\]
By construction, the vertices \(x\), \(y\), and \(w^a_{xy}\) are assigned to pairwise different parts. Therefore the
color conditions \(C=\{1,2\}\), \(C=\{1,3\}\), and \(C=\{2,3\}\) all fail. Hence the triple avoids
\(\operatorname{for}_{12}\), \(\operatorname{for}_{13}\), and \(\operatorname{for}_{23}\).

Thus the constructed solution avoids \(\mathcal F_{12\text{-}23\text{-}13}\).

The converse follows from Lemma~\ref{lem:general-12-common-converse} with \(D=\{12,23,13\}\).
\end{proof}










\begin{theorem}\label{thm:forrs-rt-given-partition-np-complete}
The \(\mathcal F_{rs\text{-}rt}\)-Avoidance problem with a given partition is NP-complete.
\end{theorem}

\begin{proof}[Proof sketch]
We reduce from the following NP-complete ordering problem. An instance is a triple
\[
(A,B,T),
\]
where \(A\) is a set of elements, \(B\) is a set of precedence constraints \((u,v)\), requiring
\(u\prec v\), and \(T\) is a family of disjoint triples \((a,b,c)\), requiring that the order
\(a\prec b\prec c\) does \emph{not} occur. In the notation of Guttmann--Maucher, this corresponds to family \(4\) with
\[
P=S_3\setminus\{(123)\}.
\]

First, we subdivide the precedence constraints. For every \((u,v)\in B\), add a new element \(z_{uv}\) and replace
\(u\prec v\) by the two constraints
\[
u\prec z_{uv},
\qquad
z_{uv}\prec v.
\]
Let \((A',B',T)\) be the resulting instance. This transformation preserves satisfiability.

We now construct a graph \(G\) and a fixed partition \(\mathcal Q\). The base vertices of \(G\) are the elements of
\(A'\). For every pair \((u,v)\in B'\), add two auxiliary vertices \(p_{uv}\) and \(q_{uv}\). Thus
\[
V(G)=A'\cup \{p_{uv},q_{uv}:(u,v)\in B'\}.
\]

The partition \(\mathcal Q\) is defined as follows. For every triple \((a,b,c)\in T\), put \(a\) and \(c\) in the
same part, and put \(b\) in a different part. Since the triples are disjoint, this creates no conflicts. Every base
vertex not appearing in a triple is put in its own part. Finally, for each \((u,v)\in B'\), put \(p_{uv}\) and
\(q_{uv}\) in the same part as \(v\).

The edges of \(G\) are the following. For every triple \((a,b,c)\in T\), add
\[
ac,\qquad bc,
\]
and do not add the edge \(ab\). For every pair \((u,v)\in B'\), add
\[
up_{uv},\quad vp_{uv},\quad uq_{uv},\quad vq_{uv}.
\]
No other edges are added.

The triple gadget forbids the order \(a\prec b\prec c\): if this order occurred, then the triple
\(a\prec b\prec c\) would have the required edges \(ac\) and \(bc\), and \(a,c\) would lie in the same part, giving an
occurrence of \(\forrt\).

The pair gadget forces \(u\prec v\). Indeed, suppose that a valid order had \(v\prec u\). If one of the auxiliaries
appears after \(u\), say
\[
v\prec u\prec p_{uv},
\]
then this triple realizes \(\forrt\), since \(v\sim p_{uv}\), \(u\sim p_{uv}\), and \(v,p_{uv}\) lie in the same part.
If both auxiliaries appear before \(u\), then two of them are ordered, say
\[
p_{uv}\prec q_{uv}\prec u,
\]
and this triple realizes \(\forrs\), since \(p_{uv},q_{uv}\) lie in the same part and both are adjacent to \(u\).
Thus every valid order must satisfy \(u\prec v\).

Conversely, assume that the ordering instance \((A',B',T)\) is satisfiable. Extend a satisfying order to the auxiliary
vertices by placing, for every \((u,v)\in B'\),
\[
u\prec \cdots \prec v \prec p_{uv}\prec q_{uv}.
\]
The subdivision step ensures that the endpoints \(u\) and \(v\) of every pair in \(B'\) lie in different parts of
\(\mathcal Q\). With this extension, the pair gadgets do not create occurrences of either \(\forrs\) or \(\forrt\),
and the triple gadgets are valid because the order \(a\prec b\prec c\) is avoided for every \((a,b,c)\in T\).

Therefore, the original ordering instance is satisfiable if and only if the constructed instance \((G,\mathcal Q)\)
admits an order avoiding both \(\forrs\) and \(\forrt\). The construction is polynomial. The problem is in NP because
a certificate is an order \(\prec\), and checking whether \((\prec,\mathcal Q)\) avoids both patterns can be done in
polynomial time. Hence the problem is NP-complete.
\end{proof}







\begin{theorem}\label{thm:forrt-st-given-order-np-complete}
The \(\mathcal F_{rt\text{-}st}\)-Avoidance problem with a given order is NP-complete.
More precisely, for every fixed \(k\ge 3\), deciding whether there exists a partition into at most \(k\) parts
that avoids both \(\forrt\) and \(\forst\) with respect to a given order is NP-complete.
\end{theorem}

\begin{proof}[Proof sketch]
We reduce from \(k\)-\textsc{Coloring}. Let \(X\) be an instance with
\[
V(X)=\{x_1,\dots,x_n\},
\]
ordered arbitrarily as
\[
x_1\prec x_2\prec\cdots\prec x_n.
\]
For each edge \(x_i x_j\in E(X)\), with \(i<j\), add one witness vertex \(w_{ij}\), adjacent only to \(x_j\).
We keep the original edge \(x_i x_j\). Thus,
\[
V(G)=V(X)\cup \{w_{ij}:x_i x_j\in E(X),\ i<j\},
\]
and
\[
E(G)=E(X)\cup \{w_{ij}x_j:x_i x_j\in E(X),\ i<j\}.
\]
The given order puts all witnesses first, in arbitrary order, and then the original vertices:
\[
\text{all } w_{ij}\prec x_1\prec x_2\prec\cdots\prec x_n.
\]
In particular, for every edge \(x_i x_j\in E(X)\), with \(i<j\), we have
\[
w_{ij}\prec x_i\prec x_j,
\qquad
w_{ij}\sim x_j,
\qquad
x_i\sim x_j.
\]

If \(X\) has a proper \(k\)-coloring, keep the colors on the original vertices and assign each witness
\(w_{ij}\) the color of \(x_i\). Then every earlier neighbor of an original vertex has color different from it,
so no ordered triple realizes either \(\forrt\) or \(\forst\).

Conversely, suppose that a partition avoiding both \(\forrt\) and \(\forst\) exists. For each edge
\(x_i x_j\in E(X)\), with \(i<j\), the triple
\[
w_{ij}\prec x_i\prec x_j
\]
has the required edges to \(x_j\). Since \(\forst\) is avoided, \(x_i\) and \(x_j\) must lie in different parts.
Therefore, the restriction of the partition to \(V(X)\) is a proper \(k\)-coloring of \(X\).

Thus \(X\) is \(k\)-colorable if and only if the constructed ordered graph admits a partition into at most \(k\)
parts avoiding \(\mathcal F_{rt\text{-}st}\). The problem is in NP, so it is NP-complete.
\end{proof}







\begin{theorem}\label{thm:forrt-st-general-np-complete}
The \(\mathcal F_{rt\text{-}st}\)-Avoidance problem is NP-complete.
More precisely, for every fixed \(k\ge 3\), deciding whether a graph admits a solution with at most \(k\) parts
that avoids both \(\forrt\) and \(\forst\) is NP-complete.
\end{theorem}

\begin{proof}[Proof sketch]
We reduce from \(k\)-\textsc{Coloring}. Given an instance \(X\), we construct \(G\) by adding one new vertex \(F\)
adjacent to every vertex of \(X\), while keeping all edges of \(X\). That is,
\[
V(G)=V(X)\cup\{F\},
\qquad
E(G)=E(X)\cup \{Fx:x\in V(X)\}.
\]

If \(X\) has a proper \(k\)-coloring, choose one color \(\alpha\), assign color \(\alpha\) to \(F\), and order the vertices as
\[
F \prec \text{vertices of color }\alpha \prec \text{all remaining vertices}.
\]
Then no vertex has two earlier neighbors one of which lies in its own part, so the solution avoids both
\(\forrt\) and \(\forst\).

Conversely, suppose that \(G\) has a solution avoiding both \(\forrt\) and \(\forst\), with at most \(k\) parts.
Let \(\alpha\) be the part containing \(F\), and write
\[
A:=\{x\in V(X):x\prec F\},
\qquad
B:=\{x\in V(X):F\prec x\}.
\]
The key observation is the following: if a vertex \(t\) has at least two earlier neighbors, then all its earlier
neighbors must lie in parts different from the part of \(t\).

From this observation, the coloring induced by the solution is already proper on \(B\), and also on the edges between
\(A\) and \(B\). The only possible monochromatic edges of \(X\) lie inside \(A\). Moreover, for each part \(\beta\),
the subgraph induced by the vertices of \(A\) in part \(\beta\) is a forest: if it contained a cycle, the last vertex
of the cycle in the order would have two earlier neighbors in its own part, contradicting the observation.

We now recolor the monochromatic forest components inside \(A\). For each part \(\beta\neq \alpha\), root every tree
of the forest induced by \(A\cap \beta\), keep one side of the bipartition in color \(\beta\), and recolor the other
side with the auxiliary color \(\alpha\). The vertices recolored with \(\alpha\) form an independent set, again by the
observation above. Thus this recoloring destroys all monochromatic edges inside \(A\), and creates no new monochromatic
edges.

Hence we obtain a proper \(k\)-coloring of \(X\). Therefore \(X\) is \(k\)-colorable if and only if \(G\) admits a
solution with at most \(k\) parts avoiding \(\mathcal F_{rt\text{-}st}\). The construction is polynomial, and the
problem is in NP, so it is NP-complete.
\end{proof}







\end{document}